% Special Continous Distributions — Further Mathematics Upper Sixth — Cours signé OnBuch+
% Official MINESEC syllabus. Compiler avec : tectonic FMA20-special-continous-distributions.tex
\newcommand{\DOCMATIERE}{Further Mathematics}
\newcommand{\DOCNIVEAU}{Upper Sixth}
\newcommand{\DOCMODULE}{Upper Sixth — Further mathematics}
\newcommand{\DOCLECON}{20}
\newcommand{\DOCTITRE}{Special Continous Distributions}
% OnBuch+ — préambule « cours signés » ENRICHI (compilé avec tectonic / XeLaTeX)
% Système visuel « Pop » d'OnBuch+ : fond crème (jamais blanc), bordures encre
% épaisses, ombres « dures » décalées, titres Archivo Black en capitales.
% Version MATHÉMATIQUES : graphes (pgfplots/tikz), boîtes pédagogiques
% (définition, propriété, méthode, attention, le savais-tu, exercice résolu,
% exercice, TP, bilan), page de garde et signature OnBuch+ ancrée en bas.
%
% Macros attendues dans main.tex AVANT \input{preamble} :
%   \DOCMATIERE  \DOCNIVEAU  \DOCMODULE  \DOCLECON (numéro)  \DOCTITRE
\documentclass[11pt]{article}

\usepackage{fontspec}
\usepackage[english]{babel}
\usepackage[a4paper,margin=2cm,top=3.4cm,bottom=2.8cm,headheight=1.4cm,footskip=1.3cm]{geometry}
\usepackage{amsmath,amssymb,amsfonts}
\usepackage{mathtools} % \xmapsto, \coloneqq... souvent utilisés par les modèles
\usepackage{cancel} % simplifications barrées (bilans de forces)
% \abs{z} (module, valeur absolue) et \norm{u} (norme), utilisés par les modèles
\providecommand{\abs}{}\let\abs\relax\DeclarePairedDelimiter{\abs}{\lvert}{\rvert}
\providecommand{\norm}{}\let\norm\relax\DeclarePairedDelimiter{\norm}{\lVert}{\rVert}
\usepackage[table]{xcolor} % [table] : fournit aussi \rowcolors (alternance de lignes)
\usepackage{pagecolor}
\usepackage{tikz}
\usetikzlibrary{arrows.meta,positioning,calc,shapes.geometric,shapes.misc,shapes.arrows,decorations.pathmorphing,decorations.pathreplacing,decorations.markings,patterns,shadows,mindmap,trees,fit,backgrounds,matrix,intersections,angles,quotes}
\usepackage{pgfplots}
\usepackage[version=4]{mhchem} % équations nucléaires \ce{^{235}_{92}U}
\usepackage[european,siunitx]{circuitikz} % schémas électriques
\pgfplotsset{compat=1.18}
\usepgfplotslibrary{fillbetween,groupplots} % aires entre courbes (calcul intégral)
\usepackage{enumitem}
\usepackage{fancyhdr}
% [most] charge toutes les bibliothèques tcolorbox (listings, minted, raster,
% documentation...) et gonfle inutilement la mémoire TeX sur un document long
% (~300 boîtes) : on ne charge que ce qui est réellement utilisé.
\usepackage[skins,breakable]{tcolorbox}
\usepackage{graphicx}
\usepackage{colortbl}
\usepackage{booktabs}
\usepackage{tabularx}
\usepackage{diagbox} % case d'en-tête diagonale des tableaux à double entrée
% Colonnes extensibles centrée / gauche / droite (C, L, R), souvent utilisées par les modèles
\newcolumntype{C}{>{\centering\arraybackslash}X}
\newcolumntype{L}{>{\raggedright\arraybackslash}X}
\newcolumntype{R}{>{\raggedleft\arraybackslash}X}
\usepackage{array}
\usepackage{multicol}
\usepackage{siunitx}
% parse-numbers=false : accepte aussi \SI{2,5 \times 10^{-3}}{...}
\sisetup{locale=UK,output-decimal-marker={.},group-minimum-digits=4,parse-numbers=false}
\DeclareSIUnit\ohm{\ensuremath{\symup{\Omega}}} % U+2126 absent de la police
\DeclareSIUnit\division{div} % réglages d'oscilloscope (ms/div, V/div)
\DeclareSIUnit\tr{tr} % tours (vitesse de rotation, ex. \SI{1500}{\tr\per\minute})
\usepackage{qrcode}
% \degree : commande fréquemment utilisée par les modèles pour un angle ou une
% température en dehors de \SI{}{} (non fournie par les paquets chargés).
\providecommand{\degree}{^{\circ}}
% \qed (fin de démonstration), utilisé par les modèles sans amsthm
\providecommand{\qed}{\hfill$\square$}
% \barre : double barre des valeurs interdites dans un tableau de variations (tkz-tab)
\providecommand{\barre}{\ensuremath{\|}}
% environnement proof (amsthm non chargé) utilisé par les modèles
\ifdefined\proof\else\newenvironment{proof}[1][Proof]{\par\noindent\textit{#1.}\ }{\qed\par}\fi
\usepackage{lastpage}
\usepackage{hyperref}
\hypersetup{hidelinks}

% ============================================================
% Palette « Pop » OnBuch+ — mêmes hex que la classe P (plus_ui.dart)
% ============================================================
\definecolor{poporange}{HTML}{F59321}
\definecolor{popdark}{HTML}{E07A0C}
\definecolor{popcream}{HTML}{FFF6E8}
\definecolor{popink}{HTML}{1C1714}
\definecolor{poppurple}{HTML}{7A5AE0}
\definecolor{popgreen}{HTML}{2E9E6A}
\definecolor{poppink}{HTML}{E0457B}
\definecolor{popblue}{HTML}{2F7DE1}
\definecolor{popgold}{HTML}{C98A12}
\definecolor{popmuted}{HTML}{8A7F76}
\definecolor{popline}{HTML}{EADFD2}
\definecolor{popgray}{HTML}{8A7F76}
\colorlet{popgrey}{popgray}
% Alias tolérants (le modèle nomme parfois les couleurs différemment)
\colorlet{popwhite}{white}
\colorlet{popblack}{popink}
\colorlet{popred}{poppink}
\colorlet{popyellow}{popgold}
% Teintes claires pour fonds de tableaux / zones
\colorlet{popblueL}{popblue!10!white}
\colorlet{poporangeL}{poporange!14!white}
\colorlet{popgreenL}{popgreen!12!white}
\colorlet{poppurpleL}{poppurple!12!white}
\colorlet{poppinkL}{poppink!12!white}
\colorlet{popgoldL}{popgold!14!white}

\pagecolor{popcream}

% ============================================================
% Typographie
% ============================================================
\defaultfontfeatures{Ligatures=TeX}
\setmainfont{PlusJakartaSans-Medium.ttf}[Path=../../../fonts/,BoldFont=PlusJakartaSans-Bold.ttf]
% Maths en sans-serif (Fira Math) avec chiffres et lettres droites de Plus
% Jakarta, pour que les formules se fondent dans le texte.
\usepackage[math-style=ISO,bold-style=ISO]{unicode-math}
\setmathfont{FiraMath-Regular.otf}[Path=../../../fonts/]
\setmathfont{PlusJakartaSans-Medium.ttf}[Path=../../../fonts/,range={up/num,up/{Latin,latin},\mathup}]
% (icomma retiré : décimales avec point en anglais)
% Caractères mathématiques Unicode absents des polices de texte (Plus Jakarta,
% Archivo Black) : rendus par leur équivalent mathématique, en texte comme en
% formule — sinon ils s'affichent en carré vide (ex. « Arithmétique dans ℤ »).
\usepackage{newunicodechar}
\newunicodechar{ℝ}{\ensuremath{\mathbb{R}}}
\newunicodechar{ℤ}{\ensuremath{\mathbb{Z}}}
\newunicodechar{ℂ}{\ensuremath{\mathbb{C}}}
\newunicodechar{ℕ}{\ensuremath{\mathbb{N}}}
\newunicodechar{ℚ}{\ensuremath{\mathbb{Q}}}
\newunicodechar{𝔻}{\ensuremath{\mathbb{D}}}
\newunicodechar{α}{\ensuremath{\alpha}}
\newunicodechar{β}{\ensuremath{\beta}}
\newunicodechar{γ}{\ensuremath{\gamma}}
\newunicodechar{δ}{\ensuremath{\delta}}
\newunicodechar{ε}{\ensuremath{\varepsilon}}
\newunicodechar{θ}{\ensuremath{\theta}}
\newunicodechar{λ}{\ensuremath{\lambda}}
\newunicodechar{μ}{\ensuremath{\mu}}
\newunicodechar{σ}{\ensuremath{\sigma}}
\newunicodechar{τ}{\ensuremath{\tau}}
\newunicodechar{φ}{\ensuremath{\varphi}}
\newunicodechar{ψ}{\ensuremath{\psi}}
\newunicodechar{ω}{\ensuremath{\omega}}
\newunicodechar{Σ}{\ensuremath{\Sigma}}
% majuscules produites par \MakeUppercase dans les titres (α-aminés -> Α)
\newunicodechar{Α}{\ensuremath{\alpha}}
\newunicodechar{Β}{\ensuremath{\beta}}
\newunicodechar{Γ}{\ensuremath{\Gamma}}
\newunicodechar{Δ}{\ensuremath{\Delta}}
\newunicodechar{Ε}{\ensuremath{\varepsilon}}
\newunicodechar{Θ}{\ensuremath{\Theta}}
\newunicodechar{Λ}{\ensuremath{\Lambda}}
\newunicodechar{Μ}{\ensuremath{\mu}}
\newunicodechar{Τ}{\ensuremath{\tau}}
\newunicodechar{Φ}{\ensuremath{\Phi}}
\newunicodechar{Ψ}{\ensuremath{\Psi}}
\newunicodechar{Ω}{\ensuremath{\Omega}}
\newunicodechar{↦}{\ensuremath{\mapsto}}
\newunicodechar{∈}{\ensuremath{\in}}
\newunicodechar{∉}{\ensuremath{\notin}}
\newunicodechar{∧}{\ensuremath{\wedge}}
\newunicodechar{⇔}{\ensuremath{\Leftrightarrow}}
\newunicodechar{⟺}{\ensuremath{\Longleftrightarrow}}
\newunicodechar{⇒}{\ensuremath{\Rightarrow}}
\newunicodechar{⟹}{\ensuremath{\Longrightarrow}}
\newunicodechar{∘}{\ensuremath{\circ}}
\newunicodechar{∪}{\ensuremath{\cup}}
\newunicodechar{∩}{\ensuremath{\cap}}
\newunicodechar{≡}{\ensuremath{\equiv}}
\newunicodechar{⊂}{\ensuremath{\subset}}
\newunicodechar{⊥}{\ensuremath{\perp}}
\newunicodechar{∃}{\ensuremath{\exists}}
\newunicodechar{∀}{\ensuremath{\forall}}
\newunicodechar{∛}{\ensuremath{\sqrt[3]{\,}}}
\newunicodechar{∜}{\ensuremath{\sqrt[4]{\,}}}
\newunicodechar{⃗}{\ensuremath{{}^{\to}}}
\newunicodechar{ⁿ}{\ifmmode^{n}\else\textsuperscript{\ensuremath{n}}\fi}
\newunicodechar{ˣ}{\ifmmode^{x}\else\textsuperscript{\ensuremath{x}}\fi}
\newunicodechar{ᵇ}{\ifmmode^{b}\else\textsuperscript{\ensuremath{b}}\fi}
\newunicodechar{ʳ}{\ifmmode^{r}\else\textsuperscript{\ensuremath{r}}\fi}
\newunicodechar{ᵘ}{\ifmmode^{u}\else\textsuperscript{\ensuremath{u}}\fi}
\newunicodechar{ʸ}{\ifmmode^{y}\else\textsuperscript{\ensuremath{y}}\fi}
\newunicodechar{ᵃ}{\ifmmode^{a}\else\textsuperscript{\ensuremath{a}}\fi}
\newunicodechar{ᵉ}{\ifmmode^{e}\else\textsuperscript{\ensuremath{e}}\fi}
\newunicodechar{ᵏ}{\ifmmode^{k}\else\textsuperscript{\ensuremath{k}}\fi}
\newunicodechar{ᵖ}{\ifmmode^{p}\else\textsuperscript{\ensuremath{p}}\fi}
\newunicodechar{ᴺ}{\ifmmode^{N}\else\textsuperscript{\ensuremath{N}}\fi}
\newunicodechar{ᶜ}{\ifmmode^{c}\else\textsuperscript{\ensuremath{c}}\fi}
\newunicodechar{ʰ}{\ifmmode^{h}\else\textsuperscript{\ensuremath{h}}\fi}
\newunicodechar{ᵠ}{\ifmmode^{q}\else\textsuperscript{\ensuremath{q}}\fi}
\newunicodechar{ₙ}{\ifmmode_{n}\else\textsubscript{\ensuremath{n}}\fi}
\newunicodechar{ₐ}{\ifmmode_{a}\else\textsubscript{\ensuremath{a}}\fi}
\newunicodechar{ᵢ}{\ifmmode_{i}\else\textsubscript{\ensuremath{i}}\fi}
\newunicodechar{ᵣ}{\ifmmode_{r}\else\textsubscript{\ensuremath{r}}\fi}
\newunicodechar{ₖ}{\ifmmode_{k}\else\textsubscript{\ensuremath{k}}\fi}
\newunicodechar{ᵦ}{\ifmmode_{\beta}\else\textsubscript{\ensuremath{\beta}}\fi}
\newunicodechar{ₓ}{\ifmmode_{x}\else\textsubscript{\ensuremath{x}}\fi}
\newfontfamily\popdisplay{ArchivoBlack-Regular.ttf}[Path=../../../fonts/]
\newfontfamily\poplogo{SpaceGrotesk-Bold.ttf}[Path=../../../fonts/]
\newfontfamily\popsemi{PlusJakartaSans-SemiBold.ttf}[Path=../../../fonts/]
\newfontfamily\popxbold{PlusJakartaSans-ExtraBold.ttf}[Path=../../../fonts/]
\newfontfamily\popxboldls{PlusJakartaSans-ExtraBold.ttf}[Path=../../../fonts/,LetterSpace=3.0]

\setlist[enumerate]{leftmargin=1.7em,itemsep=5pt,topsep=4pt}
\setlist[itemize]{leftmargin=1.7em,itemsep=4pt,topsep=4pt,label={\color{poporange}\textbullet}}
\setlength{\parindent}{0pt}
\setlength{\parskip}{5pt}
\renewcommand{\arraystretch}{1.25}

% pgfplots : style Pop par défaut
\pgfplotsset{
  every axis/.append style={axis line style={popink,line width=1pt},tick style={popink},
    grid style={popline,line width=0.5pt},label style={font=\small},tick label style={font=\footnotesize}},
  popcurve/.style={poporange,line width=1.8pt,no markers,smooth},
}
% Même style disponible dans un \draw TikZ simple (hors pgfplots)
\tikzset{popcurve/.style={poporange,line width=1.8pt}}
\tikzset{popgrid/.style={popline,very thin}}
\colorlet{popcurve}{poporange} % le nom du style est parfois utilisé comme couleur

% ============================================================
% Pill / squiggle
% ============================================================
\newcommand{\poppill}[3]{%
  \tikz[baseline=-0.55ex]\node[draw=popink,line width=1.2pt,rounded corners=2.6pt,%
    fill=#2,inner xsep=6.5pt,inner ysep=3pt]{%
    \popxboldls\fontsize{7}{7}\selectfont\color{#3}\MakeUppercase{#1}};%
}
\newcommand{\popsquiggle}[1]{%
  \tikz[baseline=-0.4ex]{
    \draw[#1,line width=2.6pt,line cap=round]
      plot [smooth] coordinates {(0,0.09) (0.34,0.01) (0.68,0.21) (1.02,0.01) (1.36,0.10)};
  }%
}
% Mot-clé surligné (marqueur orange sous le texte)
\newcommand{\cle}[1]{{\popxbold\color{popdark}#1}}

% ============================================================
% En-tête / pied
% ============================================================
\pagestyle{fancy}
\fancyhf{}
\renewcommand{\headrulewidth}{0pt}
\renewcommand{\footrulewidth}{0pt}
\lhead{\raisebox{-0.15ex}{\poplogo\fontsize{13}{13}\selectfont\color{popink}OnBuch\color{poporange}+}}
\rhead{\poppill{\DOCMATIERE\ \textbullet\ \DOCNIVEAU\ \textbullet\ Chapter \DOCLECON}{white}{popink}}
\lfoot{\popxbold\fontsize{7.6}{7.6}\selectfont\color{popmuted}\MakeUppercase{Signed course OnBuch+}\ \textbullet\ {\popsemi\DOCTITRE}}
\rfoot{\tikz[baseline=-0.6ex]\node[rounded corners=8.5pt,fill=popink,minimum height=17pt,minimum width=17pt,inner xsep=5pt,inner sep=0]{\popxbold\fontsize{8}{8}\selectfont\color{popcream}\thepage\,/\,\pageref*{LastPage}};}

% ============================================================
% Titres
% ============================================================
\newcommand{\chaptitre}[2]{%
  \begin{tcolorbox}[enhanced,colback=poporange,colframe=popink,boxrule=2.6pt,arc=11pt,
      drop shadow=popink,left=15pt,right=15pt,top=12pt,bottom=11pt,before skip=3pt,after skip=18pt]
    \poppill{#2}{popink}{popcream}\par\vspace{9pt}
    {\raggedright\hyphenpenalty=10000\exhyphenpenalty=10000\popdisplay\fontsize{22}{24}\selectfont\color{popcream}\MakeUppercase{#1}\par}
  \end{tcolorbox}
}
% Section numérotée : pastille orange avec le numéro + titre Archivo
\newcounter{popsec}
\newcommand{\coursec}[1]{%
  \stepcounter{popsec}\setcounter{popsub}{0}%
  \par\vspace{16pt}\noindent
  \begin{minipage}{\linewidth}
  \tikz[baseline=-0.7ex]\node[draw=popink,line width=1.6pt,fill=poporange,rounded corners=4pt,minimum size=22pt,inner sep=0]{\popdisplay\fontsize{12}{12}\selectfont\color{popcream}\thepopsec};%
  \hspace{8pt}{\popdisplay\fontsize{15}{17}\selectfont\color{popink}\MakeUppercase{#1}}\par
  \vspace{4pt}\noindent{\color{poporange}\rule{36pt}{3.6pt}}
  \end{minipage}\par\nopagebreak\vspace{8pt}
}
\newcounter{popsub}
\newcommand{\courssub}[1]{%
  \stepcounter{popsub}%
  \par\vspace{9pt}\noindent{\popxbold\fontsize{11.5}{13}\selectfont\color{popink}{\color{poporange}\thepopsec.\thepopsub}\hspace{6pt}#1}\par\nopagebreak\vspace{4pt}
}

% ============================================================
% Boîtes pédagogiques — même recette : fond blanc, bordure épaisse colorée,
% ombre dure encre, badge plein. Titre optionnel [#1].
% ============================================================
\tcbset{popbase/.style={breakable,enhanced,colback=white,boxrule=2.3pt,arc=9pt,
  shadow={3pt}{-3pt}{0pt}{popink},left=14pt,right=14pt,top=11pt,bottom=11pt,before skip=11pt,after skip=15pt,
  fontupper=\popsemi\fontsize{10.6}{14.5}\selectfont\color{popink},
  fontlower=\popsemi\fontsize{10.6}{14.5}\selectfont\color{popink}}}
% Coche dessinée (le glyphe ✓ n'existe pas dans les polices du document)
\newcommand{\popcheck}[1]{\tikz[baseline=-0.5ex]\draw[#1,line width=1.6pt,line cap=round,line join=round](-0.13,0.0)--(-0.04,-0.09)--(0.13,0.1);}
\providecommand{\coche}{\popcheck{popgreen}}
\providecommand{\resultat}[1]{\par\smallskip\noindent\fcolorbox{popgreen}{popgreenL}{\parbox{\dimexpr\linewidth-2\fboxsep-2\fboxrule\relax}{\textbf{Result.} #1}}\par\smallskip}
\newcommand{\popboxhead}[3]{\poppill{#1}{#2}{white}\ifx\relax#3\relax\else\hspace{6pt}{\popxbold\fontsize{10.6}{12}\selectfont\color{popink}#3}\fi\par\vspace{7pt}}

\newtcolorbox{aretenir}[1][]{popbase,colframe=popgreen,before upper={\popboxhead{Key points}{popgreen}{#1}}}
\newtcolorbox{exemplebox}[1][]{popbase,colframe=poporange,before upper={\popboxhead{Exemple}{poporange}{#1}}}
\newtcolorbox{definition}[1][]{popbase,colframe=popblue,before upper={\popboxhead{Definition}{popblue}{#1}}}
\newtcolorbox{propriete}[1][]{popbase,colframe=poppurple,before upper={\popboxhead{Property}{poppurple}{#1}}}
\newtcolorbox{methode}[1][]{popbase,colframe=popgold,before upper={\popboxhead{Method}{popgold}{#1}}}
\newtcolorbox{attention}[1][]{popbase,colframe=poppink,before upper={\popboxhead{Warning}{poppink}{#1}}}
\newtcolorbox{experience}[1][]{popbase,colframe=popblue,colback=popblueL,before upper={\popboxhead{Experiment}{popblue}{#1}}}
% « Le savais-tu ? » : carte violette pleine (contexte, culture, Cameroun)
\newtcolorbox{savaistu}[1][]{popbase,colback=poppurple,colframe=popink,
  fontupper=\popsemi\fontsize{10.4}{14.2}\selectfont\color{white},
  before upper={\poppill{Did you know?}{white}{poppurple}\ifx\relax#1\relax\else\hspace{6pt}{\popxbold\color{white}#1}\fi\par\vspace{7pt}}}
% Exercice résolu : énoncé en haut, \tcblower puis la solution sur fond crème
\newcounter{popexo}\newcounter{popexr}
\newtcolorbox{exoresolu}[1][]{popbase,colframe=poporange,colbacklower=popcream,
  segmentation style={popink,line width=1.2pt,dashed},
  before upper={\stepcounter{popexr}\popboxhead{Worked example \thepopexr}{poporange}{#1}},
  before lower={\poppill{Solution}{popink}{popcream}\par\vspace{6pt}}}
% Exercice (non résolu en place) — la correction est dans la partie Corrigés
\newtcolorbox{exercice}[1][]{popbase,colframe=popink,
  before upper={\stepcounter{popexo}\popboxhead{Exercise \thepopexo}{popink}{#1}}}
% Corrigé
\newtcolorbox{corrige}[1][]{popbase,colframe=popgreen,colback=popgreenL,
  before upper={\popboxhead{Solution}{popgreen}{#1}}}
% Objectifs / prérequis / bilan
\newtcolorbox{objectifs}{popbase,colframe=popink,colback=poporangeL,
  before upper={\popboxhead{Chapter objectives}{popink}{}}}
\newtcolorbox{prerequis}{popbase,colframe=popink,colback=popblueL,
  before upper={\popboxhead{Prerequisites}{popblue}{}}}
\newtcolorbox{bilan}{popbase,colframe=popink,colback=white,boxrule=2.8pt,
  before upper={\popboxhead{Fiche bilan}{poporange}{L'essentiel à connaître pour l'examen}}}
% Figure encadrée avec légende
\newtcolorbox{popfigure}[1][]{enhanced,breakable=false,colback=white,colframe=popline,boxrule=1.4pt,arc=8pt,
  left=10pt,right=10pt,top=10pt,bottom=8pt,before skip=10pt,after skip=12pt,halign=center,
  lower separated=false,halign lower=center,
  fontlower=\popsemi\fontsize{9}{11.5}\selectfont\color{popmuted},
  #1}
\newcommand{\legende}[1]{\par\vspace{6pt}{\popsemi\fontsize{9}{11.5}\selectfont\color{popmuted}#1}}

% ============================================================
% Signature OnBuch+
% ============================================================
\newcommand{\poplogoinline}[1]{{\poplogo\fontsize{#1}{#1}\selectfont\color{popink}OnBuch\color{poporange}+}}
\newcommand{\signatureonbuch}{%
  \begin{tcolorbox}[enhanced,colback=popink,colframe=popink,boxrule=2.2pt,arc=10pt,
      drop shadow=poporange,left=14pt,right=14pt,top=10pt,bottom=10pt,before skip=0pt,after skip=0pt]
    \tikz[baseline=-0.7ex]\node[circle,fill=popgreen,draw=popcream,line width=1.3pt,minimum size=22pt,inner sep=0]{\popcheck{white}};%
    \hspace{9pt}\begin{minipage}[c]{0.62\linewidth}
      {\popxbold\fontsize{12}{13}\selectfont\color{popcream}Signed\ \textbullet\ {\poplogo OnBuch\color{poporange}+}}\par\vspace{2pt}
      {\raggedright\popsemi\fontsize{8}{10.5}\selectfont\color{popline}\MakeUppercase{OnBuch+ teaching team}\\\MakeUppercase{Follows the official MINESEC syllabus}\par}
    \end{minipage}\hfill
    {\poplogo\fontsize{22}{22}\selectfont\color{popcream}OnBuch\color{poporange}+}
  \end{tcolorbox}%
}

% ============================================================
% Page de garde
% #1 titre · #2 matière · #3 niveau · #4 module · #5 n° leçon · #6 durée ·
% #7 description / objectifs (texte ou itemize)
% ============================================================
\newcommand{\pagegarde}[7]{%
  \thispagestyle{empty}
  \newgeometry{a4paper,margin=1.7cm,top=1.6cm,bottom=1.4cm}%
  \begin{tikzpicture}[remember picture,overlay]
    \fill[poporange!18] ([xshift=-3.2cm,yshift=-3.6cm]current page.north east) circle (4.6cm);
    \fill[poppurple!14] ([xshift=2.4cm,yshift=7.6cm]current page.south west) circle (3.8cm);
  \end{tikzpicture}%
  {\poplogo\fontsize{30}{30}\selectfont\color{popink}OnBuch\color{poporange}+}\hfill
  \raisebox{0.6ex}{\poppill{Signed course \textbullet\ Premium}{popink}{popcream}}\par
  \vspace{1pt}\popsquiggle{poppurple}\par
  \vspace{8pt}
  \begin{tcolorbox}[enhanced,colback=poporange,colframe=popink,boxrule=2.8pt,arc=13pt,
      drop shadow=popink,left=18pt,right=18pt,top=12pt,bottom=13pt,after skip=10pt]
    \poppill{#2 \textbullet\ #3}{popink}{popcream}\hspace{5pt}\poppill{#4}{white}{popink}\par\vspace{8pt}
    {\popdisplay\fontsize{14}{14}\selectfont\color{popink}CHAPTER #5}\par\vspace{4pt}
    {\raggedright\hyphenpenalty=10000\exhyphenpenalty=10000\popdisplay\fontsize{25}{27}\selectfont\color{popcream}\MakeUppercase{#1}\par}
    \vspace{8pt}
    {\popxbold\fontsize{9.5}{9.5}\selectfont\color{popink}\MakeUppercase{Suggested time: #6}}
  \end{tcolorbox}
  \begin{tcolorbox}[enhanced,colback=white,colframe=popink,boxrule=2.2pt,arc=10pt,
      drop shadow=popink,left=16pt,right=16pt,top=10pt,bottom=10pt,after skip=8pt]
    \poppill{In this chapter}{poppurple}{white}\par\vspace{5pt}
    \popsemi\fontsize{9.3}{12.6}\selectfont\color{popink}#7
  \end{tcolorbox}
  \noindent\begin{minipage}[t]{0.487\linewidth}
    \begin{tcolorbox}[enhanced,colback=popgreen,colframe=popink,boxrule=2.2pt,arc=10pt,
        drop shadow=popink,left=11pt,right=11pt,top=10pt,bottom=10pt,height=3.1cm,valign=center]
      \tcbox[colback=white,colframe=popink,boxrule=1.3pt,arc=5pt,boxsep=0pt,nobeforeafter,
        left=3pt,right=3pt,top=3pt,bottom=3pt,valign=center]{\qrcode[height=1.9cm]{https://play.google.com/store/apps/details?id=com.luvvix.onbuch&hl=en}}%
      \hspace{8pt}\begin{minipage}[c]{0.5\linewidth}
        {\popdisplay\fontsize{11}{12.5}\selectfont\color{white}\MakeUppercase{Download OnBuch}}\par\vspace{4pt}
        {\raggedright\popsemi\fontsize{8.4}{11}\selectfont\color{white}Free \textbullet\ Google Play\\AI tutor, past papers, results\par}
      \end{minipage}
    \end{tcolorbox}
  \end{minipage}\hfill
  \begin{minipage}[t]{0.487\linewidth}
    \begin{tcolorbox}[enhanced,colback=poppurple,colframe=popink,boxrule=2.2pt,arc=10pt,
        drop shadow=popink,left=11pt,right=11pt,top=10pt,bottom=10pt,height=3.1cm,valign=center]
      \tikz[baseline=-0.7ex]\node[circle,fill=white,draw=popink,line width=1.3pt,minimum size=1.6cm,inner sep=0]{\popdisplay\fontsize{24}{24}\selectfont\color{poppurple}+};%
      \hspace{8pt}\begin{minipage}[c]{0.6\linewidth}
        {\popdisplay\fontsize{11}{12.5}\selectfont\color{white}\MakeUppercase{Go OnBuch+}}\par\vspace{4pt}
        {\raggedright\popsemi\fontsize{8.4}{11}\selectfont\color{white}All signed courses, unlimited AI tutor, PDF sheets — OnBuch+ tab in the app\par}
      \end{minipage}
    \end{tcolorbox}
  \end{minipage}
  \vspace{8pt}
  \popcontenu
  \vfill
  \signatureonbuch
  \restoregeometry
  \newpage
}

% Rangée « Ce cours contient » (page de garde)
\newcommand{\poptile}[3]{% #1 couleur #2 gros texte #3 légende
  \begin{tcolorbox}[enhanced,colback=#1,colframe=popink,boxrule=2pt,arc=9pt,shadow={2.5pt}{-2.5pt}{0pt}{popink},
      left=6pt,right=6pt,top=9pt,bottom=9pt,halign=center,valign=center,height=1.75cm,width=\linewidth,nobeforeafter]
    {\popdisplay\fontsize{12.5}{13}\selectfont\color{white}\MakeUppercase{#2}}\par\vspace{3pt}
    {\centering\popsemi\fontsize{7.8}{9.5}\selectfont\color{white}#3\par}
  \end{tcolorbox}}
\newcommand{\popcontenu}{%
  \par\noindent{\popxboldls\fontsize{7.5}{7.5}\selectfont\color{popmuted}THIS SIGNED COURSE CONTAINS}\par\vspace{7pt}
  \noindent\begin{minipage}[t]{0.235\linewidth}\poptile{popblue}{Course}{complete, illustrated, step by step}\end{minipage}\hfill
  \begin{minipage}[t]{0.235\linewidth}\poptile{poporange}{Methods}{and worked examples}\end{minipage}\hfill
  \begin{minipage}[t]{0.235\linewidth}\poptile{poppink}{Exercises}{exam style + solutions}\end{minipage}\hfill
  \begin{minipage}[t]{0.235\linewidth}\poptile{popgreen}{Summary sheet}{the essentials to remember}\end{minipage}\par}

% Sommaire
\newtcolorbox{sommaire}{enhanced,colback=white,colframe=popink,boxrule=2.2pt,arc=10pt,
  drop shadow=popink,left=18pt,right=18pt,top=13pt,bottom=13pt,before skip=6pt,after skip=20pt,
  before upper={\poppill{Contents}{poppurple}{white}\par\vspace{9pt}\popsemi\fontsize{10.6}{14.5}\selectfont\color{popink}}}

% Signature de fin de cours — poussée tout en bas de la dernière page
\newcommand{\findecours}{%
  \par\vfill\nopagebreak
  \begin{center}\popsquiggle{poporange}\end{center}\vspace{4pt}
  \signatureonbuch
}
\AtBeginDocument{\def\llbracket{\mathopen{[\mkern-3.5mu[}}\def\rrbracket{\mathclose{]\mkern-3.5mu]}}} % intervalles d'entiers (glyphe ⟦ absent de la police)
% Glyphes absents de Fira Math : redéfinis à partir de glyphes disponibles
\AtBeginDocument{%
  \def\setminus{\mathbin{\backslash}}\let\smallsetminus\setminus
  \def\vdots{\mathord{\vbox{\baselineskip3.2pt\lineskiplimit0pt\kern4pt\hbox{.}\hbox{.}\hbox{.}}}}%
  \def\ddots{\mathinner{\mkern1mu\raise6.5pt\hbox{.}\mkern2mu\raise3.6pt\hbox{.}\mkern2mu\raise0.7pt\hbox{.}\mkern1mu}}%
  \def\triangle{\mathord{\scalebox{0.9}{$\symup{\Delta}$}}}%
  \def\nabla{\mathord{\raisebox{\depth}{\scalebox{1}[-1]{$\symup{\Delta}$}}}}%
  \def\checkmark{\popcheck{popdark}}%
  \def\star{\mathbin{\ast}}\def\bigstar{\mathord{\ast}}%
  \def\diamond{\mathbin{\scalebox{0.6}{$\Diamond$}}}%
  \def\bigcup{\mathop{\mathchoice{\raisebox{-0.3em}{\scalebox{1.7}{$\cup$}}}{\scalebox{1.25}{$\cup$}}{\cup}{\cup}}\displaylimits}%
  \def\bigcap{\mathop{\mathchoice{\raisebox{-0.3em}{\scalebox{1.7}{$\cap$}}}{\scalebox{1.25}{$\cap$}}{\cap}{\cap}}\displaylimits}%
  \def\lesseqqgtr{\mathrel{\lessgtr}}\def\bigoplus{\mathop{\oplus}\displaylimits}%
}
\providecommand{\ssi}{if and only if} % abréviation fréquente
\AtBeginDocument{\providecommand{\wideparen}{\overparen}} % arc de cercle

% Fonctions usuelles que les modèles utilisent sans que amsmath les fournisse
\providecommand{\arsinh}{\operatorname{arsinh}}\providecommand{\arcosh}{\operatorname{arcosh}}
\providecommand{\artanh}{\operatorname{artanh}}\providecommand{\arsech}{\operatorname{arsech}}
\providecommand{\arcsch}{\operatorname{arcsch}}\providecommand{\arcoth}{\operatorname{arcoth}}
\providecommand{\arcsinh}{\operatorname{arsinh}}\providecommand{\arccosh}{\operatorname{arcosh}}
\providecommand{\arctanh}{\operatorname{artanh}}\providecommand{\sech}{\operatorname{sech}}
\providecommand{\csch}{\operatorname{csch}}\providecommand{\arcsec}{\operatorname{arcsec}}
\providecommand{\arccsc}{\operatorname{arccsc}}\providecommand{\arccot}{\operatorname{arccot}}
\providecommand{\sgn}{\operatorname{sgn}}\providecommand{\lcm}{\operatorname{lcm}}
\providecommand{\Var}{\operatorname{Var}}\providecommand{\Cov}{\operatorname{Cov}}

\begin{document}

\pagegarde{Special Continous Distributions}{Further Mathematics}{Upper Sixth}{Upper Sixth — Further mathematics}{20}{7 periods}{This chapter studies the three special continuous distributions of the A-Level syllabus — exponential, uniform (rectangular) and normal — together with their means, variances and probability computations. It ends with the use of the normal distribution, with continuity correction, as an approximation to the binomial and Poisson distributions, a favourite GCE examination topic.
\vspace{4pt}
\begin{multicols}{2}\small
\begin{itemize}
\item By the end of this chapter I can recognise a situation modelled by an exponential distribution and state its probability density function.
\item By the end of this chapter I can compute probabilities, the mean E(X)=1/λ and the variance Var(X)=1/λ² for an exponential distribution.
\item By the end of this chapter I can define the uniform distribution on [a,b], and use its mean (a+b)/2 and variance (b−a)²/12.
\item By the end of this chapter I can state the properties of the normal distribution N(μ,σ²) and interpret its density curve.
\item By the end of this chapter I can standardise a normal variable using Z=(X−μ)/σ and read probabilities from the standard normal table Φ(z).
\item By the end of this chapter I can solve inverse problems: find a value x or a parameter μ, σ given a probability.
\end{itemize}\end{multicols}}

\begin{sommaire}
\begin{multicols}{2}
\begin{enumerate}
\item The Exponential Distribution
\item The Uniform (Rectangular) Distribution
\item The Normal Distribution: pdf, Mean and Variance
\item The Standard Normal Variable and Use of Tables
\item Normal Approximations to the Binomial and Poisson Distributions
\item Integration activity
\item Methods and worked examples
\item Exercises
\item Solutions to the exercises
\item Summary sheet
\end{enumerate}
\end{multicols}
\end{sommaire}

\begin{objectifs}
\begin{itemize}
\item By the end of this chapter I can recognise a situation modelled by an exponential distribution and state its probability density function.
\item By the end of this chapter I can compute probabilities, the mean E(X)=1/λ and the variance Var(X)=1/λ² for an exponential distribution.
\item By the end of this chapter I can define the uniform distribution on [a,b], and use its mean (a+b)/2 and variance (b−a)²/12.
\item By the end of this chapter I can state the properties of the normal distribution N(μ,σ²) and interpret its density curve.
\item By the end of this chapter I can standardise a normal variable using Z=(X−μ)/σ and read probabilities from the standard normal table Φ(z).
\item By the end of this chapter I can solve inverse problems: find a value x or a parameter μ, σ given a probability.
\item By the end of this chapter I can approximate a binomial distribution B(n,p) by a normal distribution when np and n(1−p) are large enough, applying the continuity correction.
\item By the end of this chapter I can approximate a Poisson distribution Po(λ) by N(λ,λ) for large λ, applying the continuity correction.
\end{itemize}
\end{objectifs}

\begin{prerequis}
\begin{itemize}
\item Continuous random variables: probability density functions, P(a≤X≤b) as an integral/area (Lower Sixth, Further Maths).
\item Expectation E(X) and variance Var(X) of continuous random variables.
\item The binomial distribution B(n,p) and the Poisson distribution Po(λ): conditions of use, means and variances.
\item Integration of exponential functions ∫e\^{}(−λx)dx and basic definite integrals.
\item Reading statistical tables (as done for the binomial/Poisson tables in Lower Sixth).
\end{itemize}
\end{prerequis}

\newpage

\coursec{The Exponential Distribution}

At the Mfoundi market in Yaoundé, a phone-repairer notices that the waiting time between two customers seems to follow a curious ``memoryless'' pattern, while the sacks of rice weighed by a wholesaler cluster around 50 kg with symmetric small errors. Meanwhile, an MTN engineer models the lifetime of a SIM card, and an examiner wants to know the chance that more than 60 of 100 candidates pass when each has probability 0.5 of passing. All these situations — waiting times, uniform errors, measurement fluctuations and large counts — are governed by \cle{special continuous distributions}. This chapter gives you the essential tools of the GCE Advanced Level: the \cle{exponential}, the \cle{uniform} and the \cle{normal} distributions, plus the powerful normal approximations to the binomial and Poisson laws.

We begin with a natural question: \emph{if events happen randomly in time at a steady average rate — customers arriving, phone calls received, machines breaking down — how long must we wait for the next event?} The answer is the exponential distribution.

\courssub{From Poisson events to waiting times}

Recall the Poisson distribution: if events occur independently at a constant average rate of $\lambda$ per unit time, then the \emph{number} of events $N$ in an interval of length $t$ satisfies
\[
P(N = k) = \frac{(\lambda t)^k e^{-\lambda t}}{k!}, \qquad k = 0, 1, 2, \dots
\]
Now look at time differently. Instead of counting events, measure the \emph{waiting time} $X$ until the next event. The key observation is:
\[
P(X > t) = P(\text{no event occurs in } [0, t]) = P(N = 0) = e^{-\lambda t}.
\]
So the random waiting time $X$ is a \emph{continuous} random variable on $[0, +\infty)$ whose tail probability decays exponentially. Differentiating $P(X \le t) = 1 - e^{-\lambda t}$ with respect to $t$ gives the density $\lambda e^{-\lambda t}$. This is exactly the model of the phone-repairer at Mfoundi: the gap between two customers is exponential.

\begin{popfigure}
\begin{center}
\begin{tikzpicture}[scale=1.0]
\draw[->, thick, popink] (0,0) -- (11,0) node[below left] {time (min)};
\foreach \x in {1.2, 3.0, 3.6, 6.4, 8.0, 10.2} {
  \fill[poporange] (\x,0) circle (2.2pt);
}
\node[above, poporange] at (1.2,0.1) {event};
\draw[<->, popblue, thick] (1.2,0.55) -- (3.0,0.55) node[midway, above, popblue] {$X_1$};
\draw[<->, popblue, thick] (3.0,0.55) -- (3.6,0.55) node[midway, above, popblue] {$X_2$};
\draw[<->, popblue, thick] (3.6,0.55) -- (6.4,0.55) node[midway, above, popblue] {$X_3$};
\draw[<->, popblue, thick] (6.4,0.55) -- (8.0,0.55) node[midway, above, popblue] {$X_4$};
\node[popmuted, align=center] at (5.5,-0.9) {gaps $X_1, X_2, X_3, \dots$ are independent $\mathrm{Exp}(\lambda)$};
\end{tikzpicture}
\legende{Events of a Poisson process (orange dots) on a time axis; the waiting times between consecutive events are independent exponential random variables.}
\end{center}
\end{popfigure}

\courssub{Definition and validity of the density}

\begin{definition}[Exponential distribution]
A continuous random variable $X$ follows the \cle{exponential distribution} with parameter $\lambda > 0$, written $X \sim \mathrm{Exp}(\lambda)$, if its probability density function is
\[
f(x) = \begin{cases} \lambda e^{-\lambda x} & \text{if } x \ge 0, \\ 0 & \text{if } x < 0. \end{cases}
\]
\end{definition}

Let us verify that $f$ is a genuine pdf. First, since $\lambda > 0$ and $e^{-\lambda x} > 0$, we have $f(x) \ge 0$ for all $x$. Second, the total area under the curve is
\[
\int_{-\infty}^{+\infty} f(x)\,dx = \int_{0}^{+\infty} \lambda e^{-\lambda x}\,dx = \left[ -e^{-\lambda x} \right]_{0}^{+\infty} = 0 - (-1) = 1.
\]
Both conditions hold, so $f$ is indeed a probability density function.

\begin{popfigure}
\begin{center}
\begin{tikzpicture}
\begin{axis}[
  width=11cm, height=6.5cm,
  xlabel={$x$}, ylabel={$f(x)$},
  xmin=0, xmax=6, ymin=0, ymax=2.2,
  domain=0:6, samples=100,
  legend style={at={(0.97,0.97)}, anchor=north east, font=\small},
  axis lines=left,
]
\addplot[popblue, very thick] {0.5*exp(-0.5*x)};
\addlegendentry{$\lambda = 0.5$}
\addplot[poporange, very thick] {exp(-x)};
\addlegendentry{$\lambda = 1$}
\addplot[popgreen, very thick] {2*exp(-2*x)};
\addlegendentry{$\lambda = 2$}
\addplot[fill=poporange, opacity=0.25, draw=none, domain=2:6] {exp(-x)} \closedcycle;
\node[popdark, align=center] at (axis cs:3.9,0.28) {$P(X>2)$\\ ($\lambda=1$)};
\end{axis}
\end{tikzpicture}
\legende{Densities of $\mathrm{Exp}(\lambda)$ for $\lambda = 0.5, 1, 2$. Larger $\lambda$ means faster decay and shorter typical waiting times. The shaded region is $P(X>2)=e^{-2}$ for $\lambda=1$.}
\end{center}
\end{popfigure}

\courssub{Cumulative distribution function and tail probability}

\begin{propriete}[cdf and tail probability]
If $X \sim \mathrm{Exp}(\lambda)$, then for $x \ge 0$,
\[
F(x) = P(X \le x) = 1 - e^{-\lambda x}, \qquad P(X > x) = e^{-\lambda x}.
\]
\end{propriete}

\textbf{Proof.} For $x \ge 0$,
\[
F(x) = \int_0^x \lambda e^{-\lambda t}\,dt = \left[ -e^{-\lambda t} \right]_0^x = 1 - e^{-\lambda x},
\]
and $P(X > x) = 1 - F(x) = e^{-\lambda x}$. \hfill $\square$

\begin{methode}[Computing $P(a \le X \le b)$]
For $0 \le a < b$ and $X \sim \mathrm{Exp}(\lambda)$:
\begin{enumerate}
\item Write $P(a \le X \le b) = F(b) - F(a)$.
\item Substitute: $P(a \le X \le b) = (1 - e^{-\lambda b}) - (1 - e^{-\lambda a}) = e^{-\lambda a} - e^{-\lambda b}$.
\item Evaluate numerically, keeping at least 4 decimal places.
\end{enumerate}
\end{methode}

\begin{attention}[Continuous variables: single points carry no probability]
Because $X$ is continuous, $P(X = a) = 0$ for every $a$. Hence
\[
P(a \le X \le b) = P(a < X < b) = P(a \le X < b) = P(a < X \le b).
\]
Do not hesitate between strict and non-strict inequalities: for the exponential distribution they all give the same value.
\end{attention}

\courssub{Mean and variance}

\begin{propriete}[Expectation and variance — proof examinable]
If $X \sim \mathrm{Exp}(\lambda)$, then
\[
E(X) = \frac{1}{\lambda}, \qquad \mathrm{Var}(X) = \frac{1}{\lambda^2}.
\]
\end{propriete}

\textbf{Derivation.} By definition, $E(X) = \displaystyle\int_0^{+\infty} x\,\lambda e^{-\lambda x}\,dx$. Integrate by parts with $u = x$, $dv = \lambda e^{-\lambda x}dx$, so $du = dx$, $v = -e^{-\lambda x}$:
\[
E(X) = \left[ -x e^{-\lambda x} \right]_0^{+\infty} + \int_0^{+\infty} e^{-\lambda x}\,dx = 0 + \left[ -\frac{1}{\lambda} e^{-\lambda x} \right]_0^{+\infty} = \frac{1}{\lambda}.
\]
(The bracket vanishes because $x e^{-\lambda x} \to 0$ as $x \to +\infty$.)

For the variance, compute $E(X^2) = \displaystyle\int_0^{+\infty} x^2 \lambda e^{-\lambda x}\,dx$ by parts with $u = x^2$, $dv = \lambda e^{-\lambda x}dx$:
\[
E(X^2) = \left[ -x^2 e^{-\lambda x} \right]_0^{+\infty} + 2\int_0^{+\infty} x e^{-\lambda x}\,dx = 0 + \frac{2}{\lambda}\,E(X) = \frac{2}{\lambda^2}.
\]
Therefore
\[
\mathrm{Var}(X) = E(X^2) - [E(X)]^2 = \frac{2}{\lambda^2} - \frac{1}{\lambda^2} = \frac{1}{\lambda^2}. \qquad \square
\]

\begin{attention}[$\lambda$ is a rate, not a mean]
The parameter $\lambda$ is a \cle{rate} (events per unit time), \textbf{not} the mean. The mean waiting time is $E(X) = \dfrac{1}{\lambda}$. If the mean lifetime of a bulb is 500 hours, then $\lambda = \dfrac{1}{500} = 0.002$ per hour — \emph{not} $\lambda = 500$. This confusion is the most frequent source of lost marks at the GCE. Note also that the standard deviation equals the mean: $\sigma = \dfrac{1}{\lambda} = E(X)$.
\end{attention}

\courssub{The memoryless property}

\begin{propriete}[Memoryless property]
If $X \sim \mathrm{Exp}(\lambda)$, then for all $s, t \ge 0$,
\[
P(X > s + t \mid X > s) = P(X > t).
\]
\end{propriete}

\textbf{Proof.} By the definition of conditional probability,
\[
P(X > s+t \mid X > s) = \frac{P(X > s+t \text{ and } X > s)}{P(X > s)} = \frac{P(X > s+t)}{P(X > s)} = \frac{e^{-\lambda(s+t)}}{e^{-\lambda s}} = e^{-\lambda t} = P(X > t). \qquad \square
\]

\textbf{Interpretation.} Suppose $X$ is the lifetime of a component. Given that it has already survived $s$ hours, the probability that it lasts \emph{at least $t$ more hours} is the same as the probability that a \emph{brand-new} component lasts $t$ hours. The component ``forgets'' its age: a used component is \emph{as good as new}. This is exactly the pattern the Mfoundi phone-repairer observed: however long you have already waited for a customer, the remaining wait has the same distribution as at the start.

\begin{exoresolu}[Lifetime of an electronic component]
The lifetime $X$ (in hours) of a certain electronic component follows an exponential distribution with mean 1000 hours.
\begin{enumerate}
\item[(a)] State the value of $\lambda$ and write down the pdf of $X$.
\item[(b)] Find the probability that the component lasts more than 1200 hours.
\item[(c)] Find the probability that it lasts between 500 and 1500 hours.
\item[(d)] Given that it has already lasted 800 hours, find the probability that it lasts at least 400 more hours.
\end{enumerate}
\tcblower
\textbf{Solution.}
\begin{enumerate}
\item[(a)] Since $E(X) = \dfrac{1}{\lambda} = 1000$, we get $\lambda = \dfrac{1}{1000} = 0.001$ per hour. The pdf is
\[ f(x) = 0.001\,e^{-0.001x}, \quad x \ge 0. \]
\item[(b)] $P(X > 1200) = e^{-0.001 \times 1200} = e^{-1.2} \approx 0.3012$.
\item[(c)] $P(500 \le X \le 1500) = e^{-0.001 \times 500} - e^{-0.001 \times 1500} = e^{-0.5} - e^{-1.5} \approx 0.6065 - 0.2231 = 0.3834$.
\item[(d)] By the memoryless property,
\[ P(X > 1200 \mid X > 800) = P(X > 400) = e^{-0.001 \times 400} = e^{-0.4} \approx 0.6703. \]
\end{enumerate}
\end{exoresolu}

\begin{exoresolu}[Finding the parameter from a probability]
The waiting time $T$ (in minutes) between two arrivals at a bus stop in Bamenda is exponentially distributed. Given that $P(T > 10) = 0.6$, find:
\begin{enumerate}
\item[(a)] the parameter $\lambda$ (to 4 decimal places);
\item[(b)] the mean waiting time;
\item[(c)] $P(T \le 5)$.
\end{enumerate}
\tcblower
\textbf{Solution.}
\begin{enumerate}
\item[(a)] $P(T > 10) = e^{-10\lambda} = 0.6$, so $-10\lambda = \ln 0.6$ and
\[ \lambda = -\frac{\ln 0.6}{10} = \frac{0.5108}{10} \approx 0.0511 \text{ per minute}. \]
\item[(b)] $E(T) = \dfrac{1}{\lambda} = \dfrac{10}{-\ln 0.6} \approx 19.58$ minutes.
\item[(c)] $P(T \le 5) = 1 - e^{-5\lambda} = 1 - e^{-0.2554} \approx 1 - 0.7746 = 0.2254$.
\end{enumerate}
\end{exoresolu}

\begin{aretenir}[Key points on the exponential distribution]
\begin{itemize}
\item $X \sim \mathrm{Exp}(\lambda)$ models the \cle{waiting time} between events of a Poisson process of rate $\lambda$; pdf $f(x) = \lambda e^{-\lambda x}$ for $x \ge 0$.
\item $F(x) = 1 - e^{-\lambda x}$ and $P(X > x) = e^{-\lambda x}$.
\item $P(a \le X \le b) = e^{-\lambda a} - e^{-\lambda b}$.
\item $E(X) = \dfrac{1}{\lambda}$, \quad $\mathrm{Var}(X) = \dfrac{1}{\lambda^2}$, \quad $\sigma = \dfrac{1}{\lambda}$.
\item Memoryless: $P(X > s+t \mid X > s) = P(X > t)$.
\item Median $m = \dfrac{\ln 2}{\lambda}$; the mode is $0$ (the density is decreasing).
\end{itemize}
\end{aretenir}

\begin{exercice}[Calls at a call centre]
Calls arrive at a call centre in Douala at an average rate of 4 per minute, according to a Poisson process. Let $T$ be the waiting time, in minutes, between two consecutive calls.
\begin{enumerate}
\item[(a)] Write down the distribution of $T$ and its pdf.
\item[(b)] Calculate the probability of waiting more than 30 seconds between two calls.
\item[(c)] Calculate the probability of waiting between 10 and 20 seconds.
\item[(d)] Find the median waiting time, i.e. the value $m$ such that $P(T > m) = 0.5$.
\end{enumerate}
\end{exercice}

\begin{corrige}[Exercise 1]
\begin{enumerate}
\item[(a)] $T \sim \mathrm{Exp}(4)$ with $f(t) = 4e^{-4t}$, $t \ge 0$ ($t$ in minutes).
\item[(b)] 30 seconds $= 0.5$ min: $P(T > 0.5) = e^{-4 \times 0.5} = e^{-2} \approx 0.1353$.
\item[(c)] 10 s $= \tfrac{1}{6}$ min, 20 s $= \tfrac{1}{3}$ min:
\[ P\!\left(\tfrac{1}{6} \le T \le \tfrac{1}{3}\right) = e^{-4/6} - e^{-4/3} = e^{-2/3} - e^{-4/3} \approx 0.5134 - 0.2636 = 0.2498. \]
\item[(d)] $e^{-4m} = 0.5 \Rightarrow -4m = \ln 0.5 \Rightarrow m = \dfrac{\ln 2}{4} \approx 0.1733$ min $\approx 10.4$ seconds.
\end{enumerate}
\end{corrige}

\begin{exercice}[Lifetime of a phone battery]
The lifetime $X$ (in years) of a type of phone battery is modelled by $X \sim \mathrm{Exp}(0.5)$.
\begin{enumerate}
\item[(a)] Find the mean and the standard deviation of $X$.
\item[(b)] Find the probability that a battery lasts more than 3 years.
\item[(c)] A battery has already worked for 2 years. Find the probability that it works for at least one more year.
\item[(d)] Find the value $k$ such that $P(X \le k) = 0.9$.
\end{enumerate}
\end{exercice}

\begin{corrige}[Exercise 2]
\begin{enumerate}
\item[(a)] $E(X) = \dfrac{1}{0.5} = 2$ years; $\sigma = \dfrac{1}{0.5} = 2$ years (and $\mathrm{Var}(X) = 4$).
\item[(b)] $P(X > 3) = e^{-0.5 \times 3} = e^{-1.5} \approx 0.2231$.
\item[(c)] By the memoryless property, $P(X > 3 \mid X > 2) = P(X > 1) = e^{-0.5} \approx 0.6065$.
\item[(d)] $1 - e^{-0.5k} = 0.9 \Rightarrow e^{-0.5k} = 0.1 \Rightarrow k = \dfrac{\ln 10}{0.5} = 2\ln 10 \approx 4.61$ years.
\end{enumerate}
\end{corrige}

\begin{savaistu}[Why ``memoryless'' matters in engineering]
The exponential distribution is the \emph{only} continuous distribution with the memoryless property. Engineers use it to model components whose failure is due to sudden external shocks rather than gradual wear — such as a fuse blown by a power surge. For components that age (tyres, batteries), other models are preferred, because an old tyre is definitely \emph{not} as good as new!
\end{savaistu}

\coursec{The Uniform (Rectangular) Distribution}

In the previous section we met the \cle{exponential distribution}, a continuous model whose probability density decreases as $x$ increases: short waiting times are more likely than long ones. But not every continuous situation behaves that way. Imagine you arrive at a bus stop in Yaound\'e knowing only that the bus will come \emph{at some moment within the next 15 minutes}, with no reason to expect one moment more than another. Or think of a calculator that rounds every amount to the nearest franc: the rounding error is just as likely to be $0.1$ franc as $0.4$ franc. In such situations every value in an interval is \emph{equally likely}. The model for this is the \cle{uniform distribution}, the simplest of all continuous distributions, and the subject of this section.

\courssub{Situations Leading to the Uniform Model}

Before any formula, let us collect concrete situations where the uniform model is the natural choice.

\begin{itemize}
\item \textbf{A random number in an interval.} A computer (or a fair spinner) generates a number $X$ ``at random'' between $a$ and $b$. No sub-interval of a given length is favoured over another of the same length.
\item \textbf{Rounding errors.} When a measurement is rounded to the nearest unit (nearest franc, nearest millimetre, nearest minute), the true value differs from the recorded value by an error $X$ which lies uniformly between $-0.5$ and $+0.5$ of the unit.
\item \textbf{Arrival within a known window.} A bus is known to arrive at some instant between 8:00 and 8:15, but nothing more precise is known. The arrival time $X$ is uniform on $[0,15]$ (measuring time in minutes after 8:00).
\item \textbf{The position of a defect.} A flaw occurs at a uniformly random position along a cable of length $L$.
\end{itemize}

In each case the key phrase is: \emph{all values in the interval are equally likely}. For a continuous variable this cannot mean that each single value has a positive probability (we know $P(X=x)=0$ for continuous variables); it means that the \cle{probability density function} is \emph{constant} on the interval.

\courssub{Definition and Probability Density Function}

\begin{definition}[The uniform distribution $U(a,b)$]
A continuous random variable $X$ has the \cle{uniform distribution} on the interval $[a,b]$ (with $a<b$), written $X\sim U(a,b)$, if its probability density function is
\[
f(x)=\begin{cases} \dfrac{1}{b-a} & \text{if } a\le x\le b,\\[6pt] 0 & \text{otherwise.}\end{cases}
\]
Because the graph of $f$ is a horizontal segment over $[a,b]$, the region under the curve is a \emph{rectangle}: this is why the distribution is also called the \cle{rectangular distribution}.
\end{definition}

Why must the constant height be exactly $\dfrac{1}{b-a}$? Because the total area under any pdf must equal $1$. The region under the graph is a rectangle of base $b-a$ and height $h$, so
\[
\text{Area}=(b-a)\times h=1 \quad\Longrightarrow\quad h=\frac{1}{b-a}.
\]
We can also verify this by integration:
\[
\int_{-\infty}^{+\infty} f(x)\,dx=\int_a^b \frac{1}{b-a}\,dx=\frac{1}{b-a}\big[x\big]_a^b=\frac{b-a}{b-a}=1. \checkmark
\]

\begin{popfigure}
\begin{center}
\begin{tikzpicture}
\begin{axis}[
    width=11cm, height=6cm,
    axis lines=middle,
    xlabel={$x$}, ylabel={$f(x)$},
    xmin=0, xmax=8, ymin=0, ymax=0.45,
    xtick={2,3,5,6}, ytick={0.25},
    yticklabels={$\frac{1}{4}$},
    every axis x label/.style={at={(ticklabel* cs:1.02)}, anchor=west},
    every axis y label/.style={at={(ticklabel* cs:1.05)}, anchor=south},
]
\addplot [fill=poporange!45, draw=none] coordinates {(3,0) (3,0.25) (5,0.25) (5,0)} \closedcycle;
\addplot [very thick, popblue] coordinates {(0,0) (2,0)};
\addplot [very thick, popblue] coordinates {(2,0.25) (6,0.25)};
\addplot [very thick, popblue] coordinates {(6,0) (8,0)};
\addplot [dashed, popblue] coordinates {(2,0) (2,0.25)};
\addplot [dashed, popblue] coordinates {(6,0) (6,0.25)};
\addplot [dashed, popdark] coordinates {(3,0) (3,0.25)};
\addplot [dashed, popdark] coordinates {(5,0) (5,0.25)};
\node[popdark, align=center] at (axis cs:4,0.13) {\small $P(3\le X\le 5)$};
\node[popblue] at (axis cs:4,0.32) {\small $f(x)=\dfrac{1}{4}$};
\end{axis}
\end{tikzpicture}
\legende{The pdf of $X\sim U(2,6)$. The shaded rectangle has area $P(3\le X\le 5)=\dfrac{5-3}{6-2}=\dfrac{1}{2}$.}
\end{center}
\end{popfigure}

\courssub{Probabilities as Ratios of Lengths}

Since the density is constant, every probability is the area of a rectangle, and no integration is ever needed in practice.

\begin{propriete}[Probabilities for $U(a,b)$]
If $X\sim U(a,b)$ and $[c,d]\subseteq[a,b]$, then
\[
P(c\le X\le d)=\int_c^d \frac{1}{b-a}\,dx=\frac{d-c}{b-a}.
\]
In words: \emph{probability $=$ (length of the favourable part) $\div$ (total length)}. (The proof by integration, shown here, is examinable but immediate.)
\end{propriete}

\begin{methode}[Computing a uniform probability]
\begin{enumerate}
\item Identify the interval $[a,b]$ and check that $X\sim U(a,b)$.
\item Identify the favourable interval $[c,d]$ and \textbf{truncate it to $[a,b]$} if it sticks out.
\item Compute $P=\dfrac{\text{length of favourable part}}{b-a}$.
\end{enumerate}
\end{methode}

\begin{attention}[Check the position of your bounds]
The formula $\dfrac{d-c}{b-a}$ is only valid when $[c,d]\subseteq[a,b]$. Remember that
\[
P(X\le x)=0 \ \text{for } x<a, \qquad P(X\le x)=1 \ \text{for } x>b.
\]
So, for example, if $X\sim U(2,6)$, then $P(X\le 5)=P(2\le X\le 5)=\dfrac{3}{4}$, \emph{not} $\dfrac{5}{4}$: the part of $(-\infty,5]$ below $a=2$ carries no probability at all. Always clip your interval to $[a,b]$ first.
\end{attention}

\courssub{The Cumulative Distribution Function}

Integrating the pdf from $-\infty$ to $x$ gives the cumulative distribution function. For $x<a$ nothing has been accumulated; for $a\le x\le b$ the accumulated area is a rectangle of base $x-a$ and height $\frac{1}{b-a}$; for $x>b$ the whole area $1$ is accumulated.

\begin{propriete}[CDF of $U(a,b)$]
If $X\sim U(a,b)$, its cumulative distribution function is
\[
F(x)=P(X\le x)=\begin{cases} 0 & \text{if } x<a,\\[4pt] \dfrac{x-a}{b-a} & \text{if } a\le x\le b,\\[6pt] 1 & \text{if } x>b.\end{cases}
\]
\end{propriete}

\begin{popfigure}
\begin{center}
\begin{tikzpicture}[scale=1.05]
\draw[->, thick] (-0.5,0) -- (8,0) node[right] {$x$};
\draw[->, thick] (0,-0.4) -- (0,3) node[above] {$F(x)$};
\draw[very thick, popblue] (-0.5,0) -- (2,0);
\draw[very thick, popblue] (2,0) -- (6,2.4);
\draw[very thick, popblue] (6,2.4) -- (8,2.4);
\draw[dashed, popmuted] (6,0) -- (6,2.4) -- (0,2.4);
\node[below] at (2,0) {$a$};
\node[below] at (6,0) {$b$};
\node[left] at (0,2.4) {$1$};
\node[popblue, align=center] at (4.6,1.7) {\small $F(x)=\dfrac{x-a}{b-a}$};
\fill[popblue] (2,0) circle (1.6pt);
\fill[popblue] (6,2.4) circle (1.6pt);
\end{tikzpicture}
\legende{The cumulative distribution function of $U(a,b)$: flat at $0$ before $a$, linear with slope $\frac{1}{b-a}$ on $[a,b]$, flat at $1$ after $b$.}
\end{center}
\end{popfigure}

\courssub{Mean and Variance of the Uniform Distribution}

\textbf{The mean.} Intuitively, since the density is symmetric about the midpoint of $[a,b]$, the average value must be the midpoint $\dfrac{a+b}{2}$. Let us prove it from the definition $E(X)=\displaystyle\int_{-\infty}^{+\infty} x f(x)\,dx$:
\begin{align*}
E(X)&=\int_a^b x\cdot\frac{1}{b-a}\,dx=\frac{1}{b-a}\left[\frac{x^2}{2}\right]_a^b=\frac{b^2-a^2}{2(b-a)}=\frac{(b-a)(b+a)}{2(b-a)}=\frac{a+b}{2}.
\end{align*}

\textbf{The variance.} We use the computational formula $\operatorname{Var}(X)=E(X^2)-[E(X)]^2$. First,
\begin{align*}
E(X^2)&=\int_a^b x^2\cdot\frac{1}{b-a}\,dx=\frac{1}{b-a}\left[\frac{x^3}{3}\right]_a^b=\frac{b^3-a^3}{3(b-a)}.
\end{align*}
Since $b^3-a^3=(b-a)(b^2+ab+a^2)$, this simplifies to $E(X^2)=\dfrac{a^2+ab+b^2}{3}$. Therefore
\begin{align*}
\operatorname{Var}(X)&=\frac{a^2+ab+b^2}{3}-\left(\frac{a+b}{2}\right)^2=\frac{a^2+ab+b^2}{3}-\frac{a^2+2ab+b^2}{4}\\[2pt]
&=\frac{4(a^2+ab+b^2)-3(a^2+2ab+b^2)}{12}=\frac{a^2-2ab+b^2}{12}=\frac{(b-a)^2}{12}.
\end{align*}

\begin{propriete}[Mean and variance of $U(a,b)$]
If $X\sim U(a,b)$, then
\[
E(X)=\frac{a+b}{2} \qquad\text{and}\qquad \operatorname{Var}(X)=\frac{(b-a)^2}{12},
\qquad \sigma=\frac{b-a}{\sqrt{12}}=\frac{b-a}{2\sqrt{3}}.
\]
Both derivations above are examinable: know how to obtain them from $E(X)=\int x f(x)\,dx$ and $\operatorname{Var}(X)=E(X^2)-[E(X)]^2$.
\end{propriete}

\begin{aretenir}[Key facts about $U(a,b)$]
\begin{itemize}
\item pdf: $f(x)=\dfrac{1}{b-a}$ on $[a,b]$, $0$ elsewhere; the graph is a rectangle.
\item Probabilities are ratios of lengths: $P(c\le X\le d)=\dfrac{d-c}{b-a}$ for $[c,d]\subseteq[a,b]$.
\item $E(X)=\dfrac{a+b}{2}$ (the midpoint); $\operatorname{Var}(X)=\dfrac{(b-a)^2}{12}$.
\item The variance depends only on the \emph{width} of the interval, not on its position.
\end{itemize}
\end{aretenir}

\courssub{Worked Examples}

\begin{exoresolu}[Waiting for the bus]
A bus arrives at a stop at a time uniformly distributed between 8:00 and 8:15. Let $X$ be the waiting time, in minutes after 8:00, so $X\sim U(0,15)$.
\begin{enumerate}
\item Write down the pdf of $X$ and find $E(X)$ and $\operatorname{Var}(X)$.
\item Find the probability of waiting between 3 and 9 minutes.
\item Find the probability of waiting more than 10 minutes.
\item Find the time $t$ such that $P(X\le t)=0.75$.
\end{enumerate}
\tcblower
\textbf{Solution.}
\begin{enumerate}
\item $f(x)=\dfrac{1}{15-0}=\dfrac{1}{15}$ for $0\le x\le 15$, and $0$ otherwise.
\[
E(X)=\frac{0+15}{2}=7.5\ \text{min}, \qquad \operatorname{Var}(X)=\frac{(15-0)^2}{12}=\frac{225}{12}=18.75\ \text{min}^2.
\]
\item $[3,9]\subseteq[0,15]$, so $P(3\le X\le 9)=\dfrac{9-3}{15}=\dfrac{6}{15}=\dfrac{2}{5}=0.4$.
\item $P(X>10)=\dfrac{15-10}{15}=\dfrac{5}{15}=\dfrac{1}{3}\approx 0.333$.
\item We need $F(t)=\dfrac{t-0}{15}=0.75$, giving $t=11.25$ minutes, i.e.\ 11 min 15 s after 8:00. There is a $75\%$ chance the bus has arrived by then.
\end{enumerate}
\end{exoresolu}

\begin{exoresolu}[Rounding to the nearest franc]
A shop in Douala rounds every bill to the nearest franc. The rounding error $X$ (true amount minus rounded amount) is uniformly distributed on $[-0.5,\,0.5]$ francs.
\begin{enumerate}
\item Find the mean and standard deviation of $X$.
\item Find the probability that the error exceeds $0.3$ franc in absolute value.
\end{enumerate}
\tcblower
\textbf{Solution.} Here $a=-0.5$ and $b=0.5$, so $b-a=1$.
\begin{enumerate}
\item $E(X)=\dfrac{-0.5+0.5}{2}=0$ franc: rounding is fair on average.
\[
\operatorname{Var}(X)=\frac{1^2}{12}=\frac{1}{12}, \qquad \sigma=\frac{1}{\sqrt{12}}=\frac{1}{2\sqrt{3}}\approx 0.289\ \text{franc}.
\]
\item $P(|X|>0.3)=P(X>0.3)+P(X<-0.3)=\dfrac{0.5-0.3}{1}+\dfrac{-0.3-(-0.5)}{1}=0.2+0.2=0.4$.
\end{enumerate}
\end{exoresolu}

\courssub{Finding $a$ and $b$ from the Mean and Variance}

In examination questions the endpoints are sometimes unknown and must be recovered from given information about the mean and the variance (or standard deviation). The strategy is to solve the two equations simultaneously.

\begin{methode}[Recovering $a$ and $b$ from $E(X)$ and $\operatorname{Var}(X)$]
\begin{enumerate}
\item Write $\dfrac{a+b}{2}=\mu$ and $\dfrac{(b-a)^2}{12}=\sigma^2$.
\item From the second equation, $b-a=\sqrt{12\sigma^2}=2\sqrt{3}\,\sigma$ (take the positive root since $b>a$).
\item Solve the system $a+b=2\mu$, $b-a=2\sqrt{3}\,\sigma$ by adding and subtracting:
\[
a=\mu-\sqrt{3}\,\sigma, \qquad b=\mu+\sqrt{3}\,\sigma.
\]
\end{enumerate}
\end{methode}

\begin{exoresolu}[Finding the endpoints]
A continuous random variable $X$ is uniformly distributed on $[a,b]$. Given that $E(X)=4$ and $\operatorname{Var}(X)=3$, find $a$ and $b$, and hence compute $P(2\le X\le 5)$.
\tcblower
\textbf{Solution.} We set up the two equations:
\[
\frac{a+b}{2}=4 \quad\Longrightarrow\quad a+b=8; \qquad \frac{(b-a)^2}{12}=3 \quad\Longrightarrow\quad (b-a)^2=36 \quad\Longrightarrow\quad b-a=6
\]
(taking the positive square root because $b>a$). Adding the equations: $2b=14$, so $b=7$; subtracting: $2a=2$, so $a=1$. Hence $X\sim U(1,7)$ and
\[
P(2\le X\le 5)=\frac{5-2}{7-1}=\frac{3}{6}=\frac{1}{2}.
\]
\end{exoresolu}

\begin{attention}[Trap: the sign of the square root]
From $(b-a)^2=36$ you must take $b-a=+6$, not $-6$, because by definition $a<b$. Choosing the negative root swaps $a$ and $b$ and gives an invalid interval. Also, do not confuse $\operatorname{Var}(X)$ with $\sigma$: if the question gives the standard deviation, square it before using $\dfrac{(b-a)^2}{12}$.
\end{attention}

\begin{exercice}[Practice]
\begin{enumerate}
\item $X\sim U(3,11)$. Find (a) $E(X)$, (b) $\operatorname{Var}(X)$, (c) $P(X\le 5)$, (d) $P(X>13)$, (e) the value $m$ such that $P(X\le m)=\tfrac14$.
\item The error made when a stopwatch reading is rounded to the nearest second is uniform on $[-0.5,0.5]$ seconds. Find the probability that the error is within $0.2$ second of zero.
\item $X\sim U(a,b)$ with $E(X)=10$ and $\sigma=2\sqrt{3}$. Find $a$ and $b$.
\end{enumerate}
\end{exercice}

\begin{corrige}[Exercise 1]
\begin{enumerate}
\item (a) $E(X)=\dfrac{3+11}{2}=7$. (b) $\operatorname{Var}(X)=\dfrac{(11-3)^2}{12}=\dfrac{64}{12}=\dfrac{16}{3}$. (c) $P(X\le 5)=\dfrac{5-3}{8}=\dfrac14$. (d) $P(X>13)=0$ since $13>b=11$. (e) $\dfrac{m-3}{8}=\dfrac14\Rightarrow m-3=2\Rightarrow m=5$.
\item $P(|X|\le 0.2)=P(-0.2\le X\le 0.2)=\dfrac{0.2-(-0.2)}{1}=0.4$.
\item $b-a=2\sqrt{3}\,\sigma=2\sqrt{3}\times 2\sqrt{3}=12$ and $a+b=20$. Adding: $2b=32$, $b=16$; then $a=4$. So $X\sim U(4,16)$.
\end{enumerate}
\end{corrige}

\begin{savaistu}[Why ``rectangular''?]
The uniform distribution is the continuous analogue of throwing a fair die: a fair die gives equal probability $\tfrac16$ to each of six \emph{discrete} outcomes, while $U(a,b)$ spreads the probability evenly over a whole \emph{continuum} of values. Random-number generators in calculators and computers produce (approximately) $U(0,1)$ values, and every other continuous distribution used in simulation --- including the exponential distribution of the previous section and the normal distribution of the next --- is built by transforming these uniform numbers.
\end{savaistu}

The uniform distribution models complete indifference within an interval. But most real measurements --- heights, marks, masses --- cluster around a central value and become rarer as we move away from it. Modelling that bell-shaped behaviour is the work of the \emph{normal distribution}, which we study next.

\coursec{The Normal Distribution: pdf, Mean and Variance}

\courssub{Introduction: the need for a continuous model}

The uniform distribution of this section. In the previous section, we met the uniform distribution: a continuous uniform distribution. In the previous section, we saw the uniform distribution of the uniform distribution of the uniform distribution of the uniform distribution of the uniform distribution of the uniform distribution of the uniform distribution of the uniform distribution of the uniform distribution of the uniform distribution of the uniform distribution of the uniform distribution of the uniform distribution of the uniform distribution of the uniform distribution of the uniform distribution of the uniform distribution of the uniform distribution of the uniform distribution of the uniform distribution of the uniform distribution of the uniform distribution of the uniform distribution of the uniform distribution of the uniform distribution of the uniform distribution of the uniform distribution of the uniform distribution of the uniform distribution of the uniform distribution of the uniform distribution of the uniform distribution of

\coursec{The Standard Normal Variable and Use of Tables}

In the previous section we discovered the normal distribution $N(\mu,\sigma^2)$, with its bell-shaped density curve centred on the mean $\mu$ and spread out by the standard deviation $\sigma$. We also saw that probabilities are \emph{areas under the curve}. But here is the practical difficulty: the density
\[
f(x)=\frac{1}{\sigma\sqrt{2\pi}}\,e^{-\frac{(x-\mu)^2}{2\sigma^2}}
\]
has no elementary antiderivative. We cannot integrate it by hand. So how did statisticians compute normal probabilities before computers? The answer is one of the cleverest ideas in all of statistics: instead of building a table for \emph{every} possible $\mu$ and $\sigma$ (there are infinitely many!), we build \emph{one single table} for one special normal distribution, and we transform every other normal variable into that special one. That special distribution is the \cle{standard normal distribution}, and the transformation is called \cle{standardisation}. This section is the computational heart of the whole chapter: nearly every GCE question on the normal distribution reduces to the techniques you are about to learn.

\courssub{The Standard Normal Distribution}

\begin{definition}[Standard normal distribution and the function $\Phi$]
The \cle{standard normal distribution} is the normal distribution with mean $0$ and variance $1$. A variable following it is usually denoted $Z$, and we write
\[
Z\sim N(0,1).
\]
Its probability density function is
\[
\varphi(z)=\frac{1}{\sqrt{2\pi}}\,e^{-\frac{z^2}{2}},\qquad z\in\mathbb{R}.
\]
Its cumulative distribution function is denoted by the capital Greek letter \cle{Phi}:
\[
\Phi(z)=P(Z\le z)=\int_{-\infty}^{z}\frac{1}{\sqrt{2\pi}}\,e^{-\frac{t^2}{2}}\,dt.
\]
\end{definition}

Geometrically, $\Phi(z)$ is simply the \emph{area under the standard normal curve to the left of} $z$. Since the total area under the curve is $1$, we always have $0<\Phi(z)<1$, with $\Phi(z)\to 0$ as $z\to-\infty$ and $\Phi(z)\to 1$ as $z\to+\infty$.

\begin{popfigure}
\begin{center}
\begin{tikzpicture}
\begin{axis}[
    width=12cm, height=6.2cm,
    axis lines=middle,
    xlabel={$z$}, ylabel={$\varphi(z)$},
    xmin=-3.6, xmax=3.6, ymin=0, ymax=0.45,
    xtick={-3,-2,-1,0,1,1.26,2,3},
    xticklabels={$-3$,$-2$,$-1$,$0$,$1$,$z$,$2$,$3$},
    ytick=\empty,
    enlargelimits=false,
    clip=false
]
\addplot[fill=popblueL, draw=none, domain=-3.6:1.26, samples=80] {0.3989*exp(-x^2/2)} \closedcycle;
\addplot[thick, popblue, domain=-3.6:3.6, samples=120] {0.3989*exp(-x^2/2)};
\draw[dashed, popdark] (axis cs:1.26,0) -- (axis cs:1.26,0.18);
\node[popdark, align=center] at (axis cs:-0.7,0.10) {$\Phi(z)=P(Z\le z)$};
\end{axis}
\end{tikzpicture}
\end{center}
\legende{The standard normal curve. The shaded left tail is the cumulative probability $\Phi(z)=P(Z\le z)$.}
\end{popfigure}

\begin{propriete}[Elementary properties of $\Phi$]
For the standard normal variable $Z\sim N(0,1)$:
\begin{enumerate}
\item $\Phi(0)=0.5$ (the curve is symmetric about $0$, so half the area lies to the left of $0$);
\item $\Phi(+\infty)=1$ and $\Phi(-\infty)=0$;
\item $P(Z=z)=0$ for every single value $z$ (a continuous variable has zero probability of hitting an exact point), so $P(Z\le z)=P(Z<z)$ and $P(Z\ge z)=P(Z>z)$.
\end{enumerate}
\end{propriete}

\courssub{The Standardisation Theorem}

Why is $N(0,1)$ enough for everything? Because of the following fundamental result.

\begin{propriete}[Standardisation theorem — proof instructive]
If $X\sim N(\mu,\sigma^2)$, then the variable
\[
Z=\frac{X-\mu}{\sigma}\quad\text{satisfies}\quad Z\sim N(0,1).
\]
Conversely, $X=\mu+\sigma Z$.
\end{propriete}

\textbf{Justification.} Subtracting $\mu$ shifts the centre of the distribution from $\mu$ to $0$; dividing by $\sigma$ rescales the horizontal axis so that the spread becomes $1$. More formally, a linear transformation $aX+b$ of a normal variable is still normal, with $E(aX+b)=a\mu+b$ and $\operatorname{Var}(aX+b)=a^2\sigma^2$. Taking $a=\frac{1}{\sigma}$ and $b=-\frac{\mu}{\sigma}$:
\[
E(Z)=\frac{1}{\sigma}\,\mu-\frac{\mu}{\sigma}=0,\qquad \operatorname{Var}(Z)=\frac{1}{\sigma^2}\cdot\sigma^2=1.
\]
Hence $Z$ is normal with mean $0$ and variance $1$, i.e. $Z\sim N(0,1)$. $\blacksquare$

\begin{popfigure}
\begin{center}
\begin{tikzpicture}[scale=1.0]
% Top axis: X
\draw[->, thick, popdark] (0,2.6) -- (10,2.6) node[right] {$X\sim N(\mu,\sigma^2)$};
\foreach \x/\lab in {2/{$a$}, 5/{$\mu$}, 8/{$b$}} {
  \draw[thick, popdark] (\x,2.45) -- (\x,2.75);
  \node[above, popdark] at (\x,2.75) {\lab};
}
\draw[very thick, poporange] (2,2.6) -- (8,2.6);
% Bottom axis: Z
\draw[->, thick, popblue] (0,0) -- (10,0) node[right] {$Z\sim N(0,1)$};
\foreach \x/\lab in {2/{$z_1$}, 5/{$0$}, 8/{$z_2$}} {
  \draw[thick, popblue] (\x,-0.15) -- (\x,0.15);
  \node[below, popblue] at (\x,-0.15) {\lab};
}
\draw[very thick, popblue] (2,0) -- (8,0);
% Arrows
\draw[->, thick, poppurple] (2,2.35) -- (2,0.25);
\draw[->, thick, poppurple] (8,2.35) -- (8,0.25);
\node[poppurple, align=center, text width=4.2cm] at (5,1.3) {standardise: $z=\dfrac{x-\mu}{\sigma}$};
\end{tikzpicture}
\end{center}
\legende{Standardisation maps the interval $[a,b]$ under $X$ to the interval $[z_1,z_2]$ under $Z$, where $z_1=\frac{a-\mu}{\sigma}$ and $z_2=\frac{b-\mu}{\sigma}$. The two intervals carry the same probability.}
\end{popfigure}

The practical consequence is enormous: every probability question about \emph{any} normal variable $X$ becomes a question about $Z$, whose cumulative probabilities are tabulated.

\begin{methode}[Computing normal probabilities by standardisation]
\begin{enumerate}
\item Identify $\mu$ and $\sigma$ (take the square root if the variance is given).
\item Rewrite the event in terms of $X$, then standardise each bound: replace $x$ by $z=\dfrac{x-\mu}{\sigma}$.
\item Use the table to find $\Phi(z)$:
\begin{itemize}
\item $P(X\le x)=\Phi\!\left(\dfrac{x-\mu}{\sigma}\right)$;
\item $P(X\ge x)=1-\Phi\!\left(\dfrac{x-\mu}{\sigma}\right)$;
\item $P(a\le X\le b)=\Phi\!\left(\dfrac{b-\mu}{\sigma}\right)-\Phi\!\left(\dfrac{a-\mu}{\sigma}\right)$.
\end{itemize}
\item Round the final probability to 4 decimal places, as in the table.
\end{enumerate}
\end{methode}

\courssub{Reading the Standard Normal Table}

The GCE standard normal table gives values of $\Phi(z)$ for $z\ge 0$, to \textbf{two decimal places} in $z$ and \textbf{four decimal places} in $\Phi(z)$. The table is organised as follows:

\begin{itemize}
\item the \textbf{row} gives $z$ to one decimal place (e.g. $1.2$);
\item the \textbf{column} gives the second decimal digit (e.g. $0.06$);
\item the entry at the intersection is $\Phi(z)$.
\end{itemize}

\begin{exemplebox}[Reading $\Phi(1.26)$]
To find $\Phi(1.26)$: go down the left column to the row $1.2$, then across to the column headed $0.06$ (the second decimal digit). The entry is $0.8962$. Hence
\[
\Phi(1.26)=P(Z\le 1.26)=0.8962.
\]
\end{exemplebox}

\begin{popfigure}
\begin{center}
\begin{tabular}{|c|c|c|c|c|c|}
\hline
\rowcolor{popblueL} $z$ & $0.04$ & $0.05$ & $0.06$ & $0.07$ & $0.08$ \\
\hline
$1.1$ & $0.8729$ & $0.8749$ & $0.8770$ & $0.8790$ & $0.8810$ \\
\hline
\cellcolor{popblueL} $1.2$ & $0.8925$ & $0.8944$ & \cellcolor{popgreenL} $\mathbf{0.8962}$ & $0.8980$ & $0.8997$ \\
\hline
$1.3$ & $0.9099$ & $0.9115$ & $0.9131$ & $0.9147$ & $0.9162$ \\
\hline
\end{tabular}
\end{center}
\legende{Extract-style sketch of the standard normal table. Row $1.2$, column $0.06$ gives $\Phi(1.26)=0.8962$.}
\end{popfigure}

\courssub{Symmetry Relations: Dealing with Negative $z$ and Upper Tails}

The table only lists $z\ge 0$. For negative values and for upper-tail probabilities we exploit the symmetry of the curve about $0$.

\begin{propriete}[Symmetry relations]
For any real numbers $z$, $a$, $b$ with $a\le b$:
\[
\Phi(-z)=1-\Phi(z),\qquad P(Z>z)=1-\Phi(z),\qquad P(a\le Z\le b)=\Phi(b)-\Phi(a).
\]
\end{propriete}

\textbf{Why $\Phi(-z)=1-\Phi(z)$?} By symmetry of the curve about the vertical axis, the area to the left of $-z$ equals the area to the right of $z$:
\[
\Phi(-z)=P(Z\le -z)=P(Z\ge z)=1-P(Z<z)=1-\Phi(z).
\]

\begin{attention}[The most frequent error with $\Phi$]
Never write $\Phi(-z)=-\Phi(z)$. Probabilities cannot be negative! The correct relation is
\[
\Phi(-z)=1-\Phi(z).
\]
For example $\Phi(-1.26)=1-\Phi(1.26)=1-0.8962=0.1038$, not $-0.8962$.
\end{attention}

\begin{exoresolu}[Direct probability calculations for $Z$]
Given $Z\sim N(0,1)$, find (a) $P(Z\le 1.5)$; (b) $P(Z>0.84)$; (c) $P(Z\le -2)$; (d) $P(-1.5\le Z\le 2)$; (e) $P(|Z|\le 1.96)$.
\tcblower
\textbf{Solution.}
\begin{enumerate}
\item $P(Z\le 1.5)=\Phi(1.5)=0.9332$.
\item $P(Z>0.84)=1-\Phi(0.84)=1-0.7995=0.2005$.
\item $P(Z\le -2)=\Phi(-2)=1-\Phi(2)=1-0.9772=0.0228$.
\item $P(-1.5\le Z\le 2)=\Phi(2)-\Phi(-1.5)=\Phi(2)-\bigl(1-\Phi(1.5)\bigr)=0.9772-(1-0.9332)=0.9104$.
\item $P(|Z|\le 1.96)=P(-1.96\le Z\le 1.96)=\Phi(1.96)-\Phi(-1.96)=2\Phi(1.96)-1=2(0.9750)-1=0.9500$.
\end{enumerate}
Note the useful shortcut from part (e): $P(|Z|\le z)=2\Phi(z)-1$. This is why $95\%$ of a normal distribution lies within $1.96$ standard deviations of the mean.
\end{exoresolu}

\courssub{Worked Examples with Real Data}

\begin{exoresolu}[Marks in a mock examination]
The marks obtained by Upper Sixth students in a Further Mathematics mock examination in Yaoundé are modelled by $X\sim N(58, 12^2)$. A student is chosen at random. Find the probability that the student scored (a) at most $70$ marks; (b) at least $40$ marks; (c) between $50$ and $70$ marks.
\tcblower
\textbf{Solution.} Here $\mu=58$ and $\sigma=12$. Standardise with $Z=\dfrac{X-58}{12}$.
\begin{enumerate}
\item $P(X\le 70)=P\!\left(Z\le \dfrac{70-58}{12}\right)=P(Z\le 1)=\Phi(1)=0.8413$.
\item $P(X\ge 40)=P\!\left(Z\ge \dfrac{40-58}{12}\right)=P(Z\ge -1.5)=P(Z\le 1.5)=\Phi(1.5)=0.9332$.
\item $P(50\le X\le 70)=P\!\left(\dfrac{50-58}{12}\le Z\le \dfrac{70-58}{12}\right)=P(-0.67\le Z\le 1)$
\[
=\Phi(1)-\Phi(-0.67)=\Phi(1)-\bigl(1-\Phi(0.67)\bigr)=0.8413-(1-0.7486)=0.5899.
\]
\end{enumerate}
So about $84\%$ of students scored at most $70$, about $93\%$ scored at least $40$, and about $59\%$ scored between $50$ and $70$.
\end{exoresolu}

\begin{exoresolu}[Masses of bags of cement]
The mass of a bag of cement produced by a factory in Douala is normally distributed with mean $50$ kg and standard deviation $0.4$ kg. A bag is rejected if its mass is below $49.2$ kg. Find the proportion of bags rejected.
\tcblower
\textbf{Solution.} Let $M\sim N(50, 0.4^2)$. Standardise the bound $49.2$:
\[
z=\frac{49.2-50}{0.4}=-2.
\]
Hence
\[
P(M<49.2)=P(Z<-2)=\Phi(-2)=1-\Phi(2)=1-0.9772=0.0228.
\]
About $2.28\%$ of the bags are rejected.
\end{exoresolu}

\courssub{Inverse Problems of Type 1: Finding a Value from a Probability}

So far we were given a value $x$ and asked for a probability. Inverse problems reverse the question: we are \emph{given the probability} $p$ and asked for the value $x$ such that $P(X\le x)=p$.

\begin{methode}[Inverse problem: from probability to value]
\begin{enumerate}
\item Write the statement in the form $P(X\le x)=p$ (use $1-\Phi$ first if the given event is of the type $P(X\ge x)$).
\item Read the table \textbf{backwards}: look for $p$ (or the closest entry) in the \emph{body} of the table and read off $z$ from the row and column headings. If $p<0.5$, then $z$ is negative: find $z'$ with $\Phi(z')=1-p$ and set $z=-z'$.
\item Un-standardise: $x=\mu+\sigma z$.
\end{enumerate}
\end{methode}

\begin{exoresolu}[Finding a diameter exceeded by $5\%$ of items]
The diameters of ball-bearings made by a machine are distributed as $D\sim N(25, 0.04^2)$ mm. Find the diameter $d$ exceeded by only $5\%$ of the ball-bearings.
\tcblower
\textbf{Solution.} ``Exceeded by $5\%$'' means $P(D>d)=0.05$, i.e. $P(D\le d)=0.95$. Standardising,
\[
P\!\left(Z\le \frac{d-25}{0.04}\right)=0.95 \quad\Longrightarrow\quad \Phi\!\left(\frac{d-25}{0.04}\right)=0.95.
\]
Reading the table backwards, the entry closest to $0.9500$ corresponds to $z=1.645$ (midway between $1.64$ giving $0.9495$ and $1.65$ giving $0.9505$). Therefore
\[
\frac{d-25}{0.04}=1.645 \quad\Longrightarrow\quad d=25+0.04(1.645)=25.0658\approx 25.07\text{ mm}.
\]
Only $5\%$ of the ball-bearings have a diameter greater than about $25.07$ mm.
\end{exoresolu}

\courssub{Inverse Problems of Type 2: Finding an Unknown $\mu$ or $\sigma$}

Sometimes the mean or the standard deviation (or both) is unknown, but we are given one or two probability statements. Each probability statement, after standardisation, produces \emph{one equation} involving $\mu$ and $\sigma$.

\begin{attention}[Two probabilities $\Rightarrow$ two equations]
When both $\mu$ and $\sigma$ are unknown, a single probability statement is not enough: one equation cannot determine two unknowns. You need \textbf{two} probability statements, giving \textbf{two simultaneous equations} in $\mu$ and $\sigma$. Solve them by elimination or substitution.
\end{attention}

\begin{exoresolu}[Determining $\mu$ and $\sigma$ from two probabilities]
The masses of loaves of bread sold by a bakery are normally distributed with unknown mean $\mu$ grams and unknown standard deviation $\sigma$ grams. It is known that $10\%$ of the loaves have mass less than $780$ g and $20\%$ have mass greater than $830$ g. Find $\mu$ and $\sigma$.
\tcblower
\textbf{Solution.} Translate each statement into an equation in $Z$.
\[
P(X<780)=0.10 \;\Longrightarrow\; \Phi\!\left(\frac{780-\mu}{\sigma}\right)=0.10.
\]
Since $0.10<0.5$, the $z$-value is negative. From the table, $\Phi(1.2816)\approx 0.90$, so $\Phi(-1.2816)=0.10$. Hence
\[
\frac{780-\mu}{\sigma}=-1.2816 \quad\Longrightarrow\quad \mu-1.2816\sigma=780. \tag{1}
\]
Similarly,
\[
P(X>830)=0.20 \;\Longrightarrow\; P(X\le 830)=0.80 \;\Longrightarrow\; \Phi\!\left(\frac{830-\mu}{\sigma}\right)=0.80.
\]
From the table, $\Phi(0.8416)\approx 0.80$, so
\[
\frac{830-\mu}{\sigma}=0.8416 \quad\Longrightarrow\quad \mu+0.8416\sigma=830. \tag{2}
\]
Subtracting (1) from (2):
\[
(0.8416+1.2816)\sigma=50 \;\Longrightarrow\; 2.1232\,\sigma=50 \;\Longrightarrow\; \sigma\approx 23.55.
\]
Substituting back into (2): $\mu=830-0.8416(23.55)\approx 810.2$.

Hence the mean mass is about $\mu\approx 810$ g and the standard deviation $\sigma\approx 23.6$ g.
\end{exoresolu}

\courssub{Linear Interpolation Between Table Entries (Useful Extra)}

The table gives $z$ only to two decimal places. If a calculation produces, say, $z=1.645$, the honest table reading is ambiguous between $\Phi(1.64)=0.9495$ and $\Phi(1.65)=0.9505$. \cle{Linear interpolation} estimates the value in between by assuming the curve is approximately straight over the small interval:
\[
\Phi(1.645)\approx 0.9495+\frac{1.645-1.64}{1.65-1.64}\times(0.9505-0.9495)=0.9495+0.5\times 0.0010=0.9500.
\]
In the GCE, examiners usually accept the nearest table entry, but interpolation gives extra accuracy and shows real mastery. The same idea works in reverse: to find $z$ from $\Phi(z)=p$ when $p$ falls strictly between two entries, interpolate between the corresponding $z$-values.

\begin{aretenir}[Key points of this section]
\begin{itemize}
\item $Z\sim N(0,1)$ is the standard normal variable; $\Phi(z)=P(Z\le z)$ is tabulated for $z\ge 0$.
\item Standardisation: if $X\sim N(\mu,\sigma^2)$ then $Z=\dfrac{X-\mu}{\sigma}\sim N(0,1)$, and $x=\mu+\sigma z$.
\item Core formulas: $P(X\le x)=\Phi\!\left(\frac{x-\mu}{\sigma}\right)$; $P(X\ge x)=1-\Phi\!\left(\frac{x-\mu}{\sigma}\right)$; $P(a\le X\le b)=\Phi(z_b)-\Phi(z_a)$.
\item Symmetry: $\Phi(-z)=1-\Phi(z)$ — never $-\Phi(z)$.
\item Inverse problems: read the table backwards to get $z$ from $p$, then $x=\mu+\sigma z$; two probability statements give two simultaneous equations for unknown $\mu$ and $\sigma$.
\end{itemize}
\end{aretenir}

\begin{exercice}[Practice]
\begin{enumerate}
\item Given $Z\sim N(0,1)$, find (a) $P(Z\le 2.35)$; (b) $P(Z>-1.04)$; (c) $P(0.5\le Z\le 1.87)$; (d) $P(|Z|>1.5)$.
\item The heights of students in a school in Bamenda are modelled by $H\sim N(168, 6^2)$ cm. Find the probability that a randomly chosen student is (a) taller than $180$ cm; (b) between $160$ cm and $175$ cm.
\item For $X\sim N(70, 25)$, find the value $x$ such that (a) $P(X\le x)=0.90$; (b) $P(X\ge x)=0.025$.
\item The lifetime $T$ (in hours) of an electric bulb is normal. Given $P(T<950)=0.1587$ and $P(T>1080)=0.0228$, find the mean and standard deviation of $T$.
\end{enumerate}
\end{exercice}

\begin{corrige}[Exercise 1]
\begin{enumerate}
\item (a) $\Phi(2.35)=0.9906$. (b) $P(Z>-1.04)=\Phi(1.04)=0.8508$. (c) $\Phi(1.87)-\Phi(0.5)=0.9693-0.6915=0.2778$. (d) $P(|Z|>1.5)=2\bigl(1-\Phi(1.5)\bigr)=2(1-0.9332)=0.1336$.
\item (a) $z=\frac{180-168}{6}=2$; $P(H>180)=1-\Phi(2)=0.0228$. (b) $z_1=\frac{160-168}{6}=-1.33$, $z_2=\frac{175-168}{6}=1.17$; probability $=\Phi(1.17)-\Phi(-1.33)=0.8790-(1-0.9082)=0.7872$.
\item (a) $\Phi(z)=0.90\Rightarrow z=1.2816$; $x=70+5(1.2816)=76.4$. (b) $P(X\le x)=0.975\Rightarrow z=1.96$; $x=70+5(1.96)=79.8$.
\item $\Phi(-1)=0.1587$ and $\Phi(2)=0.9772$, so $\frac{950-\mu}{\sigma}=-1$ and $\frac{1080-\mu}{\sigma}=2$. Subtracting: $3\sigma=130$, so $\sigma\approx 43.3$ hours and $\mu=950+43.3\approx 993.3$ hours.
\end{enumerate}
\end{corrige}

\begin{savaistu}[Why $1.96$?]
The value $z=1.96$ is famous worldwide: since $\Phi(1.96)=0.975$, exactly $95\%$ of any normal population lies within $1.96$ standard deviations of the mean. Doctors, engineers and pollsters use this every day when they build ``$95\%$ confidence intervals''. The number you just computed in this section is the backbone of modern data analysis — including the quality control of goods produced in Cameroonian factories.
\end{savaistu}

\coursec{Normal Approximations to the Binomial and Poisson Distributions}

In the previous section we mastered the standard normal variable $Z\sim N(0,1)$ and learned to read probabilities from the standard normal table. We now put that tool to work: when a discrete distribution (binomial or Poisson) involves a large parameter, exact calculations become prohibitively heavy. The \cle{Central Limit Theorem} guarantees that the shape of the distribution approaches a bell curve, so we can replace the discrete histogram by a continuous normal density — provided we apply a \cle{continuity correction} to account for the jump from integers to real numbers.

\begin{propriete}[Central Limit Theorem (Informal Statement)]
If $X_1, X_2, \dots, X_n$ are independent and identically distributed random variables with mean $\mu$ and variance $\sigma^2$, then the sum $S_n = \sum_{i=1}^n X_i$ (and the sample mean $\bar{X}$) is approximately normally distributed for large $n$:
\[
S_n \;\dot\sim\; N(n\mu,\, n\sigma^2), \qquad \bar{X} \;\dot\sim\; N\Bigl(\mu,\, \frac{\sigma^2}{n}\Bigr).
\]
Both the Binomial $B(n,p)$ (sum of $n$ Bernoulli trials) and the Poisson $Po(\lambda)$ (limit of Binomial, or sum of many rare events) fall under this theorem.
\end{propriete}

\courssub{Why Approximate?}

Consider $X\sim B(100,0.55)$. To find $P(X\ge 60)$ exactly we would need to sum 41 binomial terms:
\[
P(X\ge 60)=\sum_{k=60}^{100}\binom{100}{k}(0.55)^k(0.45)^{100-k}.
\]
Even with a calculator this is tedious. Similarly, for $X\sim Po(25)$, $P(X\le 20)=\sum_{k=0}^{20}e^{-25}\frac{25^k}{k!}$ requires 21 terms.  
The normal distribution offers a fast, accurate alternative \emph{when the parameters are large enough}. The approximations we study here are direct consequences of the Central Limit Theorem and are standard tools in the GCE Advanced Level examination.

\courssub{Normal Approximation to the Binomial Distribution}

\begin{propriete}[Normal Approximation to the Binomial]
Let $X\sim B(n,p)$ with $n$ large and $p$ not too close to $0$ or $1$.  
A common rule of thumb (examinable) is:
\[
np>5\quad\text{and}\quad n(1-p)>5.
\]
Then $X$ can be approximated by a normal distribution:
\[
X\;\dot\sim\; N\bigl(\mu=np,\;\sigma^2=np(1-p)\bigr).
\]
The symbol $\dot\sim$ means ``is approximately distributed as''.
\end{propriete}

\begin{attention}[Validity Conditions]
Always \textbf{state and check} the conditions $np>5$ and $n(1-p)>5$ before using the approximation. Failure to do so loses marks in the GCE.
\end{attention}

\begin{popfigure}
\begin{tikzpicture}
\begin{axis}[
    width=\linewidth, height=8cm,
    xmin=0, xmax=20,
    ymin=0, ymax=0.25,
    xtick={0,2,...,20},
    ytick={0,0.05,0.1,0.15,0.2,0.25},
    xlabel={$k$}, ylabel={Probability},
    axis lines=middle,
    enlargelimits=0.05,
    clip=false,
    legend style={at={(0.5,1.02)}, anchor=south, legend columns=2, font=\footnotesize}
]
% Binomial histogram B(20,0.5)
\addplot[ybar, bar width=0.35cm, fill=popblue!60, draw=popblue!80!black] coordinates {
(0,0.000001) (1,0.00002) (2,0.00018) (3,0.00109) (4,0.00462) (5,0.01479)
(6,0.03696) (7,0.07393) (8,0.12013) (9,0.16018) (10,0.17620)
(11,0.16018) (12,0.12013) (13,0.07393) (14,0.03696) (15,0.01479)
(16,0.00462) (17,0.00109) (18,0.00018) (19,0.00002) (20,0.000001)
};
% Normal curve N(10,5)
\addplot[domain=0:20, samples=200, thick, poporange, smooth] 
{1/sqrt(2*pi*5)*exp(-(x-10)^2/(2*5))};
\addlegendentry{$B(20,0.5)$ histogram}
\addlegendentry{$N(10,5)$ density}
\end{axis}
\end{tikzpicture}
\legende{Histogram of $B(20,0.5)$ with superimposed normal curve $N(10,5)$. The fit is already good at $n=20$.}
\end{popfigure}

\courssub{Continuity Correction for the Binomial}

Because the binomial variable takes only integer values while the normal variable is continuous, we must replace each integer $k$ by the interval $[k-0.5,\,k+0.5]$. The area under the normal curve over that interval approximates the height of the histogram bar at $k$.

\begin{methode}[Continuity Correction — Binomial]
To approximate $P(X \text{ relation } k)$ where $X\sim B(n,p)$ and $k$ is an integer:
\begin{enumerate}
\item Write the event in terms of $\le$ or $\ge$ (or $=$).
\item Replace the discrete bound by a continuous interval using $\pm0.5$:
\begin{itemize}
\item $P(X\le k) \;\approx\; P(Y\le k+0.5)$
\item $P(X< k) = P(X\le k-1) \;\approx\; P(Y\le k-0.5)$
\item $P(X\ge k) \;\approx\; P(Y\ge k-0.5)$
\item $P(X> k) = P(X\ge k+1) \;\approx\; P(Y\ge k+0.5)$
\item $P(X= k) \;\approx\; P(k-0.5 \le Y \le k+0.5)$
\end{itemize}
where $Y\sim N(np,\,np(1-p))$.
\item Standardise: $Z=\dfrac{Y-\mu}{\sigma}$ and use the standard normal table.
\end{enumerate}
\end{methode}

\begin{popfigure}
\begin{tikzpicture}[scale=1.3]
% Axes
\draw[->] (-0.5,0) -- (6.5,0) node[right] {$x$};
\draw[->] (0,-0.3) -- (0,1.6) node[above] {$f(x)$};
% Discrete bar at k=3
\draw[fill=popblue!60, draw=popblue!80!black] (2.5,0) rectangle (3.5,1.2);
\node[popblue!80!black, font=\small] at (3,1.35) {$P(X=3)$};
% Continuous normal curve (sketch, scaled to match bar height visually)
\draw[thick, poporange, domain=0:6, samples=100, smooth] plot (\x,{1.2*exp(-(\x-3)^2/1.5)});
\node[poporange, font=\small] at (5.5,1.1) {$N(\mu,\sigma^2)$};
% Continuity correction interval
\draw[<->, thick, popgreen!70!black] (2.5,0.1) -- (3.5,0.1);
\node[popgreen!70!black, above, font=\small] at (3,0.15) {$[2.5,\,3.5]$};
% Tick marks
\foreach \x in {2.5,3,3.5} \draw (\x,-0.05) -- (\x,0.05);
\node[below, font=\small] at (2.5,0) {$k-0.5$};
\node[below, font=\small] at (3,0) {$k$};
\node[below, font=\small] at (3.5,0) {$k+0.5$};
\end{tikzpicture}
\legende{Continuity correction: the discrete probability $P(X=k)$ (blue bar) is approximated by the area under the normal curve (orange) from $k-0.5$ to $k+0.5$ (green interval).}
\end{popfigure}

\begin{exoresolu}[Germination of Maize Seeds]
An agricultural station in Buea tests a new variety of maize with a germination probability of $0.55$. A farmer plants $100$ seeds.  
Use a normal approximation to find the probability that \textbf{at least 60} seeds germinate.
\tcblower
\textbf{Solution.}
Let $X$ be the number of seeds that germinate. $X\sim B(100,0.55)$.
\medskip
\textbf{1. Check conditions:}
$np = 100\times0.55 = 55 > 5$,\quad $n(1-p)=100\times0.45 = 45 > 5$. \quad $\checkmark$
\medskip
\textbf{2. Normal approximation:}
$X\;\dot\sim\; Y\sim N(\mu=55,\;\sigma^2=100\times0.55\times0.45=24.75)$,\quad $\sigma=\sqrt{24.75}\approx4.975$.
\medskip
\textbf{3. Continuity correction:}
``At least 60'' means $X\ge 60$.  
We use $P(X\ge 60)\approx P(Y\ge 59.5)$.
\medskip
\textbf{4. Standardise:}
\[
z = \frac{59.5 - 55}{\sqrt{24.75}} = \frac{4.5}{4.975} \approx 0.9045.
\]
\medskip
\textbf{5. Use standard normal table:}
$P(Y\ge 59.5) = P(Z\ge 0.9045) = 1 - \Phi(0.9045)$.
From tables, $\Phi(0.90)=0.8159$, $\Phi(0.91)=0.8186$. Interpolating: $\Phi(0.9045)\approx0.8171$.
\[
P(X\ge 60) \approx 1 - 0.8171 = 0.1829.
\]
\medskip
\textbf{Answer:} The approximate probability is $\mathbf{0.183}$ (3 s.f.).
\end{exoresolu}

\courssub{Normal Approximation to the Poisson Distribution}

\begin{propriete}[Normal Approximation to the Poisson]
Let $X\sim Po(\lambda)$ with $\lambda$ large. A common rule of thumb (examinable) is:
\[
\lambda > 10.
\]
Then $X$ can be approximated by a normal distribution:
\[
X\;\dot\sim\; N(\mu=\lambda,\;\sigma^2=\lambda).
\]
\end{propriete}

\begin{attention}[Validity Condition]
Always state $\lambda>10$ before approximating a Poisson by a normal. For $\lambda\le10$ the approximation is poor; use exact Poisson tables or the binomial approximation (if applicable).
\end{attention}

The continuity correction works \emph{exactly the same way} because the Poisson variable is also discrete (integer-valued).

\begin{methode}[Continuity Correction — Poisson]
For $X\sim Po(\lambda)$ approximated by $Y\sim N(\lambda,\lambda)$:
\begin{itemize}
\item $P(X\le k) \approx P(Y\le k+0.5)$
\item $P(X< k) \approx P(Y\le k-0.5)$
\item $P(X\ge k) \approx P(Y\ge k-0.5)$
\item $P(X> k) \approx P(Y\ge k+0.5)$
\item $P(X= k) \approx P(k-0.5\le Y\le k+0.5)$
\end{itemize}
Then standardise $Z=\dfrac{Y-\lambda}{\sqrt{\lambda}}$ and read from the normal table.
\end{methode}

\begin{exoresolu}[Call Centre — CAMTEL Douala]
A CAMTEL call centre in Douala receives an average of $25$ calls per hour. Assuming a Poisson model, use a normal approximation to find the probability that in a given hour the centre receives \textbf{at most 20} calls.
\tcblower
\textbf{Solution.}
Let $X$ be the number of calls in one hour. $X\sim Po(25)$.
\medskip
\textbf{1. Check condition:} $\lambda=25>10$. \quad $\checkmark$
\medskip
\textbf{2. Normal approximation:}
$X\;\dot\sim\; Y\sim N(\mu=25,\;\sigma^2=25)$, so $\sigma=5$.
\medskip
\textbf{3. Continuity correction:}
``At most 20'' means $X\le 20$.  
We use $P(X\le 20)\approx P(Y\le 20.5)$.
\medskip
\textbf{4. Standardise:}
\[
z = \frac{20.5 - 25}{5} = \frac{-4.5}{5} = -0.9.
\]
\medskip
\textbf{5. Use standard normal table:}
$P(Y\le 20.5) = P(Z\le -0.9) = \Phi(-0.9) = 1-\Phi(0.9)$.
From tables, $\Phi(0.9)=0.8159$.
\[
P(X\le 20) \approx 1 - 0.8159 = 0.1841.
\]
\medskip
\textbf{Answer:} The approximate probability is $\mathbf{0.184}$ (3 s.f.).
\end{exoresolu}

\courssub{Choosing the Correct Continuity Correction}

This is a classic GCE trap. The key is to rewrite the required probability in terms of $\le$ or $\ge$ \emph{before} applying the $\pm0.5$ shift.

\begin{aretenir}[Continuity Correction Quick Reference]
\begin{center}
\begin{tabularx}{\linewidth}{|l|X|X|}\hline
\textbf{Required probability} & \textbf{Rewrite as} & \textbf{Normal interval} \\ \hline
$P(X \le k)$ & $P(X \le k)$ & $P(Y \le k+0.5)$ \\
$P(X < k)$ & $P(X \le k-1)$ & $P(Y \le k-0.5)$ \\
$P(X \ge k)$ & $P(X \ge k)$ & $P(Y \ge k-0.5)$ \\
$P(X > k)$ & $P(X \ge k+1)$ & $P(Y \ge k+0.5)$ \\
$P(X = k)$ & $P(X = k)$ & $P(k-0.5 \le Y \le k+0.5)$ \\ \hline
\end{tabularx}
\end{center}
\end{aretenir}

\begin{attention}[Common Mistakes]
\begin{itemize}
\item Forgetting the continuity correction entirely (the \#1 cause of lost marks).
\item Applying the correction in the wrong direction (e.g. using $k+0.5$ for $X\ge k$).
\item Not rewriting strict inequalities ($<$, $>$) in terms of $\le$ or $\ge$ first.
\item Using the normal approximation without checking the validity conditions ($np>5$, $n(1-p)>5$, $\lambda>10$).
\end{itemize}
\end{attention}

\courssub{Summary of Approximations (Synthesis)}

You have now seen three approximations across the Further Mathematics course. The table below summarises them for revision.

\begin{center}
\begin{tabularx}{\linewidth}{|l|X|X|X|}\hline
\rowcolor{popblueL}
\textbf{Approximation} & \textbf{Original Distribution} & \textbf{Approximating Distribution} & \textbf{Validity Conditions} \\ \hline
Poisson $\to$ Binomial & $B(n,p)$ with $n$ large, $p$ small & $Po(\lambda=np)$ & $n\ge50$, $p\le0.1$, $np\le5$ (Lower Sixth) \\ \hline
Normal $\to$ Binomial & $B(n,p)$ & $N(np,\,np(1-p))$ & $np>5$ and $n(1-p)>5$ \\ \hline
Normal $\to$ Poisson & $Po(\lambda)$ & $N(\lambda,\,\lambda)$ & $\lambda>10$ \\ \hline
\end{tabularx}
\end{center}

\begin{savaistu}[Historical Note]
The normal approximation to the binomial was first derived by Abraham de Moivre in 1733 (published in \emph{The Doctrine of Chances}). He used it to avoid summing huge binomial coefficients — exactly the problem we face today! The continuity correction was later refined by Laplace and Gauss.
\end{savaistu}

\begin{aretenir}[Exam Tip — The Five-Line Chain]
In the GCE, present your solution as five clear lines:
\begin{enumerate}
\item \textbf{Distribution:} $X\sim B(n,p)$ or $X\sim Po(\lambda)$.
\item \textbf{Check conditions:} $np>5$, $n(1-p)>5$ or $\lambda>10$.
\item \textbf{Approximation:} $X\dot\sim Y\sim N(\mu,\sigma^2)$.
\item \textbf{Continuity correction:} Rewrite the probability and apply $\pm0.5$.
\item \textbf{Standardise \& table:} $Z=\frac{Y-\mu}{\sigma}$, read $\Phi(z)$.
\end{enumerate}
This structured approach guarantees you earn all method marks even if the final arithmetic has a slip.
\end{aretenir}

\begin{exercice}[Practice]
\begin{enumerate}
\item A factory in Douala produces light bulbs with a defect rate of $2\%$. In a batch of $500$ bulbs, use a normal approximation to find the probability that there are \textbf{more than 15} defective bulbs.
\item The number of accidents per month at a certain junction in Yaoundé follows a Poisson distribution with mean $12$. Use a normal approximation to find the probability that in a given month there are \textbf{between 8 and 15 accidents inclusive}.
\item A multiple-choice test has $80$ questions, each with $4$ options. A student guesses all answers. Use a normal approximation to find the probability that the student gets \textbf{at least 25} correct answers.
\end{enumerate}
\end{exercice}

\begin{corrige}[Exercise 1]
$X\sim B(500,0.02)$. $np=10>5$, $n(1-p)=490>5$.  
$X\dot\sim Y\sim N(10,\,9.8)$, $\sigma\approx3.1305$.  
$P(X>15)=P(X\ge16)\approx P(Y\ge15.5)$.  
$z=\frac{15.5-10}{3.1305}\approx1.757$. $P(Z\ge1.757)=1-\Phi(1.757)\approx1-0.9605=0.0395$.
\end{corrige}

\begin{corrige}[Exercise 2]
$X\sim Po(12)$. $\lambda=12>10$.  
$X\dot\sim Y\sim N(12,\,12)$, $\sigma=\sqrt{12}\approx3.464$.  
$P(8\le X\le15)\approx P(7.5\le Y\le15.5)$.  
$z_1=\frac{7.5-12}{3.464}\approx-1.30$, $z_2=\frac{15.5-12}{3.464}\approx1.01$.  
$P(-1.30\le Z\le1.01)=\Phi(1.01)-\Phi(-1.30)=0.8438-(1-0.9032)=0.7470$.
\end{corrige}

\begin{corrige}[Exercise 3]
$X\sim B(80,0.25)$. $np=20>5$, $n(1-p)=60>5$.  
$X\dot\sim Y\sim N(20,\,15)$, $\sigma=\sqrt{15}\approx3.873$.  
$P(X\ge25)\approx P(Y\ge24.5)$.  
$z=\frac{24.5-20}{3.873}\approx1.162$. $P(Z\ge1.162)=1-\Phi(1.162)\approx1-0.8770=0.1230$.
\end{corrige}

\newpage

\coursec{Integration activity}

\coursec{Integration Activity: Special Continuous Distributions}

\courssub{Aims of the activity}

This integration activity closes the chapter by bringing together, in a single realistic study, all the special continuous distributions you have met: the \cle{exponential distribution}, the \cle{uniform (rectangular) distribution}, the \cle{normal distribution}, and the \cle{normal approximation} to a discrete distribution with \cle{continuity correction}.

By the end of the activity, you should be able to:

\begin{itemize}
\item choose, and \emph{justify} the choice of, an appropriate probability model for a real situation;
\item use $P(X>x)=e^{-\lambda x}$ for an exponential model and interpret the parameter $\lambda$;
\item compute the mean $\dfrac{a+b}{2}$ and variance $\dfrac{(b-a)^2}{12}$ of a uniform distribution on $[a,b]$;
\item standardise a normal variable using $Z=\dfrac{X-\mu}{\sigma}$ and read probabilities from the standard normal table, including \emph{inverse} readings (finding a value from a probability);
\item check the validity conditions ($np>5$, $nq>5$; $\lambda>10$) before replacing a binomial or Poisson distribution by a normal one, and apply the continuity correction correctly;
\item communicate a complete statistical argument in a group presentation, as expected in the GCE Advanced Level examination.
\end{itemize}

\begin{methode}[How to work]
Form groups of three or four. Each group keeps a written record of: the model chosen, the reason for the choice, the calculations, and the conclusion stated in everyday language. Allow about 15 minutes per task, then 15 minutes for the group presentations. The standard normal table and a non-programmable calculator are allowed.
\end{methode}

\courssub{Situation}

The \cle{SONEL} water-pumping station of Bonabéri, in Douala, supplies a residential neighbourhood. The station manager, Mrs Etame, collects data in order to plan maintenance and billing. Her records give the following information.

\textbf{Breakdowns.} Breakdowns of the main pump occur \emph{at random}, independently of one another, at a constant average rate of \textbf{one breakdown every 40 days}. The time $T$ (in days) between two consecutive breakdowns may therefore be modelled by an exponential distribution.

\textbf{Repair crew.} When a breakdown is reported by telephone, the repair crew's arrival time $U$ (in hours) after the call is \emph{equally likely to be anywhere between 1 hour and 3 hours}. It may therefore be modelled by a uniform distribution on $[1,3]$.

\textbf{Consumption.} The daily water consumption $X$ of the neighbourhood (in $\mathrm{m^3}$) is known from long experience to be normally distributed with mean $\mu=850\ \mathrm{m^3}$ and standard deviation $\sigma=60\ \mathrm{m^3}$.

\textbf{Bill payment.} The neighbourhood contains $200$ households supplied by the station. Each household, independently of the others, pays its monthly bill on time with probability $0.6$. Let $Y$ be the number of households paying on time in a given month.

Mrs Etame asks your group of Upper Sixth Further Mathematics students to answer four questions that she needs for her annual report:

\begin{enumerate}
\item What is the probability that the pump runs for \emph{more than 60 days} without any breakdown?
\item What are the mean and the variance of the crew's arrival time?
\item What is the probability that the daily consumption exceeds $950\ \mathrm{m^3}$? Above what level is the consumption on only $5\%$ of days (the station must be able to supply this level)?
\item What is the probability that \emph{at least 110 households} pay their bill on time in a given month?
\end{enumerate}

\courssub{Tasks}

\textbf{Task 1 — Modelling the breakdowns (exponential distribution).}\par
\begin{enumerate}
\item Explain why the exponential distribution is a reasonable model for $T$. Which property of the breakdowns justifies this choice?
\item Show that the rate parameter is $\lambda=\dfrac{1}{40}=0.025$ per day, and write down the probability density function of $T$.
\item Compute $P(T>60)$. Give your answer to 3 decimal places and interpret it in one sentence for Mrs Etame.
\item Compute the mean time between breakdowns and check that it agrees with the data.
\end{enumerate}

\textbf{Task 2 — Modelling the repair crew (uniform distribution).}\par
\begin{enumerate}
\item Write down the probability density function of $U$ on $[1,3]$ and sketch its graph.
\item Compute $E(U)$ and $\mathrm{Var}(U)$ using the formulae for the uniform distribution.
\item Compute $P(U<1.5)$, the probability that the crew arrives within an hour and a half.
\end{enumerate}

\textbf{Task 3 — Modelling the consumption (normal distribution).}\par
\begin{enumerate}
\item Standardise $X$ and compute $P(X>950)$, giving your answer to 4 decimal places.
\item Let $c$ be the consumption level exceeded on only $5\%$ of days. Write the equation $P(X>c)=0.05$ in terms of the standard normal variable $Z$, find the corresponding value $z$ from the table (use $z\approx 1.645$), and deduce $c$ to the nearest $\mathrm{m^3}$.
\item Explain in one sentence why this value $c$ matters for the sizing of the pump.
\end{enumerate}

\textbf{Task 4 — Modelling bill payments (normal approximation to the binomial).}\par
\begin{enumerate}
\item State the exact distribution of $Y$ and its parameters.
\item Verify that the conditions for a normal approximation are satisfied, and give the mean and variance of the approximating normal distribution.
\item Using the continuity correction, compute an approximate value of $P(Y\geq 110)$.
\item Justify why the continuity correction is needed here, and say what would go wrong without it.
\end{enumerate}

\textbf{Task 5 — Group presentation.}\par
Prepare a five-minute presentation of your four models. For each one, state: the distribution chosen, the reason for the choice, the validity conditions checked, the numerical result, and a conclusion in plain language that Mrs Etame — who is not a mathematician — can understand.

\courssub{Model answer}

\textbf{Task 1.} Breakdowns occur at random, independently, at a constant average rate: these are exactly the conditions under which the waiting time between events follows an exponential distribution (the continuous analogue of the geometric distribution, with the memoryless property).

The mean of an exponential distribution is $\dfrac{1}{\lambda}$, and the mean time between breakdowns is $40$ days, so
\[
\frac{1}{\lambda}=40 \quad\Longrightarrow\quad \lambda=\frac{1}{40}=0.025\ \text{per day}.
\]
The probability density function is
\[
f(t)=0.025\,e^{-0.025t},\quad t\geq 0 \quad (\text{and }0\text{ otherwise}).
\]
Then
\[
P(T>60)=e^{-\lambda \times 60}=e^{-0.025\times 60}=e^{-1.5}\approx 0.223.
\]
\emph{Interpretation:} there is about a $22\%$ chance that the pump runs for more than 60 days without any breakdown. The mean is $E(T)=\dfrac{1}{\lambda}=40$ days, which matches the data, as expected.

\textbf{Task 2.} The arrival time is uniform on $[1,3]$, so
\[
f(u)=\frac{1}{3-1}=\frac{1}{2},\quad 1\leq u\leq 3,
\]
and $f(u)=0$ otherwise. The graph is a horizontal rectangle of height $0.5$ between $u=1$ and $u=3$ (total area $=1$).

\begin{popfigure}
\begin{tikzpicture}[scale=1.2]
\draw[->] (0,0) -- (4,0) node[right] {$u$ (hours)};
\draw[->] (0,0) -- (0,0.7) node[above] {$f(u)$};
\draw[thick, popblue] (1,0) -- (1,0.5) -- (3,0.5) -- (3,0);
\draw[dashed, popmuted] (1,0) node[below] {$1$} -- (1,0.5);
\draw[dashed, popmuted] (3,0) node[below] {$3$} -- (3,0.5);
\draw[dashed, popmuted] (0,0.5) node[left] {$0.5$} -- (1,0.5);
\draw[poporange, thick, <->] (1,0.55) -- (3,0.55) node[midway, above] {$b-a=2$};
\draw[poporange, thick, <->] (3.2,0) -- (3.2,0.5) node[midway, right] {height $=1/2$};
\end{tikzpicture}
\legende{Probability density function of $U\sim \mathrm{Uniform}[1,3]$}
\end{popfigure}

Mean and variance:
\[
E(U)=\frac{a+b}{2}=\frac{1+3}{2}=2\ \text{hours},\qquad
\mathrm{Var}(U)=\frac{(b-a)^2}{12}=\frac{(3-1)^2}{12}=\frac{4}{12}=\frac{1}{3}\ \text{hours}^2.
\]
The standard deviation is $\sqrt{1/3}\approx 0.577$ hour, about 35 minutes.

Probability of arrival within 1.5 hours:
\[
P(U<1.5)=\frac{1.5-1}{3-1}=\frac{0.5}{2}=0.25.
\]

\textbf{Task 3.} Here $X\sim N(850,\,60^2)$. Standardising with $Z=\dfrac{X-850}{60}$:
\[
P(X>950)=P\!\left(Z>\frac{950-850}{60}\right)=P(Z>1.6667).
\]
From the standard normal table, $\Phi(1.67)\approx 0.9525$, so
\[
P(X>950)=1-0.9525=0.0475.
\]
So consumption exceeds $950\ \mathrm{m^3}$ on roughly $4.75\%$ of days, about one day in twenty.

For the level $c$ exceeded on only $5\%$ of days:
\[
P(X>c)=0.05 \iff P\!\left(Z>\frac{c-850}{60}\right)=0.05
\iff \Phi\!\left(\frac{c-850}{60}\right)=0.95.
\]
The table gives $\Phi(1.645)\approx 0.95$, hence
\[
\frac{c-850}{60}=1.645 \quad\Longrightarrow\quad c=850+1.645\times 60=948.7\approx 949\ \mathrm{m^3}.
\]
\emph{Interpretation:} the pump and pipes must be able to deliver at least about $949\ \mathrm{m^3}$ per day if the station is to meet demand on $95\%$ of days.

\textbf{Task 4.} Exactly, $Y\sim B(200,\,0.6)$, since there is a fixed number of independent households, each with the same probability $p=0.6$ of paying on time.

Validity conditions:
\[
np=200\times 0.6=120>5,\qquad nq=200\times 0.4=80>5,
\]
so the normal approximation is appropriate. The approximating distribution is
\[
Y\approx N(np,\,npq)=N(120,\,48),\qquad \sigma=\sqrt{48}\approx 6.928.
\]
``At least 110'' means $Y\geq 110$, i.e.\ $Y=110,111,\dots,200$. With the continuity correction, the integer $110$ is represented by the interval $[109.5,\,110.5]$, so
\[
P(Y\geq 110)\approx P(Y_N\geq 109.5)
=P\!\left(Z\geq \frac{109.5-120}{6.928}\right)
=P(Z\geq -1.516).
\]
By symmetry and the table,
\[
P(Z\geq -1.516)=\Phi(1.516)\approx 0.935.
\]
So there is approximately a $93.5\%$ chance that at least 110 households pay on time.

The continuity correction is essential because we approximate a \emph{discrete} variable (whole households) by a \emph{continuous} one: without it we would compute $P(Z\geq -1.66)$ using the boundary $110$ itself, which treats the integer $110$ as a point of zero width and systematically distorts the probability.

\textbf{Task 5 (guidance).} A good presentation states, for each model: \emph{exponential} because breakdowns are random, independent events at a constant rate; \emph{uniform} because every arrival time in $[1,3]$ is equally likely; \emph{normal} because consumption results from many small independent influences; \emph{normal approximation} because $n$ is large with $np$ and $nq$ both greater than $5$, with continuity correction because the binomial is discrete.

\begin{attention}[Pitfalls to avoid in this activity]
\begin{itemize}
\item Do not confuse $\lambda$ with the mean: for the exponential distribution the mean is $\dfrac{1}{\lambda}$, here $40$, not $\lambda=40$.
\item $P(T>60)=e^{-\lambda \times 60}$: remember that $\lambda$ must be in \emph{matching units} (per day, since $60$ is in days).
\item For the uniform distribution, $\mathrm{Var}(U)=\dfrac{(b-a)^2}{12}$, not $\dfrac{(b-a)^2}{2}$.
\item In Task 3(b), the $5\%$ is in the \emph{upper} tail, so use $\Phi(z)=0.95$, i.e.\ $z=+1.645$, not $-1.645$.
\item Never apply the normal approximation without first checking $np>5$ and $nq>5$; and never forget the continuity correction when moving from a discrete to a continuous model.
\end{itemize}
\end{attention}

\begin{savaistu}[Why this matters in Cameroon]
Water and electricity utilities such as ENEO and Camwater plan maintenance teams, reservoir sizes and billing forecasts using exactly these tools: exponential models for time between failures, uniform models when only bounds are known, and normal models for aggregated consumption. Mastering them opens the way to careers in engineering, actuarial science and data analysis — sectors that are growing rapidly in Douala and Yaoundé.
\end{savaistu}

\newpage

\coursec{Methods and worked examples}

\coursec{Special Continuous Distributions}

\courssub{Introduction and Overview}

In the previous chapter we studied the general theory of continuous random variables: probability density functions, cumulative distribution functions, expectation and variance. We now turn to three \emph{specific} distributions that appear constantly in modelling real phenomena — the exponential, the uniform (rectangular) and the normal distributions. The first two have simple algebraic forms that allow exact calculations; the normal distribution, while not integrable in elementary terms, is the cornerstone of statistical inference because of the Central Limit Theorem. We shall also learn how the normal distribution can approximate the binomial and Poisson distributions when the parameters are large, a technique that remains essential for GCE examination questions.

\begin{savaistu}[Why these three?]
The exponential distribution models \emph{waiting times} between random events (phone calls arriving at a switchboard, radioactive decays, component lifetimes). The uniform distribution models \emph{complete ignorance} over a finite interval (a random arrival time, a rounding error). The normal distribution arises whenever many small independent effects add up (heights, measurement errors, examination marks). Together they cover a vast range of applications in engineering, economics, biology and the social sciences.
\end{savaistu}

\courssub{The Exponential Distribution}

\begin{definition}[Exponential distribution]
A continuous random variable $X$ follows an \emph{exponential distribution} with rate parameter $\lambda > 0$, written $X \sim \mathrm{Exp}(\lambda)$, if its probability density function is
\[
f(x) = \begin{cases}
\lambda e^{-\lambda x}, & x \geq 0, \\[4pt]
0, & x < 0.
\end{cases}
\]
The cumulative distribution function is
\[
F(x) = P(X \leq x) = \begin{cases}
1 - e^{-\lambda x}, & x \geq 0, \\[4pt]
0, & x < 0.
\end{cases}
\]
\end{definition}

\begin{propriete}[Mean and variance of the exponential distribution]
For $X \sim \mathrm{Exp}(\lambda)$:
\[
E(X) = \frac{1}{\lambda}, \qquad \mathrm{Var}(X) = \frac{1}{\lambda^{2}}, \qquad \sigma = \frac{1}{\lambda}.
\]
\textbf{Proof (examinable).} 
\begin{align*}
E(X) &= \int_{0}^{\infty} x \lambda e^{-\lambda x}\,dx 
= \lambda \left[ -\frac{x}{\lambda}e^{-\lambda x} \right]_{0}^{\infty} + \int_{0}^{\infty} \frac{1}{\lambda}e^{-\lambda x}\,dx 
= 0 + \frac{1}{\lambda} = \frac{1}{\lambda}. \\[6pt]
E(X^{2}) &= \int_{0}^{\infty} x^{2} \lambda e^{-\lambda x}\,dx 
= \frac{2}{\lambda^{2}} \quad\text{(integration by parts twice)}. \\[6pt]
\mathrm{Var}(X) &= E(X^{2}) - [E(X)]^{2} = \frac{2}{\lambda^{2}} - \frac{1}{\lambda^{2}} = \frac{1}{\lambda^{2}}.
\end{align*}
\end{propriete}

\begin{propriete}[Memoryless property]
The exponential distribution is the \emph{only} continuous distribution with the memoryless property: for all $s, t \geq 0$,
\[
P(X > s + t \mid X > s) = P(X > t) = e^{-\lambda t}.
\]
This means that a component that has already survived $s$ hours is \emph{as good as new} — its remaining lifetime has exactly the same exponential distribution.
\end{propriete}

\begin{methode}[Handling an exponential model]
\begin{enumerate}
\item \textbf{Identify the parameter $\lambda$.} If the mean $\mu$ is given, $\lambda = 1/\mu$. If $\lambda$ is given directly, $\mu = 1/\lambda$ and $\sigma^{2} = 1/\lambda^{2}$.
\item \textbf{Use the survival function} $P(X > x) = e^{-\lambda x}$ for $x \geq 0$. This is faster than integrating the pdf.
\item \textbf{For an interval} $a < X < b$ ($0 \leq a < b$): $P(a < X < b) = e^{-\lambda a} - e^{-\lambda b}$.
\item \textbf{Apply the memoryless property} for conditional probabilities of the form $P(X > s+t \mid X > s)$.
\item \textbf{Check units}: $x$ and $\lambda$ must use the same time unit (both in minutes, or both in hours, etc.).
\end{enumerate}
\end{methode}

\begin{exemplebox}[Lifetime of an electronic component]
The lifetime $T$ (in hours) of an electronic component used in a phone repair shop in Douala is modelled by an exponential distribution with mean $500$ hours.
\begin{enumerate}
\item[(a)] Find the probability that a component lasts more than $600$ hours.
\item[(b)] Find the probability that a component lasts between $200$ and $400$ hours.
\item[(c)] Given that a component has already worked for $300$ hours, find the probability that it will work for at least $200$ more hours.
\item[(d)] Find the median lifetime.
\end{enumerate}
\end{exemplebox}

\begin{exoresolu}[Lifetime of an electronic component]
\textbf{Solution.} Mean $= 500$ hours $\Rightarrow \lambda = \dfrac{1}{500} = 0.002$ per hour.

\textbf{(a)} $P(T > 600) = e^{-\lambda \times 600} = e^{-0.002 \times 600} = e^{-1.2} \approx 0.3012$.

\textbf{(b)} $P(200 < T < 400) = e^{-0.002 \times 200} - e^{-0.002 \times 400} = e^{-0.4} - e^{-0.8}$.
Numerically, $e^{-0.4} \approx 0.6703$, $e^{-0.8} \approx 0.4493$, so $P \approx 0.6703 - 0.4493 = 0.2210$.

\textbf{(c)} By the memoryless property:
$P(T > 300 + 200 \mid T > 300) = P(T > 200) = e^{-0.002 \times 200} = e^{-0.4} \approx 0.6703$.

\textbf{(d)} Median $m$ satisfies $P(T \leq m) = \frac{1}{2}$:
$1 - e^{-0.002 m} = \frac{1}{2} \;\Longrightarrow\; e^{-0.002m} = \frac{1}{2} \;\Longrightarrow\; 0.002\,m = \ln 2$.
Hence $m = \dfrac{\ln 2}{0.002} = 500 \ln 2 \approx 346.6$ hours.
Note: median ($346.6$ h) $<$ mean ($500$ h) — typical for a right-skewed distribution.
\end{exoresolu}

\begin{attention}[Common mistakes with the exponential distribution]
\begin{itemize}
\item Confusing $\lambda$ (rate) with $1/\lambda$ (mean). Always write $\lambda = 1/\text{mean}$ explicitly.
\item Forgetting that the pdf is zero for $x < 0$; probabilities for negative values are zero.
\item Using the memoryless property incorrectly: it applies \emph{only} to $P(X > s+t \mid X > s)$, not to $P(X < s+t \mid X < s)$.
\item Mixing units (e.g., $\lambda$ in per hour but $x$ in minutes). Convert everything to the same unit first.
\end{itemize}
\end{attention}

\courssub{The Uniform (Rectangular) Distribution}

\begin{definition}[{Uniform distribution on $[a,b]$}]
A continuous random variable $X$ follows a \emph{uniform} (or \emph{rectangular}) distribution on the interval $[a,b]$ ($a < b$), written $X \sim U(a,b)$, if its probability density function is
\[
f(x) = \begin{cases}
\dfrac{1}{b-a}, & a \leq x \leq b, \\[6pt]
0, & \text{otherwise}.
\end{cases}
\]
The cumulative distribution function is
\[
F(x) = P(X \leq x) = \begin{cases}
0, & x < a, \\[4pt]
\dfrac{x-a}{b-a}, & a \leq x \leq b, \\[6pt]
1, & x > b.
\end{cases}
\]
\end{definition}

\begin{propriete}[Mean and variance of the uniform distribution]
For $X \sim U(a,b)$:
\[
E(X) = \frac{a+b}{2}, \qquad \mathrm{Var}(X) = \frac{(b-a)^{2}}{12}, \qquad \sigma = \frac{b-a}{\sqrt{12}}.
\]
\textbf{Proof (examinable).}
\begin{align*}
E(X) &= \int_{a}^{b} x \cdot \frac{1}{b-a}\,dx = \frac{1}{b-a} \left[ \frac{x^{2}}{2} \right]_{a}^{b} = \frac{b^{2}-a^{2}}{2(b-a)} = \frac{a+b}{2}. \\[6pt]
E(X^{2}) &= \int_{a}^{b} x^{2} \cdot \frac{1}{b-a}\,dx = \frac{1}{b-a} \left[ \frac{x^{3}}{3} \right]_{a}^{b} = \frac{b^{3}-a^{3}}{3(b-a)} = \frac{a^{2}+ab+b^{2}}{3}. \\[6pt]
\mathrm{Var}(X) &= E(X^{2}) - [E(X)]^{2} = \frac{a^{2}+ab+b^{2}}{3} - \frac{(a+b)^{2}}{4} = \frac{(b-a)^{2}}{12}.
\end{align*}
\end{propriete}

\begin{methode}[Handling a uniform (rectangular) model]
\begin{enumerate}
\item \textbf{Write down the key results}:
$E(X) = \dfrac{a+b}{2}$, $\quad \mathrm{Var}(X) = \dfrac{(b-a)^{2}}{12}$, $\quad \sigma = \dfrac{b-a}{\sqrt{12}}$.
\item \textbf{Probabilities are ratios of lengths} (areas of rectangles):
$P(c < X < d) = \dfrac{d-c}{b-a}$ for $a \leq c < d \leq b$.
\item \textbf{If an endpoint is unknown}, set up an equation from the given mean or variance and solve for $a$ or $b$.
\item \textbf{Always sketch} the rectangle of height $\frac{1}{b-a}$: the required probability is the shaded area.
\end{enumerate}
\end{methode}

\begin{exemplebox}[Waiting time at a bus stop]
Buses leave a motor park in Bamenda every $20$ minutes. A passenger arrives at random. Her waiting time $X$ (in minutes) is uniformly distributed on $[0, 20]$.
\begin{enumerate}
\item[(a)] Write down the probability density function of $X$.
\item[(b)] Calculate the mean and the standard deviation of the waiting time.
\item[(c)] Find the probability that she waits more than $12$ minutes.
\item[(d)] Find the probability that she waits between $5$ and $15$ minutes.
\end{enumerate}
\end{exemplebox}

\begin{exoresolu}[Waiting time at a bus stop]
\textbf{Solution.} Here $a = 0$, $b = 20$, so $b-a = 20$.

\textbf{(a)} $f(x) = \begin{cases} \frac{1}{20} & 0 \leq x \leq 20, \\ 0 & \text{otherwise.} \end{cases}$

\textbf{(b)} $E(X) = \frac{0+20}{2} = 10$ minutes.
$\mathrm{Var}(X) = \frac{(20-0)^{2}}{12} = \frac{400}{12} = \frac{100}{3}$.
$\sigma = \sqrt{\frac{100}{3}} = \frac{10}{\sqrt{3}} \approx 5.77$ minutes.

\textbf{(c)} $P(X > 12) = \frac{20-12}{20} = \frac{8}{20} = 0.4$.

\textbf{(d)} $P(5 < X < 15) = \frac{15-5}{20} = \frac{10}{20} = 0.5$.
Each answer is the area of a rectangle of height $\frac{1}{20}$ over the relevant interval.
\end{exoresolu}

\begin{exemplebox}[Finding an unknown endpoint]
A continuous random variable $X$ is uniformly distributed on $[2, b]$ with $b > 2$. Given that $\mathrm{Var}(X) = 12$, find
\begin{enumerate}
\item[(a)] the value of $b$;
\item[(b)] $E(X)$ and $P(X > 5)$.
\end{enumerate}
\end{exemplebox}

\begin{exoresolu}[Finding an unknown endpoint]
\textbf{Solution. (a)} For $X \sim U(2, b)$:
$\mathrm{Var}(X) = \frac{(b-2)^{2}}{12} = 12 \;\Longrightarrow\; (b-2)^{2} = 144 \;\Longrightarrow\; b-2 = 12$ (positive root since $b>2$).
Hence $\boxed{b = 14}$.

\textbf{(b)} $E(X) = \frac{2+14}{2} = 8$.
Density height $= \frac{1}{14-2} = \frac{1}{12}$.
$P(X > 5) = \frac{14-5}{12} = \frac{9}{12} = 0.75$.
\end{exoresolu}

\begin{attention}[Common mistakes with the uniform distribution]
\begin{itemize}
\item Using $\mathrm{Var}(X) = \frac{(b-a)^{2}}{2}$ instead of $\frac{(b-a)^{2}}{12}$.
\item Forgetting that the pdf is constant \emph{only} on $[a,b]$; outside this interval it is zero.
\item In probability calculations, using the wrong interval length (e.g., using $b-a$ instead of $d-c$).
\item Not checking that the computed $a$ and $b$ satisfy $a < b$.
\end{itemize}
\end{attention}

\courssub{The Normal Distribution}

\begin{definition}[Normal distribution]
A continuous random variable $X$ follows a \emph{normal distribution} with mean $\mu$ and variance $\sigma^{2}$ ($\sigma > 0$), written $X \sim N(\mu, \sigma^{2})$, if its probability density function is
\[
f(x) = \frac{1}{\sigma\sqrt{2\pi}}\, e^{-\frac{1}{2}\left(\frac{x-\mu}{\sigma}\right)^{2}}, \qquad -\infty < x < \infty.
\]
The curve is symmetric about $x = \mu$, bell-shaped, with inflection points at $x = \mu \pm \sigma$.
\end{definition}

\begin{propriete}[Mean and variance]
For $X \sim N(\mu, \sigma^{2})$: $E(X) = \mu$, $\mathrm{Var}(X) = \sigma^{2}$.
\textbf{Proof (not examinable in full, but the symmetry argument for the mean is expected).}
By symmetry of $f(x)$ about $\mu$, $E(X) = \mu$. The variance derivation uses the substitution $z = \frac{x-\mu}{\sigma}$ and the standard Gaussian integral $\int_{-\infty}^{\infty} z^{2} e^{-z^{2}/2}\,dz = \sqrt{2\pi}$.
\end{propriete}

\begin{aretenir}[Key features of the normal curve]
\begin{itemize}
\item Total area under the curve $= 1$.
\item Symmetric about $\mu$: $P(X < \mu) = P(X > \mu) = 0.5$.
\item Approximately $68\%$ of the area lies within $\mu \pm \sigma$, $95\%$ within $\mu \pm 2\sigma$, $99.7\%$ within $\mu \pm 3\sigma$ (empirical rule).
\item The standard normal distribution is $Z \sim N(0,1)$ with pdf $\phi(z) = \frac{1}{\sqrt{2\pi}} e^{-z^{2}/2}$ and cdf $\Phi(z) = P(Z \leq z)$.
\end{itemize}
\end{aretenir}

\begin{popfigure}
\begin{tikzpicture}[scale=1.2, domain=-3.5:3.5]
\begin{axis}[
    axis lines=middle,
    xlabel=$z$, ylabel=$\phi(z)$,
    xmin=-3.5, xmax=3.5, ymin=0, ymax=0.45,
    xtick={-3,-2,-1,0,1,2,3},
    ytick={0.1,0.2,0.3,0.4},
    width=12cm, height=6cm,
    samples=200,
    clip=false
]
\addplot[popblue, thick] {exp(-x^2/2)/sqrt(max(0,2*pi))};
\addplot[popblue!30, fill=popblue!20, domain=-1:1] {exp(-x^2/2)/sqrt(max(0,2*pi))} \closedcycle;
\addplot[popblue!30, fill=popblue!20, domain=-2:2] {exp(-x^2/2)/sqrt(max(0,2*pi))} \closedcycle;
\addplot[popblue!30, fill=popblue!20, domain=-3:3] {exp(-x^2/2)/sqrt(max(0,2*pi))} \closedcycle;
\node[popblue, anchor=south] at (axis cs:0,0.4) {$\approx 99.7\%$};
\node[popblue, anchor=south] at (axis cs:0,0.3) {$\approx 95\%$};
\node[popblue, anchor=south] at (axis cs:0,0.2) {$\approx 68\%$};
\draw[dashed, popmuted] (axis cs:-1,0) -- (axis cs:-1,0.242);
\draw[dashed, popmuted] (axis cs:1,0) -- (axis cs:1,0.242);
\draw[dashed, popmuted] (axis cs:-2,0) -- (axis cs:-2,0.054);
\draw[dashed, popmuted] (axis cs:2,0) -- (axis cs:2,0.054);
\draw[dashed, popmuted] (axis cs:-3,0) -- (axis cs:-3,0.0044);
\draw[dashed, popmuted] (axis cs:3,0) -- (axis cs:3,0.0044);
\node[below] at (axis cs:-1,0) {$-1$};
\node[below] at (axis cs:1,0) {$1$};
\node[below] at (axis cs:-2,0) {$-2$};
\node[below] at (axis cs:2,0) {$2$};
\node[below] at (axis cs:-3,0) {$-3$};
\node[below] at (axis cs:3,0) {$3$};
\end{axis}
\end{tikzpicture}
\legende{Standard normal density $\phi(z)$ with empirical rule areas}
\end{popfigure}

\courssub{The Standard Normal Variable and Use of Tables}

\begin{propriete}[Standardisation]
If $X \sim N(\mu, \sigma^{2})$, then $Z = \dfrac{X - \mu}{\sigma} \sim N(0,1)$.
For any $x$,
\[
P(X < x) = P\!\left(Z < \frac{x-\mu}{\sigma}\right) = \Phi\!\left(\frac{x-\mu}{\sigma}\right),
\]
where $\Phi(z)$ is the cumulative distribution function of the standard normal variable, tabulated for $z \geq 0$.
\end{propriete}

\begin{methode}[Computing $P(X < x)$ for $X \sim N(\mu, \sigma^{2})$]
\begin{enumerate}
\item \textbf{Standardise}: compute $z = \dfrac{x - \mu}{\sigma}$.
\item \textbf{Rewrite the probability} in terms of $Z$: $P(X < x) = P(Z < z) = \Phi(z)$.
\item \textbf{Read the table} for $z \geq 0$. For negative $z$ use symmetry:
\[
\Phi(-z) = 1 - \Phi(z), \quad P(Z > z) = 1 - \Phi(z), \quad P(a < Z < b) = \Phi(b) - \Phi(a).
\]
\item \textbf{Round} $z$ to two decimal places before using the table; give final probability to four decimal places.
\item \textbf{Sanity check}: if $x < \mu$ then $z < 0$ and the probability must be $< 0.5$; if $x > \mu$ then $z > 0$ and probability $> 0.5$.
\end{enumerate}
\end{methode}

\begin{exemplebox}[Masses of bags of rice]
The mass $X$ (in kg) of bags of rice produced by a mill is normally distributed with mean $50.2$ kg and standard deviation $0.4$ kg. A bag is chosen at random. Find the probability that its mass is
\begin{enumerate}
\item[(a)] less than $50.8$ kg;
\item[(b)] more than $49.5$ kg;
\item[(c)] between $49.6$ kg and $50.6$ kg.
\end{enumerate}
(Use $\Phi(1.50) = 0.9332$, $\Phi(1.75) = 0.9599$, $\Phi(1.00) = 0.8413$.)
\end{exemplebox}

\begin{exoresolu}[Masses of bags of rice]
\textbf{Solution.} $\mu = 50.2$, $\sigma = 0.4$.

\textbf{(a)} $z = \frac{50.8 - 50.2}{0.4} = 1.50$.
$P(X < 50.8) = P(Z < 1.50) = \Phi(1.50) = 0.9332$.

\textbf{(b)} $z = \frac{49.5 - 50.2}{0.4} = -1.75$.
$P(X > 49.5) = P(Z > -1.75) = \Phi(1.75) = 0.9599$.

\textbf{(c)} $z_1 = \frac{49.6 - 50.2}{0.4} = -1.50$, $z_2 = \frac{50.6 - 50.2}{0.4} = 1.00$.
\begin{align*}
P(49.6 < X < 50.6) &= \Phi(1.00) - \Phi(-1.50) \\
&= \Phi(1.00) - [1 - \Phi(1.50)] \\
&= 0.8413 - (1 - 0.9332) = 0.8413 - 0.0668 = 0.7745.
\end{align*}
About $77.45\%$ of bags have mass between $49.6$ kg and $50.6$ kg.
\end{exoresolu}

\begin{exemplebox}[More practice with normal probabilities]
The heights of adult males in a certain population are normally distributed with mean $172$ cm and standard deviation $7$ cm. Find:
\begin{enumerate}
\item[(a)] $P(\text{height} > 185 \text{ cm})$;
\item[(b)] $P(165 \text{ cm} < \text{height} < 180 \text{ cm})$;
\item[(c)] The height exceeded by the tallest $10\%$ of males (use $\Phi(1.28) = 0.8997$).
\end{enumerate}
\end{exemplebox}

\begin{exoresolu}[More practice with normal probabilities]
\textbf{Solution.} $X \sim N(172, 7^{2})$.

\textbf{(a)} $z = \frac{185-172}{7} = \frac{13}{7} \approx 1.86$.
$P(X > 185) = P(Z > 1.86) = 1 - \Phi(1.86) \approx 1 - 0.9686 = 0.0314$.

\textbf{(b)} $z_1 = \frac{165-172}{7} = -1.00$, $z_2 = \frac{180-172}{7} \approx 1.14$.
$P(165 < X < 180) = \Phi(1.14) - \Phi(-1.00) = \Phi(1.14) - [1 - \Phi(1.00)]$.
Using tables: $\Phi(1.14) \approx 0.8729$, $\Phi(1.00) = 0.8413$.
$P \approx 0.8729 - 0.1587 = 0.7142$.

\textbf{(c)} Let $h$ be the height. $P(X > h) = 0.10 \Rightarrow P(X < h) = 0.90$.
$\frac{h-172}{7} = 1.28 \Rightarrow h = 172 + 1.28 \times 7 = 172 + 8.96 = 180.96 \approx 181$ cm.
\end{exoresolu}

\courssub{Inverse Normal Problems — Finding $x$, $\mu$ or $\sigma$}

\begin{methode}[Working backwards from a given probability]
\begin{enumerate}
\item \textbf{Find $z$ from the table}: locate the given probability in the body of the table and read off $z$. If the probability is $< 0.5$ and the event is $P(X < x)$, the $z$-value is negative.
\item \textbf{Translate into an equation} using $z = \dfrac{x - \mu}{\sigma}$.
\item \textbf{Solve} for the unknown: $x = \mu + z\sigma$, or $\mu = x - z\sigma$, or $\sigma = \dfrac{x - \mu}{z}$.
\item \textbf{Two probabilities} give a pair of simultaneous equations in $\mu$ and $\sigma$: solve by elimination.
\end{enumerate}
\textbf{Useful table values:} $\Phi(1.645) = 0.95$, $\Phi(1.96) = 0.975$, $\Phi(2.326) = 0.99$, $\Phi(0.84) \approx 0.80$, $\Phi(1.28) \approx 0.90$.
\end{methode}

\begin{exemplebox}[Finding a cut-off mark]
Scores in a Further Mathematics mock examination in a school in Buea are modelled by $X \sim N(58, 12^{2})$. The top $10\%$ of candidates receive a distinction. Find the minimum score needed for a distinction, giving your answer to the nearest integer.
(Use $\Phi(1.28) = 0.8997$.)
\end{exemplebox}

\begin{exoresolu}[Finding a cut-off mark]
\textbf{Solution.} Let $d$ be the cut-off score. $P(X > d) = 0.10 \Rightarrow P(X < d) = 0.90$.
Standardising: $P\!\left(Z < \frac{d-58}{12}\right) = 0.90$.
From tables, $\Phi(1.28) = 0.8997 \approx 0.90$, so $\frac{d-58}{12} = 1.28$.
$d = 58 + 1.28 \times 12 = 58 + 15.36 = 73.36$.
Minimum score for distinction $\approx \boxed{74}$ marks (rounding up to guarantee top $10\%$).
\end{exoresolu}

\begin{exemplebox}[Finding $\mu$ and $\sigma$ simultaneously]
The masses of pineapples sold at a market in Yaoundé are normally distributed. It is observed that $15\%$ of the pineapples have mass less than $0.8$ kg and $8\%$ have mass greater than $1.6$ kg. Find the mean and standard deviation of the masses.
(Use $\Phi(1.04) = 0.8508$, $\Phi(1.41) = 0.9207$.)
\end{exemplebox}

\begin{exoresolu}[Finding $\mu$ and $\sigma$ simultaneously]
\textbf{Solution.} Let $X \sim N(\mu, \sigma^{2})$ be the mass in kg.

\textbf{First condition.} $P(X < 0.8) = 0.15$. Since $0.15 < 0.5$, $z$ is negative:
$P\!\left(Z < \frac{0.8-\mu}{\sigma}\right) = 0.15 \Rightarrow \frac{0.8-\mu}{\sigma} = -1.04$,
because $\Phi(1.04) = 0.8508 \approx 0.85$ and $\Phi(-1.04) = 1 - 0.85 = 0.15$.
Thus $\mu - 1.04\sigma = 0.8$. \quad (1)

\textbf{Second condition.} $P(X > 1.6) = 0.08 \Rightarrow P(X < 1.6) = 0.92$:
$\frac{1.6-\mu}{\sigma} = 1.41 \Rightarrow \mu + 1.41\sigma = 1.6$. \quad (2)

\textbf{Solving.} Subtract (1) from (2):
$(1.41 + 1.04)\sigma = 1.6 - 0.8 \Rightarrow 2.45\sigma = 0.8 \Rightarrow \sigma = \frac{0.8}{2.45} \approx 0.327$ kg.
Substitute into (2): $\mu = 1.6 - 1.41 \times 0.327 \approx 1.6 - 0.461 = 1.139$ kg.
Mean $\approx \boxed{1.14 \text{ kg}}$, standard deviation $\approx \boxed{0.33 \text{ kg}}$.
\end{exoresolu}

\begin{attention}[Inverse normal problems — common errors]
\begin{itemize}
\item Forgetting the negative sign on $z$ when the given probability is $< 0.5$ for a ``less than'' event.
\item Using $\sigma^{2}$ instead of $\sigma$ in the standardisation formula.
\item Rounding $z$ too early; keep at least 3 decimal places during calculation.
\item In simultaneous equations, mixing up which equation corresponds to which tail.
\end{itemize}
\end{attention}

\courssub{Normal Approximation to the Binomial Distribution}

\begin{propriete}[Normal approximation to the binomial]
If $X \sim B(n, p)$ with $n$ large and $np > 5$, $n(1-p) > 5$ (some texts use $npq > 5$), then
\[
X \approx Y \sim N\bigl(np,\; np(1-p)\bigr).
\]
\end{propriete}

\begin{methode}[Approximating $B(n,p)$ by a normal distribution]
\begin{enumerate}
\item \textbf{Check validity}: $n$ large, $np > 5$ and $n(1-p) > 5$.
\item \textbf{Match parameters}: $\mu = np$, $\sigma^{2} = np(1-p)$.
\item \textbf{Apply continuity correction}: the binomial is discrete, the normal continuous. Each integer $k$ corresponds to the interval $[k-0.5,\; k+0.5]$:
\begin{align*}
P(X = k) &\approx P(k-0.5 < Y < k+0.5), \\
P(X \leq k) &\approx P(Y < k+0.5), \\
P(X \geq k) &\approx P(Y > k-0.5), \\
P(X < k) = P(X \leq k-1) &\approx P(Y < k-0.5), \\
P(X > k) &\approx P(Y > k+0.5).
\end{align*}
\item \textbf{Standardise} $Y$ and read the standard normal table.
\end{enumerate}
\end{methode}

\begin{attention}[The continuity correction is compulsory]
Forgetting the $\pm 0.5$ adjustment is the most common error at the GCE. Always ask: does the event \emph{include} the boundary integer?
\begin{itemize}
\item If $X \leq k$, integer $k$ is included $\rightarrow$ extend to $k+0.5$.
\item If $X < k$, integer $k$ is excluded $\rightarrow$ stop at $k-0.5$.
\item If $X \geq k$, integer $k$ is included $\rightarrow$ extend to $k-0.5$.
\item If $X > k$, integer $k$ is excluded $\rightarrow$ start at $k+0.5$.
\end{itemize}
\end{attention}

\begin{popfigure}
\begin{tikzpicture}[scale=0.9]
\begin{axis}[
    axis lines=middle,
    xlabel=$x$, ylabel=Probability,
    xmin=0, xmax=15, ymin=0, ymax=0.25,
    xtick={0,1,...,15},
    ytick={0.05,0.1,0.15,0.2},
    width=13cm, height=6cm,
    samples=200
]
% Binomial bars (n=10, p=0.5)
\addplot[ybar, bar width=0.4, fill=popblue!60, draw=popblue] coordinates {
    (0,0.001) (1,0.01) (2,0.044) (3,0.117) (4,0.205) (5,0.246)
    (6,0.205) (7,0.117) (8,0.044) (9,0.01) (10,0.001)
};
% Normal curve
\addplot[poporange, thick, domain=0:10] {1/sqrt(max(0,2*pi*2.5))*exp(-(x-5)^2/(2*2.5))};
% Continuity correction shading for P(X=5)
\addplot[poporange!30, fill=poporange!20, domain=4.5:5.5] {1/sqrt(max(0,2*pi*2.5))*exp(-(x-5)^2/(2*2.5))} \closedcycle;
\node[poporange, anchor=south] at (axis cs:5,0.22) {$P(4.5 < Y < 5.5)$};
\node[popblue, anchor=north] at (axis cs:5,0.246) {$P(X=5)$};
\draw[dashed, popmuted] (axis cs:4.5,0) -- (axis cs:4.5,0.15);
\draw[dashed, popmuted] (axis cs:5.5,0) -- (axis cs:5.5,0.15);
\node[below, font=\small] at (axis cs:4.5,0) {$4.5$};
\node[below, font=\small] at (axis cs:5.5,0) {$5.5$};
\end{axis}
\end{tikzpicture}
\legende{Continuity correction: $P(X=5)$ for $B(10,0.5)$ approximated by $P(4.5 < Y < 5.5)$ for $N(5, 2.5)$}
\end{popfigure}

\begin{exemplebox}[Defective items in a factory]
A machine produces items of which $4\%$ are defective. A batch of $200$ items is selected at random. Using a normal approximation, find the probability that the batch contains
\begin{enumerate}
\item[(a)] fewer than $10$ defective items;
\item[(b)] exactly $8$ defective items.
\end{enumerate}
(Use $\Phi(0.54) = 0.7054$, $\Phi(0.18) = 0.5714$, $\Phi(0.90) = 0.8159$.)
\end{exemplebox}

\begin{exoresolu}[Defective items in a factory]
\textbf{Solution.} Let $X \sim B(200, 0.04)$ be the number of defective items.

\textbf{Validity.} $np = 200 \times 0.04 = 8 > 5$, $n(1-p) = 200 \times 0.96 = 192 > 5$: approximation valid.

\textbf{Parameters.}
$\mu = np = 8$, $\sigma^{2} = np(1-p) = 8 \times 0.96 = 7.68$, $\sigma = \sqrt{7.68} \approx 2.771$.
$X \approx Y \sim N(8, 7.68)$.

\textbf{(a)} ``Fewer than $10$'' means $X \leq 9$; integer $9$ included $\rightarrow$ boundary $9.5$:
$P(X < 10) = P(X \leq 9) \approx P(Y < 9.5)$.
$z = \frac{9.5 - 8}{2.771} = \frac{1.5}{2.771} \approx 0.54$.
$P(X < 10) \approx \Phi(0.54) = 0.7054$.

\textbf{(b)} Exactly $8$ corresponds to interval $[7.5, 8.5]$:
$P(X = 8) \approx P(7.5 < Y < 8.5)$.
$z_1 = \frac{7.5 - 8}{2.771} \approx -0.18$, $z_2 = \frac{8.5 - 8}{2.771} \approx 0.18$.
$P(X = 8) \approx \Phi(0.18) - \Phi(-0.18) = 2\Phi(0.18) - 1 = 2(0.5714) - 1 = 0.1428$.
\end{exoresolu}

\begin{exemplebox}[Another binomial approximation]
A fair coin is tossed $100$ times. Use a normal approximation to find the probability of obtaining between $45$ and $55$ heads inclusive.
(Use $\Phi(1.00) = 0.8413$, $\Phi(0.00) = 0.5000$.)
\end{exemplebox}

\begin{exoresolu}[Another binomial approximation]
\textbf{Solution.} $X \sim B(100, 0.5)$. $np = 50 > 5$, $n(1-p) = 50 > 5$.
$\mu = 50$, $\sigma^{2} = 25$, $\sigma = 5$.
$X \approx Y \sim N(50, 25)$.
``Between $45$ and $55$ inclusive'' $\Rightarrow 45 \leq X \leq 55$; both integers included $\rightarrow$ interval $[44.5, 55.5]$.
$P(45 \leq X \leq 55) \approx P(44.5 < Y < 55.5)$.
$z_1 = \frac{44.5-50}{5} = -1.10$, $z_2 = \frac{55.5-50}{5} = 1.10$.
$P \approx \Phi(1.10) - \Phi(-1.10) = 2\Phi(1.10) - 1$.
Using $\Phi(1.10) \approx 0.8643$: $P \approx 2(0.8643) - 1 = 0.7286$.
\end{exoresolu}

\courssub{Normal Approximation to the Poisson Distribution}

\begin{propriete}[Normal approximation to the Poisson]
If $X \sim \mathrm{Po}(\lambda)$ with $\lambda$ large (in practice $\lambda > 10$, some texts say $\lambda \geq 15$), then
\[
X \approx Y \sim N(\lambda, \lambda).
\]
\end{propriete}

\begin{methode}[Approximating $\mathrm{Po}(\lambda)$ by a normal distribution]
\begin{enumerate}
\item \textbf{Check validity}: $\lambda > 10$ (or $\lambda \geq 15$).
\item \textbf{Match parameters}: $\mu = \lambda$, $\sigma^{2} = \lambda$, $\sigma = \sqrt{\lambda}$.
\item \textbf{Apply continuity correction} exactly as for the binomial:
$P(X \leq k) \approx P(Y < k+0.5)$, $P(X \geq k) \approx P(Y > k-0.5)$, etc.
\item \textbf{Standardise} with $\mu = \lambda$, $\sigma = \sqrt{\lambda}$, then read the standard normal table.
\end{enumerate}
\end{methode}

\begin{exemplebox}[Calls at a call centre]
A customer-service centre in Douala receives calls at an average rate of $36$ calls per hour. The number of calls received in one hour is modelled by a Poisson distribution. Using a normal approximation, find the probability that in a given hour the centre receives
\begin{enumerate}
\item[(a)] at most $40$ calls;
\item[(b)] more than $30$ calls;
\item[(c)] between $30$ and $40$ calls inclusive.
\end{enumerate}
(Use $\Phi(0.75) = 0.7734$, $\Phi(0.92) = 0.8212$, $\Phi(1.08) = 0.8599$.)
\end{exemplebox}

\begin{exoresolu}[Calls at a call centre]
\textbf{Solution.} $X \sim \mathrm{Po}(36)$. $\lambda = 36 > 10$, so normal approximation appropriate:
$X \approx Y \sim N(36, 36)$, $\mu = 36$, $\sigma = \sqrt{36} = 6$.

\textbf{(a)} ``At most $40$'' means $X \leq 40$; integer $40$ included $\rightarrow$ boundary $40.5$:
$P(X \leq 40) \approx P(Y < 40.5)$.
$z = \frac{40.5 - 36}{6} = \frac{4.5}{6} = 0.75$.
$P(X \leq 40) \approx \Phi(0.75) = 0.7734$.

\textbf{(b)} ``More than $30$'' means $X \geq 31$; integer $31$ included $\rightarrow$ boundary $30.5$:
$P(X > 30) = P(X \geq 31) \approx P(Y > 30.5)$.
$z = \frac{30.5 - 36}{6} = \frac{-5.5}{6} \approx -0.92$.
$P(X > 30) \approx P(Z > -0.92) = \Phi(0.92) = 0.8212$.

\textbf{(c)} ``Between $30$ and $40$ inclusive'' means $30 \leq X \leq 40$; both integers included $\rightarrow$ interval $[29.5, 40.5]$:
$P(30 \leq X \leq 40) \approx P(29.5 < Y < 40.5)$.
$z_1 = \frac{29.5 - 36}{6} = \frac{-6.5}{6} \approx -1.08$, $z_2 = \frac{40.5 - 36}{6} = 0.75$.
\begin{align*}
P(30 \leq X \leq 40) &\approx \Phi(0.75) - \Phi(-1.08) \\
&= \Phi(0.75) - [1 - \Phi(1.08)] \\
&= 0.7734 - (1 - 0.8599) = 0.7734 - 0.1401 = 0.6333.
\end{align*}
\end{exoresolu}

\begin{exemplebox}[Poisson approximation with smaller $\lambda$]
The number of accidents per month at a certain junction follows a Poisson distribution with mean $12$. Use a normal approximation to find the probability that in a given month there are fewer than $8$ accidents.
(Use $\Phi(1.30) = 0.9032$.)
\end{exemplebox}

\begin{exoresolu}[Poisson approximation with smaller $\lambda$]
\textbf{Solution.} $X \sim \mathrm{Po}(12)$. $\lambda = 12 > 10$, approximation valid.
$X \approx Y \sim N(12, 12)$, $\sigma = \sqrt{12} \approx 3.464$.
``Fewer than $8$'' means $X \leq 7$; integer $7$ included $\rightarrow$ boundary $7.5$:
$P(X < 8) = P(X \leq 7) \approx P(Y < 7.5)$.
$z = \frac{7.5 - 12}{\sqrt{12}} = \frac{-4.5}{3.464} \approx -1.30$.
$P(X < 8) \approx \Phi(-1.30) = 1 - \Phi(1.30) = 1 - 0.9032 = 0.0968$.
\end{exoresolu}

\courssub{Choosing the Correct Model and Approximation}

\begin{methode}[Deciding which model and which approximation to use]
\begin{enumerate}
\item \textbf{Fixed number $n$ of independent trials, two outcomes, constant $p$?} $\rightarrow$ Binomial $B(n,p)$.
If $n$ large with $np > 5$ and $n(1-p) > 5$ $\rightarrow$ approximate by $N(np, np(1-p))$.
\item \textbf{Events occurring randomly in time/space at an average rate $\lambda$?} $\rightarrow$ Poisson $\mathrm{Po}(\lambda)$.
If $\lambda > 10$ $\rightarrow$ approximate by $N(\lambda, \lambda)$.
\item \textbf{In both cases}, apply the continuity correction and standardise.
\item \textbf{State the approximation clearly}: examiners award marks for writing ``$X \approx Y \sim N(\ldots)$'' and for justifying the conditions ($np>5$, $n(1-p)>5$, or $\lambda > 10$).
\end{enumerate}
\end{methode}

\begin{exemplebox}[GCE-style two-part problem]
A large secondary school in the North West Region has $1200$ students. Records show that $30\%$ of students eat at the school canteen on any given day, independently of one another.
\begin{enumerate}
\item[(a)] Using a suitable approximation, find the probability that on a given day at least $350$ students eat at the canteen.
\item[(b)] The canteen receives on average $25$ special meal orders per day, modelled by a Poisson distribution. Using a suitable approximation, find the probability that on a given day fewer than $20$ special orders are received.
\end{enumerate}
(Use $\Phi(0.63) = 0.7357$, $\Phi(1.10) = 0.8643$.)
\end{exemplebox}

\begin{exoresolu}[GCE-style two-part problem]
\textbf{Solution. (a)} Let $X$ be the number of students eating at the canteen. $X \sim B(1200, 0.3)$.

\textbf{Validity:} $np = 1200 \times 0.3 = 360 > 5$, $n(1-p) = 840 > 5$ $\rightarrow$ normal approximation applies.
$\mu = 360$, $\sigma^{2} = np(1-p) = 360 \times 0.7 = 252$, $\sigma = \sqrt{252} \approx 15.87$.
$X \approx Y \sim N(360, 252)$.

``At least $350$'' means $X \geq 350$; integer $350$ included $\rightarrow$ boundary $349.5$:
$P(X \geq 350) \approx P(Y > 349.5)$.
$z = \frac{349.5 - 360}{15.87} = \frac{-10.5}{15.87} \approx -0.66$.
$P(X \geq 350) \approx P(Z > -0.66) = \Phi(0.66) \approx 0.7454$ (using $\Phi(0.63)=0.7357$ gives similar method).

\textbf{(b)} Let $W \sim \mathrm{Po}(25)$ be the number of special orders. $\lambda = 25 > 10$:
$W \approx V \sim N(25, 25)$, $\sigma = \sqrt{25} = 5$.

``Fewer than $20$'' means $W \leq 19$; integer $19$ included $\rightarrow$ boundary $19.5$:
$P(W < 20) = P(W \leq 19) \approx P(V < 19.5)$.
$z = \frac{19.5 - 25}{5} = \frac{-5.5}{5} = -1.10$.
$P(W < 20) \approx \Phi(-1.10) = 1 - \Phi(1.10) = 1 - 0.8643 = 0.1357$.

\textbf{Conclusion.} About $74.5\%$ chance that at least $350$ students eat at the canteen; about $13.6\%$ chance of fewer than $20$ special orders.
\end{exoresolu}

\begin{attention}[Final checklist before leaving the examination room]
\begin{itemize}
\item \textbf{Exponential:} Did you use $\lambda = 1/E(X)$ correctly and keep units consistent?
\item \textbf{Uniform:} Did you use $\mathrm{Var}(X) = (b-a)^{2}/12$ and not $(b-a)^{2}/2$?
\item \textbf{Normal:} Did you divide by $\sigma$ (not $\sigma^{2}$) when standardising?
\item \textbf{Inverse problems:} Did you attach a \emph{negative} sign to $z$ when the probability is below $0.5$?
\item \textbf{Approximations:} Did you justify the conditions ($np>5$, $n(1-p)>5$, or $\lambda > 10$) and apply the continuity correction in the right direction?
\item \textbf{Notation:} Did you write ``$X \approx Y \sim N(\ldots)$'' and define $Y$?
\item \textbf{Rounding:} $z$-values to 2 decimal places for table use; final probabilities to 4 decimal places.
\end{itemize}
\end{attention}

\begin{savaistu}[Historical note]
The normal distribution was first derived by Abraham de Moivre in 1733 as an approximation to the binomial distribution — exactly the technique we use in this chapter! Gauss later popularised it in astronomy (hence ``Gaussian distribution''), and it became the foundation of modern statistics through the work of Laplace, Galton, Pearson and Fisher. The exponential distribution was studied by Laplace in the context of error theory, and its memoryless property makes it the continuous analogue of the geometric distribution.
\end{savaistu}

\begin{exercice}[Mixed practice on special continuous distributions]
\begin{enumerate}
\item The time between arrivals of taxis at a rank in Yaoundé follows an exponential distribution with mean $8$ minutes. Find the probability that a passenger waits (a) more than $10$ minutes, (b) between $5$ and $15$ minutes, (c) at least $5$ more minutes given they have already waited $10$ minutes.
\item A random variable $X \sim U(a, b)$ has mean $10$ and variance $48/12 = 4$. Find $a$ and $b$, and hence $P(X > 12)$.
\item The masses of packets of sugar are normally distributed with mean $1.02$ kg and standard deviation $0.015$ kg. Find the probability that a randomly chosen packet has mass (a) less than $1$ kg, (b) between $1.00$ kg and $1.03$ kg. Find the mass exceeded by the heaviest $5\%$ of packets.
\item A factory produces bolts with $2\%$ defectives. In a batch of $500$ bolts, use a normal approximation to find the probability that there are (a) at most $15$ defectives, (b) exactly $10$ defectives.
\item The number of emails received per hour by a company follows a Poisson distribution with mean $18$. Use a normal approximation to find the probability that in a given hour the company receives (a) more than $25$ emails, (b) between $15$ and $20$ emails inclusive.
\end{enumerate}
\end{exercice}

\begin{corrige}[Exercise 1]
\textbf{1.} $\lambda = 1/8 = 0.125$ per minute.
(a) $P(T > 10) = e^{-0.125 \times 10} = e^{-1.25} \approx 0.2865$.
(b) $P(5 < T < 15) = e^{-0.625} - e^{-1.875} \approx 0.5353 - 0.1534 = 0.3819$.
(c) Memoryless property: $P(T > 15 \mid T > 10) = P(T > 5) = e^{-0.625} \approx 0.5353$.

\textbf{2.} $E(X) = (a+b)/2 = 10 \Rightarrow a+b = 20$.
$\mathrm{Var}(X) = (b-a)^2/12 = 4 \Rightarrow (b-a)^2 = 48 \Rightarrow b-a = \sqrt{48} = 4\sqrt{3} \approx 6.928$.
Solving: $b = 10 + 2\sqrt{3} \approx 13.464$, $a = 10 - 2\sqrt{3} \approx 6.536$.
$P(X > 12) = (b-12)/(b-a) = (1.464)/(6.928) \approx 0.2113$.

\textbf{3.} $X \sim N(1.02, 0.015^2)$.
(a) $z = (1 - 1.02)/0.015 = -1.333$. $P(X < 1) = \Phi(-1.33) = 1 - \Phi(1.33) \approx 1 - 0.9082 = 0.0918$.
(b) $z_1 = -1.33$, $z_2 = (1.03-1.02)/0.015 = 0.667$. $P = \Phi(0.67) - \Phi(-1.33) \approx 0.7486 - 0.0918 = 0.6568$.
(c) Top $5\%$: $z = 1.645$. $x = 1.02 + 1.645 \times 0.015 = 1.02 + 0.024675 = 1.0447$ kg.

\textbf{4.} $X \sim B(500, 0.02)$. $np = 10 > 5$, $n(1-p) = 490 > 5$.
$\mu = 10$, $\sigma^2 = 9.8$, $\sigma \approx 3.1305$.
(a) $P(X \leq 15) \approx P(Y < 15.5)$. $z = (15.5-10)/3.1305 = 1.757$. $\Phi(1.76) \approx 0.9608$.
(b) $P(X = 10) \approx P(9.5 < Y < 10.5)$. $z_1 = -0.16$, $z_2 = 0.16$. $P \approx 2\Phi(0.16)-1 \approx 2(0.5636)-1 = 0.1272$.

\textbf{5.} $X \sim \mathrm{Po}(18)$. $\lambda = 18 > 10$. $Y \sim N(18, 18)$, $\sigma = \sqrt{18} \approx 4.243$.
(a) $P(X > 25) = P(X \geq 26) \approx P(Y > 25.5)$. $z = (25.5-18)/4.243 = 1.768$. $P \approx 1 - \Phi(1.77) \approx 1 - 0.9616 = 0.0384$.
(b) $P(15 \leq X \leq 20) \approx P(14.5 < Y < 20.5)$. $z_1 = -0.825$, $z_2 = 0.589$. $P \approx \Phi(0.59) - \Phi(-0.83) \approx 0.7224 - 0.2033 = 0.5191$.
\end{corrige}

\newpage

\coursec{Exercises}

\courssub{I check my knowledge}

\begin{exercice}[Multiple choice: the exponential distribution]
The lifetime $T$ (in hours) of an electronic component follows an exponential distribution with parameter $\lambda = 0.01$. Only one of the following statements is correct. Which one?
\begin{enumerate}
\item $E(T) = 0.01$ hour;
\item $P(T > 100) = e^{-1}$;
\item $P(T > 100) = 1 - e^{-1}$;
\item $\mathrm{Var}(T) = 100$ hours$^2$.
\end{enumerate}
\end{exercice}

\begin{exercice}[True or false?]
Answer \emph{true} or \emph{false} to each of the following statements, justifying your answer carefully.
\begin{enumerate}
\item If $X \sim U(0, 10)$, then $E(X) = 5$ and $\mathrm{Var}(X) = \dfrac{100}{12}$.
\item If $X \sim N(\mu, \sigma^2)$, then $P(X < \mu) = 0.5$.
\item If $Z \sim N(0,1)$, then $P(Z < -1.5) = P(Z < 1.5)$.
\item The exponential distribution has the memoryless property: $P(T > s + t \mid T > s) = P(T > t)$.
\end{enumerate}
\end{exercice}

\begin{exercice}[Definitions and key formulas]
\begin{enumerate}
\item Write down the probability density function of a random variable $X$ following the uniform distribution on $[a, b]$, and sketch its graph.
\item State the probability density function of the standard normal distribution $N(0,1)$ and give two properties of its curve.
\item State the conditions under which a binomial distribution $B(n, p)$ may be approximated by a normal distribution, and give the parameters of the approximating distribution.
\end{enumerate}
\end{exercice}

\begin{exercice}[Direct calculations with the standard normal table]
Let $Z \sim N(0,1)$. Using the standard normal table, find:
\begin{enumerate}
\item $P(Z < 1.25)$;
\item $P(Z > 0.84)$;
\item $P(-1.64 < Z < 2.05)$;
\item the value of $z$ such that $P(Z < z) = 0.95$;
\item the value of $z$ such that $P(|Z| < z) = 0.90$.
\end{enumerate}
\end{exercice}

\courssub{I practise}

\begin{exercice}[Lifetime of a phone battery]
The lifetime $T$ (in hours) of a phone battery follows an exponential distribution with mean $500$ hours.
\begin{enumerate}
\item Write down the probability density function of $T$.
\item Find the probability that the battery lasts more than $800$ hours.
\item Find the probability that the battery fails between $200$ and $400$ hours.
\item Find the median lifetime of the battery.
\end{enumerate}
\end{exercice}

\begin{exercice}[Uniform distribution]
Let $X \sim U(2, 10)$.
\begin{enumerate}
\item Find $E(X)$ and $\mathrm{Var}(X)$.
\item Find $P(X > 5)$.
\item Find the value of $k$ such that $P(X \leq k) = 0.75$.
\item Find $P(3 < X < 7 \mid X > 4)$.
\end{enumerate}
\end{exercice}

\begin{exercice}[Masses of bags of cement]
The masses of bags of cement produced by a factory are normally distributed with mean $50$ kg and variance $0.16$ kg$^2$.
\begin{enumerate}
\item Find the percentage of bags weighing less than $49.5$ kg.
\item Find the percentage of bags weighing between $49.6$ kg and $50.5$ kg.
\item Find the mass exceeded by only $2\%$ of the bags.
\end{enumerate}
\end{exercice}

\begin{exercice}[Waiting time at an Express Union counter]
The waiting time $T$ (in minutes) between the arrivals of two consecutive customers at an Express Union counter follows an exponential distribution with parameter $\lambda = 0.2$.
\begin{enumerate}
\item Find the mean waiting time between two customers.
\item Find the probability of waiting more than $5$ minutes between two customers.
\item Given that no customer has arrived for $3$ minutes, find the probability of waiting at least $2$ more minutes. Comment on your result.
\end{enumerate}
\end{exercice}

\courssub{I prepare for the GCE}

\begin{exercice}[Mock examination marks \hfill (13 marks)]
The marks obtained by candidates in a Further Mathematics mock examination are normally distributed with mean $52$ and standard deviation $\sigma$. It is observed that $15\%$ of the candidates score above $70$.
\begin{enumerate}
\item Show that $\sigma \approx 17.3$, giving your answer to one decimal place. \hfill (4 marks)
\item Find the percentage of candidates who score between $40$ and $60$. \hfill (3 marks)
\item Find the mark $m$ such that $10\%$ of the candidates score below $m$. \hfill (3 marks)
\item Three candidates are chosen at random. Find the probability that exactly two of them score above $70$. \hfill (3 marks)
\end{enumerate}
\end{exercice}

\begin{exercice}[Throwing a die: normal approximation to the binomial \hfill (12 marks)]
A fair die is thrown $180$ times and $X$ denotes the number of sixes obtained.
\begin{enumerate}
\item State the exact distribution of $X$ and calculate $E(X)$ and $\mathrm{Var}(X)$. \hfill (3 marks)
\item Justify the use of a normal approximation for $X$, and state the parameters of the approximating distribution. \hfill (3 marks)
\item Using the normal approximation with a continuity correction, find the probability of obtaining between $25$ and $35$ sixes inclusive. \hfill (4 marks)
\item Using the same approximation, find the probability of obtaining more than $40$ sixes. \hfill (2 marks)
\end{enumerate}
\end{exercice}

\begin{exercice}[Filling machine and telephone calls \hfill (16 marks)]
\textbf{Part A.} A machine fills bottles with a volume that is normally distributed with mean $\mu$ ml and variance $4$ ml$^2$. It is found that $10\%$ of the bottles contain less than $498$ ml.
\begin{enumerate}
\item Show that $\mu \approx 500.6$ ml. \hfill (4 marks)
\item Find the probability that a randomly chosen bottle contains more than $502$ ml. \hfill (2 marks)
\item A bottle is said to be \emph{under-filled} if it contains less than $498$ ml. Using a suitable approximation (which you must justify), find the probability that in a random sample of $100$ bottles, at least $12$ are under-filled. \hfill (5 marks)
\end{enumerate}
\textbf{Part B.} Calls arrive at a police station at an average rate of $30$ per hour, and the number of calls received in one hour follows a Poisson distribution.
\begin{enumerate}
\item[(d)] Justify the use of a normal approximation and, applying a continuity correction, find an approximate value for the probability of receiving more than $40$ calls in one hour. \hfill (5 marks)
\end{enumerate}
\end{exercice}

\newpage

\coursec{Solutions to the exercises}

\coursec{Examination Practice and Worked Examples}

\begin{corrige}[Exercise 1 — Multiple choice: the exponential distribution]
For $T \sim \mathrm{Exp}(\lambda)$ with $\lambda = 0.01$:
\[
E(T) = \frac{1}{\lambda} = 100 \text{ hours}, \qquad \mathrm{Var}(T) = \frac{1}{\lambda^2} = 10\,000 \text{ hours}^2,
\]
and the survival function is $P(T > t) = e^{-\lambda t}$.

\begin{itemize}
\item Statement 1 is \textbf{false}: $E(T) = 100$, not $0.01$.
\item Statement 2: $P(T > 100) = e^{-0.01 \times 100} = e^{-1}$. This is \textbf{correct}.
\item Statement 3 is \textbf{false}: $1 - e^{-1} = P(T \le 100)$.
\item Statement 4 is \textbf{false}: $\mathrm{Var}(T) = 10\,000$, not $100$.
\end{itemize}

\textbf{Answer: statement 2.}
\end{corrige}

\begin{corrige}[Exercise 2 — True or false?]
\begin{enumerate}
\item \textbf{True.} For $X \sim U(a,b)$, $E(X) = \dfrac{a+b}{2} = \dfrac{0+10}{2} = 5$ and $\mathrm{Var}(X) = \dfrac{(b-a)^2}{12} = \dfrac{10^2}{12} = \dfrac{100}{12}$.

\item \textbf{True.} The normal curve is symmetric about $x = \mu$, so exactly half of the area lies on each side of the mean: $P(X < \mu) = P(X > \mu) = 0.5$.

\item \textbf{False.} By symmetry of the standard normal curve, $P(Z < -1.5) = P(Z > 1.5) = 1 - P(Z < 1.5) = 1 - 0.9332 = 0.0668$, whereas $P(Z < 1.5) = 0.9332$. The two probabilities are \emph{not} equal.

\item \textbf{True.} This is the memoryless property. Indeed,
\[
P(T > s+t \mid T > s) = \frac{P(T > s+t)}{P(T > s)} = \frac{e^{-\lambda(s+t)}}{e^{-\lambda s}} = e^{-\lambda t} = P(T > t).
\]
\end{enumerate}
\end{corrige}

\begin{corrige}[Exercise 3 — Definitions and key formulas]
\begin{enumerate}
\item If $X \sim U(a,b)$, its probability density function is
\[
f(x) = \begin{cases} \dfrac{1}{b-a} & \text{if } a \le x \le b, \\[4pt] 0 & \text{otherwise.} \end{cases}
\]
The graph is a horizontal segment at height $\dfrac{1}{b-a}$ between $a$ and $b$ (a rectangle of total area $1$), and $0$ elsewhere.

\begin{popfigure}
\begin{tikzpicture}[scale=0.9]
\draw[->] (-0.5,0) -- (7,0) node[below] {$x$};
\draw[->] (0,0) -- (0,2.5) node[left] {$f(x)$};
\draw[very thick, popblue] (1.5,0) -- (1.5,1.5) -- (5,1.5) -- (5,0);
\draw[dashed, popmuted] (1.5,0) node[below] {$a$} -- (5,0) node[below] {$b$};
\draw[dashed, popmuted] (0,1.5) node[left] {$\frac{1}{b-a}$} -- (1.5,1.5);
\fill[pattern=north east lines, pattern color=popblue!50] (1.5,0) rectangle (5,1.5);
\end{tikzpicture}
\legende{Density of the uniform distribution $U(a,b)$.}
\end{popfigure}

\item The density of $Z \sim N(0,1)$ is
\[
\varphi(z) = \frac{1}{\sqrt{2\pi}}\, e^{-z^2/2}, \qquad z \in \mathbb{R}.
\]
Two properties of its curve: it is \textbf{symmetric} about the vertical axis $z = 0$ (so $\varphi(-z) = \varphi(z)$), and it is \textbf{bell-shaped} with a maximum at $z = 0$, the total area under the curve being $1$.

\item The binomial distribution $B(n,p)$ may be approximated by a normal distribution when $n$ is large and $p$ is not too close to $0$ or $1$; the standard conditions are
\[
np > 5 \quad \text{and} \quad n(1-p) > 5 \qquad (\text{some texts use } np \ge 5 \text{ and } nq \ge 5).
\]
The approximating distribution is $N(np,\, np(1-p))$, and a \cle{continuity correction} must be applied.
\end{enumerate}
\end{corrige}

\begin{corrige}[Exercise 4 — Direct calculations with the standard normal table]
We write $\Phi(z) = P(Z < z)$ and use the table of $\Phi$.

\begin{enumerate}
\item $P(Z < 1.25) = \Phi(1.25) = \boxed{0.8944}$.

\item $P(Z > 0.84) = 1 - \Phi(0.84) = 1 - 0.7995 = \boxed{0.2005}$.

\item $P(-1.64 < Z < 2.05) = \Phi(2.05) - \Phi(-1.64) = \Phi(2.05) - \bigl(1 - \Phi(1.64)\bigr)$
\[
= 0.9798 - (1 - 0.9495) = 0.9798 - 0.0505 = \boxed{0.9293}.
\]

\item We need $\Phi(z) = 0.95$. Reading the table in reverse, $\Phi(1.64) = 0.9495$ and $\Phi(1.65) = 0.9505$, so $\boxed{z \approx 1.645}$.

\item $P(|Z| < z) = 0.90 \iff 2\Phi(z) - 1 = 0.90 \iff \Phi(z) = 0.95$, hence $\boxed{z \approx 1.645}$.
\end{enumerate}
\end{corrige}

\begin{corrige}[Exercise 5 — Lifetime of a phone battery]
Since $E(T) = \dfrac{1}{\lambda} = 500$, we have $\lambda = \dfrac{1}{500} = 0.002$.

\begin{enumerate}
\item The density is
\[
f(t) = \begin{cases} 0.002\, e^{-0.002 t} & \text{if } t \ge 0, \\ 0 & \text{if } t < 0. \end{cases}
\]

\item $P(T > 800) = e^{-0.002 \times 800} = e^{-1.6} \approx \boxed{0.2019}$.

\item $P(200 < T < 400) = P(T>200) - P(T>400) = e^{-0.4} - e^{-0.8}$
\[
= 0.6703 - 0.4493 = \boxed{0.2210}.
\]

\item The median $m$ satisfies $P(T > m) = 0.5$, i.e. $e^{-0.002 m} = 0.5$, so
\[
m = \frac{\ln 2}{0.002} = 500 \ln 2 \approx \boxed{346.6 \text{ hours}}.
\]
Note that the median is smaller than the mean: the exponential distribution is skewed to the right.
\end{enumerate}
\end{corrige}

\begin{corrige}[Exercise 6 — Uniform distribution]
Here $a = 2$, $b = 10$, so $b - a = 8$ and $f(x) = \dfrac{1}{8}$ on $[2,10]$.

\begin{enumerate}
\item $E(X) = \dfrac{a+b}{2} = \dfrac{2+10}{2} = \boxed{6}$ ; \quad $\mathrm{Var}(X) = \dfrac{(b-a)^2}{12} = \dfrac{64}{12} = \boxed{\dfrac{16}{3} \approx 5.333}$.

\item $P(X > 5) = \dfrac{10 - 5}{10 - 2} = \dfrac{5}{8} = \boxed{0.625}$.

\item $P(X \le k) = \dfrac{k - 2}{8} = 0.75 \implies k - 2 = 6 \implies \boxed{k = 8}$.

\item By the definition of conditional probability,
\[
P(3 < X < 7 \mid X > 4) = \frac{P(4 < X < 7)}{P(X > 4)} = \frac{\frac{7-4}{8}}{\frac{10-4}{8}} = \frac{3}{6} = \boxed{0.5}.
\]
\end{enumerate}
\end{corrige}

\begin{corrige}[Exercise 7 — Masses of bags of cement]
Let $X \sim N(50, 0.16)$, so $\mu = 50$ and $\sigma = \sqrt{0.16} = 0.4$. We standardise with $Z = \dfrac{X - 50}{0.4}$.

\begin{enumerate}
\item $P(X < 49.5) = P\!\left(Z < \dfrac{49.5 - 50}{0.4}\right) = P(Z < -1.25) = 1 - \Phi(1.25)$
\[
= 1 - 0.8944 = 0.1056.
\]
About $\boxed{10.56\%}$ of the bags weigh less than $49.5$ kg.

\item $P(49.6 < X < 50.5) = P\!\left(\dfrac{-0.4}{0.4} < Z < \dfrac{0.5}{0.4}\right) = P(-1 < Z < 1.25)$
\[
= \Phi(1.25) - \Phi(-1) = 0.8944 - (1 - 0.8413) = 0.8944 - 0.1587 = 0.7357.
\]
About $\boxed{73.57\%}$ of the bags weigh between $49.6$ kg and $50.5$ kg.

\item We need $m$ with $P(X > m) = 0.02$, i.e. $P(X < m) = 0.98$. From the table, $\Phi^{-1}(0.98) \approx 2.0537$. Hence
\[
m = 50 + 0.4 \times 2.0537 \approx \boxed{50.82 \text{ kg}}.
\]
\end{enumerate}
\end{corrige}

\begin{corrige}[Exercise 8 — Waiting time at an Express Union counter]
$T \sim \mathrm{Exp}(\lambda)$ with $\lambda = 0.2$ per minute.

\begin{enumerate}
\item $E(T) = \dfrac{1}{\lambda} = \dfrac{1}{0.2} = \boxed{5 \text{ minutes}}$.

\item $P(T > 5) = e^{-0.2 \times 5} = e^{-1} \approx \boxed{0.3679}$.

\item By the memoryless property of the exponential distribution,
\[
P(T > 3 + 2 \mid T > 3) = P(T > 2) = e^{-0.2 \times 2} = e^{-0.4} \approx \boxed{0.6703}.
\]
\textbf{Comment:} the probability of waiting at least $2$ more minutes is exactly the same as the probability of waiting at least $2$ minutes from the start. The fact that $3$ minutes have already elapsed without any arrival gives \emph{no information} about the remaining waiting time: the process ``forgets'' the past. This is the characteristic memoryless property of the exponential distribution.
\end{enumerate}
\end{corrige}

\begin{corrige}[Exercise 9 — Mock examination marks (13 marks)]
Let $X \sim N(52, \sigma^2)$ be the mark of a candidate.

\begin{enumerate}
\item $P(X > 70) = 0.15 \iff P(X < 70) = 0.85$. From the standard normal table, $\Phi^{-1}(0.85) \approx 1.04$ (since $\Phi(1.04) = 0.8508$). Standardising:
\[
\frac{70 - 52}{\sigma} = 1.04 \implies \sigma = \frac{18}{1.04} \approx 17.3 \quad (1 \text{ d.p.}). \qquad \blacksquare
\]
\emph{Mark scheme (4 marks):} M1 standardising; A1 correct $z$-value $1.04$ (accept $1.036$); M1 solving for $\sigma$; A1 $\sigma \approx 17.3$ correctly rounded.

\item With $\sigma = 17.3$:
\[
P(40 < X < 60) = P\!\left(\frac{40-52}{17.3} < Z < \frac{60-52}{17.3}\right) = P(-0.69 < Z < 0.46)
\]
\[
= \Phi(0.46) - \Phi(-0.69) = 0.6772 - (1 - 0.7549) = 0.4321.
\]
Approximately $\boxed{43.2\%}$ of the candidates score between $40$ and $60$.
\emph{Mark scheme (3 marks):} M1 both standardisations; M1 correct use of $\Phi$ with the negative value; A1 $0.432$ (accept $43\%$–$43.3\%$).

\item We need $P(X < m) = 0.10$, so the standardised value is $\Phi^{-1}(0.10) = -1.2816$:
\[
\frac{m - 52}{17.3} = -1.2816 \implies m = 52 - 1.2816 \times 17.3 \approx 52 - 22.17 = \boxed{29.8}.
\]
\emph{Mark scheme (3 marks):} B1 $z = -1.2816$; M1 equation in $m$; A1 $m \approx 29.8$.

\item Each candidate scores above $70$ with probability $p = 0.15$, independently. The number $Y$ of candidates (out of $3$) scoring above $70$ follows $B(3, 0.15)$:
\[
P(Y = 2) = \binom{3}{2}(0.15)^2(0.85) = 3 \times 0.0225 \times 0.85 = \boxed{0.0574}.
\]
\emph{Mark scheme (3 marks):} B1 binomial model identified; M1 correct formula; A1 $0.0574$.
\end{enumerate}
\end{corrige}

\begin{corrige}[Exercise 10 — Throwing a die: normal approximation to the binomial (12 marks)]
\begin{enumerate}
\item Each throw is a Bernoulli trial with $p = P(\text{six}) = \dfrac{1}{6}$, and the $180$ throws are independent. Hence
\[
X \sim B\!\left(180, \tfrac{1}{6}\right), \qquad E(X) = np = 180 \times \tfrac{1}{6} = \boxed{30},
\]
\[
\mathrm{Var}(X) = np(1-p) = 180 \times \tfrac{1}{6} \times \tfrac{5}{6} = \boxed{25}.
\]
\emph{Mark scheme (3 marks):} B1 distribution; B1 mean; B1 variance.

\item We check the validity conditions:
\[
np = 30 > 5 \qquad \text{and} \qquad n(1-p) = 150 > 5.
\]
Both exceed $5$, so the normal approximation is appropriate. The approximating distribution is
\[
X \approx Y \sim N(30, 25), \quad \text{i.e. } \mu = 30,\ \sigma = \sqrt{25} = 5.
\]
\emph{Mark scheme (3 marks):} M1 both conditions checked numerically; A1 conclusion; B1 parameters $N(30, 25)$.

\item ``Between $25$ and $35$ inclusive'' means $25 \le X \le 35$. With the continuity correction this becomes $P(24.5 < Y < 35.5)$:
\[
P(24.5 < Y < 35.5) = P\!\left(\frac{24.5 - 30}{5} < Z < \frac{35.5 - 30}{5}\right) = P(-1.1 < Z < 1.1)
\]
\[
= 2\Phi(1.1) - 1 = 2(0.8643) - 1 = \boxed{0.7286}.
\]
\emph{Mark scheme (4 marks):} M1 continuity correction $24.5$ and $35.5$; M1 standardising; M1 symmetry $2\Phi - 1$; A1 $0.7286$.

\item ``More than $40$ sixes'' means $X \ge 41$, corrected to $Y > 40.5$:
\[
P(Y > 40.5) = P\!\left(Z > \frac{40.5 - 30}{5}\right) = P(Z > 2.1) = 1 - \Phi(2.1) = 1 - 0.9821 = \boxed{0.0179}.
\]
\emph{Mark scheme (2 marks):} M1 continuity correction $40.5$; A1 $0.0179$.
\end{enumerate}
\end{corrige}

\begin{corrige}[Exercise 11 — Filling machine and telephone calls (16 marks)]
\textbf{Part A.} Let $X \sim N(\mu, 4)$, so $\sigma = 2$ ml.

\begin{enumerate}
\item $P(X < 498) = 0.10$. From the table, $\Phi^{-1}(0.10) = -1.2816$. Standardising:
\[
\frac{498 - \mu}{2} = -1.2816 \implies \mu = 498 + 2 \times 1.2816 = 498 + 2.5632 = 500.5632 \approx 500.6 \text{ ml}. \qquad \blacksquare
\]
\emph{Mark scheme (4 marks):} M1 standardising; B1 $z = -1.2816$ (accept $-1.28$); M1 solving for $\mu$; A1 $\mu \approx 500.6$.

\item $P(X > 502) = P\!\left(Z > \dfrac{502 - 500.6}{2}\right) = P(Z > 0.7) = 1 - \Phi(0.7) = 1 - 0.7580 = \boxed{0.2420}$.
\emph{Mark scheme (2 marks):} M1 standardising; A1 $0.242$.

\item The number $Y$ of under-filled bottles in a sample of $100$ follows $B(100, 0.1)$. Check the conditions:
\[
np = 10 > 5 \qquad \text{and} \qquad n(1-p) = 90 > 5,
\]
so the normal approximation $Y \approx W \sim N(10, 9)$ (with $\sigma = 3$) is valid.

``At least $12$'' means $Y \ge 12$, corrected to $W > 11.5$:
\[
P(W > 11.5) = P\!\left(Z > \frac{11.5 - 10}{3}\right) = P(Z > 0.5) = 1 - \Phi(0.5) = 1 - 0.6915 = \boxed{0.3085}.
\]
\emph{Mark scheme (5 marks):} B1 binomial model $B(100, 0.1)$; M1 validity conditions checked; B1 parameters $N(10, 9)$; M1 continuity correction $11.5$; A1 $0.3085$.
\end{enumerate}

\textbf{Part B.}
\begin{enumerate}
\item[(d)] Let $C \sim \mathrm{Poisson}(30)$. Since $\lambda = 30 > 15$ (the usual rule of thumb: the Poisson distribution is well approximated by a normal distribution when its mean is large, typically $\lambda \ge 15$), we may use
\[
C \approx V \sim N(30, 30), \quad \text{i.e. } \mu = 30,\ \sigma = \sqrt{30} \approx 5.477.
\]
(Recall that for a Poisson distribution both the mean and the variance equal $\lambda$.)

``More than $40$ calls'' means $C \ge 41$, corrected to $V > 40.5$:
\[
P(V > 40.5) = P\!\left(Z > \frac{40.5 - 30}{\sqrt{30}}\right) = P(Z > 1.92) = 1 - \Phi(1.92) = 1 - 0.9726 = \boxed{0.0274}.
\]
\emph{Mark scheme (5 marks):} B1 justification $\lambda = 30$ large; B1 parameters $N(30, 30)$; M1 continuity correction $40.5$; M1 standardising with $\sqrt{30}$; A1 $0.0274$ (accept $0.027$–$0.028$).
\end{enumerate}
\end{corrige}

\newpage

\coursec{Summary sheet}

\begin{popfigure}
\begin{tikzpicture}[
  mindmap,
  grow cyclic,
  every node/.style={concept},
  root concept/.append style={concept color=popblue, font=\bfseries\small, text=white},
  level 1/.append style={level distance=3.4cm, sibling angle=60, font=\scriptsize},
  level 1 concept/.append style={minimum size=2.6cm, text width=2.4cm, align=center},
]
\node[root concept]{Special Continuous Distributions}
  child[concept color=poporange]{ node {Exponential $f(x)=\lambda e^{-\lambda x}$} }
  child[concept color=popgreen]{ node {Uniform on $[a,b]$} }
  child[concept color=poppurple]{ node {Normal $N(\mu,\sigma^{2})$} }
  child[concept color=popgold]{ node {Standardising $Z=\dfrac{X-\mu}{\sigma}$} }
  child[concept color=poppink]{ node {Normal approx.\ to Binomial} }
  child[concept color=popblue!70]{ node {Normal approx.\ to Poisson} };
\end{tikzpicture}
\legende{Mind map of the chapter: the three special continuous distributions, standardisation, and the two normal approximations with continuity correction.}
\end{popfigure}

\begin{bilan}
\textbf{Key definitions}
\begin{itemize}
\item A \cle{continuous random variable} $X$ has a \cle{probability density function} (pdf) $f$ with $f(x)\geq 0$ and $\displaystyle\int_{-\infty}^{+\infty} f(x)\,dx = 1$; probabilities are areas: $P(a\leq X\leq b)=\displaystyle\int_a^b f(x)\,dx$.
\item \cle{Exponential distribution} with parameter $\lambda>0$: $f(x)=\lambda e^{-\lambda x}$ for $x\geq 0$ (and $0$ otherwise). It models waiting times (lifetime of a bulb, time between arrivals).
\item \cle{Uniform (rectangular) distribution} on $[a,b]$: $f(x)=\dfrac{1}{b-a}$ for $a\leq x\leq b$. Every sub-interval of equal length has equal probability.
\item \cle{Normal distribution} $X\sim N(\mu,\sigma^{2})$: bell-shaped pdf $f(x)=\dfrac{1}{\sigma\sqrt{2\pi}}\,e^{-\frac{(x-\mu)^{2}}{2\sigma^{2}}}$, symmetric about $x=\mu$.
\item \cle{Standard normal variable}: $Z\sim N(0,1)$; its cdf is tabulated as $\Phi(z)=P(Z\leq z)$.
\end{itemize}

\textbf{Formulas and results to know}
\begin{itemize}
\item Exponential: $E(X)=\dfrac{1}{\lambda}$, $\mathrm{Var}(X)=\dfrac{1}{\lambda^{2}}$, $P(X>x)=e^{-\lambda x}$, cdf $F(x)=1-e^{-\lambda x}$.
\item Uniform on $[a,b]$: $E(X)=\dfrac{a+b}{2}$, $\mathrm{Var}(X)=\dfrac{(b-a)^{2}}{12}$, $P(c\leq X\leq d)=\dfrac{d-c}{b-a}$.
\item Normal: $E(X)=\mu$, $\mathrm{Var}(X)=\sigma^{2}$; standardise with $Z=\dfrac{X-\mu}{\sigma}$.
\item Table use: $\Phi(-z)=1-\Phi(z)$; $P(Z>z)=1-\Phi(z)$; $P(a<Z<b)=\Phi(b)-\Phi(a)$; $P(|Z|<z)=2\Phi(z)-1$.
\item Normal approximation to binomial: if $X\sim B(n,p)$ with $np>5$ and $n(1-p)>5$ (large $n$), then $X\approx N(np,\,np(1-p))$.
\item Normal approximation to Poisson: if $X\sim \mathrm{Po}(\lambda)$ with $\lambda$ large ($\lambda>10$ as a rule of thumb), then $X\approx N(\lambda,\lambda)$.
\item \cle{Continuity correction}: a discrete integer $k$ becomes the continuous interval $[k-0.5,\,k+0.5]$. Thus $P(X\leq k)\to P(Y<k+0.5)$, $P(X<k)\to P(Y<k-0.5)$, $P(X\geq k)\to P(Y>k-0.5)$, $P(X=k)\to P(k-0.5<Y<k+0.5)$.
\end{itemize}

\textbf{Methods}
\begin{itemize}
\item To find a normal probability: write the inequality, standardise both sides, rewrite using $\Phi$, read the table, subtract where necessary.
\item To find a value $x$ given a probability: find $z$ from the table (inverse reading), then $x=\mu+\sigma z$.
\item For approximations: check the validity conditions, compute $\mu$ and $\sigma^{2}$ from the original distribution, apply the continuity correction, then standardise.
\item Sketching the bell curve and shading the required area helps choose between $\Phi(z)$ and $1-\Phi(z)$.
\end{itemize}

\textbf{Common mistakes}
\begin{itemize}
\item Forgetting the continuity correction, or applying it in the wrong direction ($P(X\geq k)$ needs $k-0.5$, not $k+0.5$).
\item Confusing $\sigma$ and $\sigma^{2}$: the second parameter of $N(\mu,\sigma^{2})$ is the \emph{variance}; standardising divides by $\sigma$.
\item Using $\Phi(z)$ directly for negative $z$ instead of $1-\Phi(-z)$.
\item Applying the normal approximation without checking $np>5$, $nq>5$ or $\lambda$ large.
\item For the uniform distribution, writing the variance as $\dfrac{(b-a)^{2}}{2}$ or forgetting the factor $12$.
\item For the exponential, mixing up $E(X)=\lambda$ with the correct $E(X)=\dfrac{1}{\lambda}$.
\item Treating $P(X=k)$ as non-zero for a continuous variable: for continuous $X$, $P(X=k)=0$.
\end{itemize}

\textbf{GCE tips}
\begin{itemize}
\item Always state the distribution and its parameters clearly, e.g.\ $X\sim N(50,16)$, before standardising.
\item Show the standardisation step explicitly: method marks are awarded for $Z=\dfrac{X-\mu}{\sigma}$.
\item Quote table values to 4 decimal places and give final answers to 3 or 4 significant figures.
\item In approximation questions, one mark is usually reserved for the continuity correction and one for the validity conditions: write both.
\item Draw a quick sketch of the normal curve for every probability question; it prevents sign errors with $\Phi$.
\item Check plausibility: a probability must lie in $[0,1]$, and $P(X>\mu)=0.5$ for any normal variable.
\end{itemize}
\end{bilan}

\findecours

\end{document}
